\section{Good towers: models over Hodge loci and geometry}\label{sec:pf:towergeo}

\providecommand{\towergeock}{\operatorname{ck}}
\providecommand{\towergeola}{\operatorname{la}}
\providecommand{\towergeocyc}{\operatorname{cyc}}
\providecommand{\towergeoco}{\operatorname{co}}
\providecommand{\towergeodiv}{\operatorname{div}}
\providecommand{\towergeoRI}{\mathrm{RI}}
\providecommand{\towergeoHC}{\mathcal M}
\providecommand{\towergeoG}{\mathbf G}
\providecommand{\towergeoFC}{\mathrm{FC}}
\providecommand{\towergeoLC}{\mathrm{LC}}
\providecommand{\towergeoRel}{\mathrm{Rel}}
\providecommand{\towergeoTriv}{\mathrm{Triv}}
\providecommand{\towergeoTw}{\mathrm{Tw}}
\providecommand{\towergeoreg}{\phi^{\mathrm{reg}}}

This section proves the following theorem. The geometric package $\towergeoG$ and the statements
(PI), (Av), (W), (F), (CP) are described in \cref{pf:towergeo:sec-statements}; the statements
$\sfname{GoodTowerExistsX}$ and $\sfname{VolumeTheoremX}$ are those of \cref{def:covolume:tower},
recalled in \cref{pf:towergeo:sec-towers}.

\begin{theorem}\label{pf:towergeo:main}
Assume the geometric package $\towergeoG$, the integrability of points \textup{(PI)}, the avoidance
statement \textup{(Av)}, the Weil pairing statement \textup{(W)}, the positivity of families
\textup{(F)} and the comparison of pairs \textup{(CP)}. Then $\sfname{GoodTowerExistsX}$ and
$\sfname{VolumeTheoremX}$ hold.
\formallabel{GlobalTheory.CompanionTheory.Ar.Top.goodTower_and_volume}
\end{theorem}

It is the central step of the proof of \cref{thm:tower}: the six hypotheses are derived from the
inputs of \cref{thm:tower} in the proof of \lref{GlobalTheory.T.good_tower_volume_theorem}.

\emph{Structure of the proof.} The volume half is the volume theorem
\lref{GlobalTheory.T.volume_theorem_of_inputs}, which takes, besides inputs supplied elsewhere,
four geometric statements as hypotheses. This section proves them:
\begin{enumerate}[label=(\arabic*)]
\item the twisting statement (TA): automorphisms of the CM field fixing $\sqrt{-d}$ act on the unitary
  moduli spaces (\cref{pf:towergeo:sec-twist});
\item the trivial flag statement (TF): on every component of the special fibre at a Hodge--Tate free
  prime, the root lines of a Hodge flag are the graded pieces of a flag of the trivial bundle; it follows
  from Fargues' form of the Hodge--Tate theorem (\cref{pf:towergeo:sec-triv});
\item the Hodge comparison (HC): an identity between the Chern numbers of the Hodge--Tate and de Rham
  flags on models over discrete valuation rings, deduced from the positivity of families
  (\cref{pf:towergeo:sec-hodge});
\item the height constants statement: there are constants $C_h,C_u$, uniform in all data, giving
  positivity, lower bounds for heights of points and an upper bound for the mean of the cusp weight
  for regularised Hodge bundles (\cref{pf:towergeo:sec-nine}).
\end{enumerate}
The existence half is proved in \cref{pf:towergeo:sec-exist}: the components of the tower are the
generic points of the global model, and its levels are root models built from Hodge components with
prescribed sets of primes. \Cref{pf:towergeo:sec-nef} proves the nefness of root lines on the special fibres,
at Hodge--Tate free primes from the trivial flags and at the other primes from a statement on semi-abelian
schemes over curves; it is used for the positivity in (GH) and in the proof of the volume theorem. Two
auxiliary statements on intersection numbers, used in the proof of the volume theorem, are proved in
\cref{pf:towergeo:sec-aux}.

\emph{Scope.} The formal proofs of the results of this section are long (about $150{,}000$ lines in
\cite{LeanSource2026}). We prove in full the assembly of each statement above and the arguments which
are self-contained: the twisting statement, the passage from local relations to Chern number identities
and the descent through ramified extensions, the reduction of the relations for lattice families to refined
coordinates, the nefness of trivial flags and the reduction of nefness to curves, the integral of the cusp weight
over fundamental sets and the mean of the logarithm of first minima, the bound for the mean of the cusp weight,
the positivity of the regularised lines and the height identity at integral points. The remaining lemmas are
stated, and cited by the formal names of their proofs.

% ---------------------------------------------------------------------------
\subsection{Coordinates, cycles and Chern numbers}\label{pf:towergeo:setting}

\subsubsection*{Multiprojective coordinates}
Let $K$ be a nonempty finite set and $(J_k)_{k\in K}$ finite sets. For a commutative ring $R$ put
$R[x]=R[x_{k,i}:k\in K,\,i\in J_k]$, graded by multidegrees $a\in\Z^K$. For a field $\Omega$ and a
family $y=(y_{k,i})$ of elements of $\Omega$ the \emph{coordinate ideal} $P_y\subset\Z[x]$ is the kernel
of the multicone map attached to $y$; it is the multihomogeneous ideal of the closure in
$\prod_k\mathbb P^{J_k}_\Z$ of the point with coordinates $y$. We say that $y$ \emph{defines an
arithmetic variety of dimension $d+1$} if $P_y$ is one in the sense of \cref{def:covolume:hs}(a): a
relevant multihomogeneous prime with $P_y\cap\Z=0$ and $\dim\Z[x]/P_y=d+1+|K|$. For
$s\in\prod_kJ_k$ the \emph{Segre coordinate} is $y^s=\prod_ky_{k,s_k}$, and for a multidegree $b$ we put
$y_s^b=\prod_ky_{k,s_k}^{b_k}$. For a ring $o$ with $o\to\Omega$ and $y^s\ne0$, the \emph{chart ring} is
$o[y]_s=o[y^t/y^s:t\in\prod_kJ_k]\subseteq\Omega$.

Let $O$ be a discrete valuation ring with an $O$-algebra structure on $\Omega$. The ideal
$P^O_y\subset O[x]$ is defined in the same way. A relevant multihomogeneous prime $P\subset O[x]$ is a
\emph{model over $O$ of relative dimension $d$} if $P\cap O=0$ and $\dim O[x]/P=d+1+|K|$. For
$P\subset O[x]$ and $a\in O[x]$ the \emph{divisor} $\towergeodiv_P(a)$ is the function on prime ideals
$W$ which is the length of $(O[x]/(P+(a)))_W$ if $W$ is a relevant minimal prime of $P+(a)$, and $0$
otherwise. For a relevant prime $W$ containing a uniformiser and multidegrees $b_1,\dots,b_n$,
$\langle b_1\cdots b_n\rangle_W$ is the intersection number, on the variety over the residue field of $O$
defined by $W$, of the line bundles of multidegrees $b_1,\dots,b_n$.

\begin{definition}[Cycle numbers, local pairs and Chern numbers]\label{pf:towergeo:def-numbers}
Let $P\subset O[x]$ and let $Z$ be a finitely supported function on the primes of $O[x]$.
\begin{enumerate}[label=(\alph*)]
\item For a finite set $\sigma$, classes $c\colon\sigma\to\Z^K$, a polynomial
  $H=\sum_\mu H_\mu X^\mu\in\R[X_s:s\in\sigma]$ and a list $Q$ of multidegrees, the \emph{cycle number}
  is $\towergeocyc(Z,c,H,Q)=\sum_\mu H_\mu\sum_WZ(W)\,\langle c^\mu\cdot Q\rangle_W$, where $c^\mu$ is the
  list in which $c(s)$ occurs $\mu_s$ times.
\item For $a\in O$, classes $x\colon\{1,\dots,r\}\to\Z^K$ and $k\ge1$, the \emph{root number} is
  $\towergeola_P(a,x,k,Q)=k\cdot\towergeocyc(\towergeodiv_P(a),x,\sum_iX_i^{k-1},Q)$.
\item A \emph{local pair} $T=(x,x',f,g)$ indexed by $\sigma$ consists of classes
  $x,x'\colon\sigma\to\Z^K$ and polynomials $f,g\colon\sigma\to O[x]$. For $\sigma=\{1,\dots,r\}$ its
  \emph{$k$-th Chern number} is
  \[
  \towergeock(T,P,k,Q)=\sum_{i=1}^r\sum_W\bigl(\towergeodiv_P(f_i)-\towergeodiv_P(g_i)\bigr)(W)\,
  \Bigl\langle\sum_{a+b=k-1}x_i^a\,{x'_i}^b\cdot Q\Bigr\rangle_W .
  \]
  Restriction along a map $e\colon\sigma'\to\sigma$ is written $T\circ e$, and $T\otimes O'$ denotes the
  image of $T$ under a ring map $O\to O'$. The pair $T$ is \emph{vertical} (for a uniformiser $\varpi$) if
  for every $i$, $f_i,g_i$ are multihomogeneous of degrees with difference $x_i-x'_i$, $f_i,g_i\notin P$,
  the divisor of $f_i/g_i$ is supported on primes containing $\varpi$, and $f_i/g_i$ is a unit at the
  relevant primes containing $P$ but not $\varpi$.
\item For a permutation $\varphi$ of a finite set $\Sigma$, a prime $p$, $k\ge1$ and
  $\tau,\sigma\in\Sigma$,
  \[
  \towergeoco(\varphi,p,k;\tau,\sigma)=\sum_{n\ge0}[\varphi^{n+1}\tau=\sigma]\,p^{-k(n+1)}
  -\frac{[\sigma=\tau]}{p-1}.
  \]
\end{enumerate}
\formal{GlobalTheory.CompanionTheory.Ar.Num.cycNum}
\formal{GlobalTheory.CompanionTheory.Ar.Num.laNum}
\formal{GlobalTheory.CompanionTheory.Ar.Num.LocPair.ck}
\formal{GroundTheory.Heights₂.Heights.ModelLeaves.co}
\end{definition}

If $\towergeodiv(f_i)-\towergeodiv(g_i)$ represents the difference of the classes of the line bundles of
multidegrees $x_i$ and $x'_i$, then $\towergeock(T,P,k,Q)$ is the intersection of $Q$ with the difference
of the $k$-th power sums of the two families of Chern roots, since $a^k-b^k=(a-b)\sum_{j<k}a^jb^{k-1-j}$;
this is the meaning of the Chern numbers below. Three properties are immediate from the definitions:
$\towergeocyc(Z,c,H,Q)$ is $\R$-linear in $H$; renaming the variables of $H$ along a map $e$ of index sets
amounts to replacing $c$ by $c\circ e$; and $\towergeock$ and $\towergeola$ are sums over the indices,
hence additive under a decomposition of the index set and invariant under its permutations.

% ---------------------------------------------------------------------------
\subsection{Unitary moduli, models and Hodge components}\label{pf:towergeo:sec-unitary}

The following data are those of the formal treatment of unitary Shimura varieties; compare
\cite{Kottwitz1992,Lan2013}.
\begin{enumerate}[label=(\alph*)]
\item A \emph{host} $H$ on a CM field $E$ consists of a prime $d\equiv7\bmod16$, an element
  $\sqrt{-d}\in\OO_E$ with square $-d$, the element $w_0=(1+\sqrt{-d})/2\in\OO_E$, complex conjugation on
  $\OO_E$, the condition that $E$ is unramified over its maximal real subfield, and a square root
  $\rho_0$ of $-d$ in $\C$. Its \emph{CM type} is $\Psi_H=\{\tau\colon E\to\C:\tau(\sqrt{-d})=\rho_0\}$;
  $f=|\Psi_H|$. We write $\mathrm{rd}_E=|d_E|^{1/[E:\Q]}$, $L_H=\Re\frac{L'}{L}(1,\chi_d)$ with $\chi_d$
  the quadratic character modulo $d$, and $\ell_H(B)=\log(B\,\mathrm{rd}_E)+|L_H|$ for $B>0$.
\item A \emph{datum} $D$ of signature $(p,q)$ is a finitely generated projective $\OO_E$-module $\Lambda$
  of rank $N=p+q$ (we write $p_D=p$, $q_D=q$ where $p$ also denotes a prime) with a unimodular alternating $\Z$-bilinear form $\langle\,,\rangle$ satisfying
  $\langle ax,y\rangle=\langle x,\bar ay\rangle$, admitting a positive complex structure on
  $\Lambda_\R$ for which the action of $a\in\OO_E$ has trace $\sum_{\tau\in\Psi_H}(p\tau(a)+q\bar\tau(a))$.
  The \emph{root indices} are $\towergeoRI(D)=\Psi_H\times\{0,\dots,N-1\}$; the \emph{labels} are the
  elements $(\tau,1)$, $(\tau,2)$ of $\Psi_H\times\{1,2\}$, of ranks $p$ and $q$, indexing the embeddings
  $\tau$ and $\bar\tau$ of $\OO_E$; the swap $(\tau,1)\leftrightarrow(\tau,2)$ is written $s\mapsto s^\vee$.
\item A \emph{level} $l$ consists of an odd $M\ge3$, an ideal $\mathfrak M\subseteq\OO_E$ with
  $\mathfrak M\bar{\mathfrak M}=(M)$, $\mathfrak M+\bar{\mathfrak M}=\OO_E$ (generated by elements of
  $\Z[w_0]$) and an integer $A\ge1$; $M'=\lcm(M,A)$, and $K_A=\Q(\rho_0,\zeta_A)\subset\C$.
\item An \emph{object} over a scheme $S$ (over a base ring containing a square root of $-d$ and a
  primitive $A$-th root of unity) is an abelian scheme with $\OO_E$-action, a symmetric rigidified line
  bundle with cubical structure compatible with complex conjugation, a level-$\mathfrak M$ structure $\eta$
  and a level-$A$ structure $\eta_A$ (homomorphisms on $\Lambda$ killing $\mathfrak M\Lambda$, resp.\
  $A\Lambda$), satisfying the determinant condition of signature $(p,q)$ and the existence, at geometric
  points, of a Tate isomorphism compatible with the Weil pairings and the level structures. A \emph{map
  of objects} over $f\colon S'\to S$ is a cartesian $\OO_E$-equivariant morphism of the semi-abelian
  schemes over $f$, compatible with the cubical structures and with $\eta$, $\eta_A$.
\item A \emph{moduli space} $X^o$ for $(D,l)$ is a smooth quasi-projective $K_A$-scheme with a universal
  object $\mathrm{univ}$, such that for every reduced $K_A$-scheme $S$ and object $o$ over $S$ there is a
  unique morphism $f\colon S\to X^o$ over $K_A$ admitting a map of objects $o\to\mathrm{univ}$ over $f$,
  together with an infinitesimal rigidity condition.
\item A \emph{coefficient field} $c$ for $l$ is a number field $F_c\subset\C$ containing $\tau(E)$ for all
  $\tau$, $\zeta_{M'}$ and $K_A$. A \emph{model} $Y$ of $X^o$ over a ring $O$ with $\beta\colon O\to F_c$ is
  a normal scheme $Y.X$, projective and flat over $\Spec O$, with a dense open immersion of $X^o_{F_c}$
  into $Y.X\times_OF_c$ and a semi-abelian scheme $Y.t$ with $\OO_E$-action over $Y.X$ extending the
  universal object, whose abelian locus in the generic fibre is exactly the image of $X^o_{F_c}$.
\item For a family of moduli spaces $X^o=(X^o_l)_l$, a model $Y$ of $X^o_l$ over $\OO_{F_c}$ is
  \emph{global} if there are compatible systems of models over all valuation rings of $\C$ above all
  primes, satisfying the comparison with Siegel moduli spaces and the splitting conditions on fibres of
  the formal development, $l$ is good at every prime, $A=1$, $Y$ descends to $K_A$, and $Y$ dominates all
  these local models. $Y$ has \emph{split components} if for every generic point of a component of $Y.X$,
  with residue field $\Omega$ and constant map $F_c\to\Omega$, every element of $\Omega$ algebraic over
  $\Z$ lies in the image of $F_c$.
\end{enumerate}

\begin{definition}[Hodge components]\label{pf:towergeo:def-hodgecpt}
Let $Y$ be a model over $O_0$, $S$ a finite set of primes, $\Omega$ a field and $y$ coordinates in
$\Omega$. A \emph{Hodge component} $\towergeoHC$ of $Y$ with primes $S$ and coordinates $y$ consists of:
\begin{enumerate}[label=(\roman*)]
\item a normal scheme $\towergeoHC.X$ with a point $\xi\colon\Spec\Omega\to\towergeoHC.X$ of which every
  point is a specialisation, whose local rings are given by the Segre coordinates of $y$, and a morphism
  $f\colon\towergeoHC.X\to Y.X$; a field $\Omega_0$, a generic point $\xi_0$ of a component of $Y.X$ with
  residue field $\Omega_0$, dominated by $\xi$ through an embedding $\Omega_0\to\Omega$, and an embedding
  $F_c\to\Omega$ of the constants;
\item a semi-abelian scheme $t$ over $\towergeoHC.X$ with $\OO_E$-action and a map to $Y.t$ over $f$;
  the Lie algebra $\mathrm{Li}$ of $t$ at $\xi$, an $\Omega$-vector space with $\OO_E$-action $\rho$,
  modelled by tangent vectors;
\item a function $R$ on primes with $R(p)\ge4$ and full $p^{R(p)}$-torsion of $t$ at $\xi$ for $p\in S$;
\item for each label $s$, the $\rho$-eigenspace $V^{\mathrm{gen}}_s\subseteq\mathrm{Li}$ of the embedding
  indexed by $s$, the de Rham lattice $\mathrm{dR}_s$ (the eigen-lattice of the Lie algebra of $t$) and,
  for $p\in S$, the Hodge--Tate lattice $\mathrm{HT}_{p,s}$ built from the $p^{R(p)}$-torsion and the
  polarisation; the \emph{patched lattice} $V_s$, equal to $\mathrm{HT}_{p,s}$ at the points over $p\in S$
  and to $\mathrm{dR}_s$ elsewhere, is a vector bundle of the rank of $s$.
\end{enumerate}
$\towergeoHC$ is \emph{Hodge--Tate free at $p$} if the Hodge--Tate $\omega$-module at $p$ and its
restrictions to the eigenspaces are vector bundles of the expected ranks on the fibre over $p$; it is
\emph{$\omega$-regular} if this holds, for the modules built from the $p^{R(p)}$-torsion, at all $p\in S$.
\formal{GlobalTheory.CompanionTheory.Ar.HodgeCpt}
\formal{GlobalTheory.CompanionTheory.Ar.HT.HTFree}
\end{definition}

\begin{definition}[Flags and root lines]\label{pf:towergeo:def-flag}
A \emph{flag} $\mathrm{Fl}$ of a Hodge component $\towergeoHC$ consists of bases $b_s$ of the eigenspaces
$V^{\mathrm{gen}}_s$ which define full flags of subbundles of $V_s$ and descend to a model of the Lie
algebra at $\xi_0$, together with a degree $\deg\mathrm{Fl}\in\N$, the degree of the generic flag over
$\Omega_0$. For a root index $r=(\tau,i)$, the \emph{root lattice} of $r$ is the graded piece of the flag
of $V_{(\tau,1)}$ at position $i$ if $i<p$, and the dual of the graded piece of the flag of $V_{(\tau,2)}$
at position $N-1-i$ if $i\ge p$. For multidegrees $x\colon\towergeoRI(D)\to\Z^K$ and generators
$e\colon\towergeoRI(D)\to\Omega$, $\mathrm{Fl}$ \emph{has root lines $(x,e)$} if in each Segre chart $s$
the root lattice of $r$ is generated by $y_s^{x_r}e_r$, i.e.\ is the line bundle of multidegree $x_r$ with
generic generator $e_r$. For $S_0\subseteq S$ and a local pair $T$ over $\Z$ indexed by $\towergeoRI(D)$,
$\mathrm{Fl}$ \emph{has two levels $(S_0,T)$} if the flag restricts to the lattices patched only at the
primes of $S_0$ with root lines $(T.x,e^0)$, the flag itself has root lines $(T.x',e)$,
$T.f_r(y)=(e_r/e^0_r)\,T.g_r(y)$ for all $r$, and the divisors of $T.f_r/T.g_r$ are supported over the
primes of $S\setminus S_0$.
\formal{GlobalTheory.CompanionTheory.Ar.HodgeCpt.Flag}
\formal{GlobalTheory.CompanionTheory.Ar.HodgeCpt.Flag.TwoLevels}
\end{definition}

% ---------------------------------------------------------------------------
\subsection{Root towers and their origins}\label{pf:towergeo:sec-towers}

We recall the objects of \cref{def:covolume:tower}. Write $\dim\mathrm{Fl}(\Psi,N)=|\Psi|\,N(N-1)/2$.

\begin{definition}[Root models, towers and origins]\label{pf:towergeo:def-towers}
\begin{enumerate}[label=(\alph*)]
\item A \emph{root model} $R$ of type $(\Psi,N)$ consists of coordinates $R.y$ in a field $R.\Omega$
  defining a birational arithmetic variety of dimension $\dim\mathrm{Fl}(\Psi,N)+1$, multidegrees $x_r$
  ($r\in\Psi\times\{0,\dots,N-1\}$) with families of smooth conjugation-invariant Hermitian lines $x(t,r)$
  of multidegree $x_r$ ($t\in\R$), integers $k_A,\delta>0$ and a smooth conjugation-invariant reference
  line $A$.
\item A \emph{root tower} $T$ of type $(\Psi,N)$ consists of $1\le j<N$, a finite nonempty set
  $\mathrm{Comp}$ of components, a finite set $\mathrm{bad}\ni2$ of primes, a rational weight
  $\nu\colon\Psi\times\{0,\dots,N-1\}\to\Q$ and root models $T.M(B,a)$ for $B\in\N$ and
  $a\in\mathrm{Comp}$. For $B\in\N$, $S_B$ denotes the set of primes $r$ with $r\le B$ or
  $r\in\mathrm{bad}$.
\item \emph{Fixed data} $F_x$ on a model $Y$, relative to an \emph{analytic system} $\mathrm{Sys}$ (a
  compact complex manifold of dimension $\dim\mathrm{Fl}(\Psi_H,N)$ with a holomorphic atlas, potentials,
  a normal crossings section and regularised weights indexed by the root indices and one further index),
  consist of a point $\xi_f$ of $Y.X$ with values in a field, a model of the Lie algebra there with
  bases, projective coordinates $y_f\in\Omega_f^{m+1}$ generating the field, a map $\mathrm{pt}$ from
  $\mathrm{Sys}$ to $\C^{m+1}\setminus0$ injective up to scalars, a positive reference potential, integers
  $k_A,q'>0$, weights $n\colon\towergeoRI(D)\to\Z$ and a nonzero element $s_E$.
\item A \emph{root origin} of a root model $R$ over $F_x$ with primes $S$ consists of a Hodge component
  $\towergeoHC$ of $Y$ with primes $S$ and coordinates $R.y$, a flag of it with root lines of multidegrees
  $R.x$, singular Hermitian metrics on the root lines realising the Hodge norms on the interior, an
  embedding of the field of $F_x$ into $R.\Omega$ of degree equal to the degree of the flag, an
  identification of a coordinate block of $R$ with the coordinates $y_f$, an effectivity condition for
  $s_E$, and a realisation of $\mathrm{Sys}$ by a logarithmic resolution of $R.y$ with exceptional ideal
  $J$, compatible with the fundamental weights.
\item A \emph{tower origin} of a root tower $T$ of type $(\Psi_H,p+q)$ on a model $Y$ over
  $\OO_{F_c}$ consists of: $T.j=p$; $r\in T.\mathrm{bad}$ if and only if $r$ is a prime dividing
  $2d_EM'$; for each component $a$ an analytic system $\mathrm{Sys}(a)$ and fixed data $F_x(a)$, whose
  points $\xi_f$ have pairwise distinct generic points exhausting the generic points of $Y.X$, with
  $F_x(a).n(r)=k_Aq'\,\nu(r)$; and root origins $\mathrm{Or}(B,a)$ of $T.M(B,a)$ over $F_x(a)$ with
  primes $S_B$.
\item A \emph{good origin} is a tower origin such that the Hodge components of all $\mathrm{Or}(B,a)$
  are $\omega$-regular, all their realisations have $J=(1)$ and smooth generic fibre, the potentials of the
  last roots $(\tau,N-1)$ are pluriharmonic, and there are reference root models $R(a)$ with root origins
  with empty set of primes, flags of degree $1$, $J=(1)$ and ample reference multidegree.
\item A good origin $o$ satisfies the conditions \textup{(Ex)} if (i) every point of every
  $\mathrm{Sys}(a)$ has a neighbourhood $U$ with a holomorphic map $\lambda\colon U\to\C^{m+1}$,
  $\lambda(y)\in\C^\times\mathrm{pt}(y)$, whose affine ratios $\lambda_i/\lambda_p$ are immersions in the
  charts, and (ii) $R(p)=4$ for the Hodge components of all $\mathrm{Or}(B,a)$ and all $p\in S_B$.
\end{enumerate}
\formal{GlobalTheory.CompanionTheory.Ar.Top.RootTower}
\formal{GlobalTheory.CompanionTheory.Ar.TowerOrigin}
\formal{GlobalTheory.CompanionTheory.Ar.Top.GoodOrigin}
\end{definition}

$\sfname{GoodTowerExistsX}$ asserts: for every CM field $E$, host $H$, datum $D$ with $p,q\ge1$, family of
moduli spaces $X^o$, level $l$, coefficient field $c$ and global model $Y$ of $X^o_l$ over $\OO_{F_c}$ with
split components, there are a root tower $T$ of type $(\Psi_H,p+q)$ and a good origin of $T$ on $Y$
satisfying (Ex). $\sfname{VolumeTheoremX}$ is the statement recalled in \cref{def:covolume:tower}.

% ---------------------------------------------------------------------------
\subsection{The statements and the proof of the main theorem}\label{pf:towergeo:sec-statements}

\subsubsection*{The twisting statement}
Let $H$ be a host on $E$. A \emph{host automorphism} $g$ is a field automorphism $\sigma_g$ of $E$ with
$\sigma_g(\sqrt{-d})=\sqrt{-d}$. For a datum $D$ let $\Lambda^g$ be $\Lambda$ with the action of $\OO_E$
twisted by $\sigma_g$, with the same form. A \emph{twisting isomorphism} is a \emph{pro-isomorphism}
$a\colon\Lambda\to\Lambda^g$: a family $(a_n)_{n\ge1}$ of $\OO_E$-linear maps such that for every $n\ge1$,
$a_n$ induces an isomorphism $\Lambda/n\Lambda\to\Lambda^g/n\Lambda^g$ (i.e.\
$a_n(x)\in n\Lambda^g\Leftrightarrow x\in n\Lambda$, and every $y$ is $a_n(x)+ne$), the forms satisfy
$\langle a_nx,a_ny\rangle\equiv\langle x,y\rangle\bmod n$, and $a_{nm}(x)\equiv a_n(x)\bmod n$. Their set
is $\towergeoTw(D,g)$. For an object $o$ and $a\in\towergeoTw(D,g)$ the \emph{twist} $o^{g,a}$ has the same
abelian scheme, line bundle and cubical structure, the $\OO_E$-action composed with $\sigma_g$, and the
level structures $\eta\circ a_M$, $\eta_A\circ a_A$; it is an object provided the twist is
\emph{admissible}: $\eta\circ a_M$ kills $\mathfrak M\Lambda$, and the determinant condition and the Tate
isomorphisms exist for the twisted action.

\begin{definition}[Twisting statement]\label{pf:towergeo:def-twistaut}
The \emph{twisting statement} \textup{(TA)} asserts: for every host $H$ on a CM field $E$, datum $D$,
level $l$, moduli space $X^o$ for $(D,l)$ and host automorphism $g$, the set $\towergeoTw(D,g)$ is
nonempty, and for every $a\in\towergeoTw(D,g)$ the twist $\mathrm{univ}^{g,a}$ is admissible and there are
an automorphism $\Phi$ of $X^o$ over $K_A$ and a map of objects $\mathrm{univ}^{g,a}\to\mathrm{univ}$ over
$\Phi$.
\formal{GlobalTheory.CompanionTheory.Ar.TwistAut}
\end{definition}

\subsubsection*{Trivial flags and the Hodge comparison}
In the next statements $E,H,D,l,X^o,c$ are as above, $Y$ is a model of $X^o$ over a ring $O_0$, $S$ is
a finite set of primes, $y$ are coordinates in a field $\Omega$ defining an arithmetic variety of dimension
$d+1$, $\towergeoHC$ is a Hodge component of $Y$ with primes $S$ and coordinates $y$, and $A$ is the integral
closure of $\Z$ in $\Omega$. A property holds \emph{at all places of $\Omega$ over $p$} if it holds for
every prime $\mathfrak m$ of $A$ containing $p$ and every discrete valuation ring $O$ which is the
localisation of $A$ at $\mathfrak m$, with compatible algebra structures on $\Omega$. A \emph{fibre
component} at such $O$ is a relevant minimal prime $W$ of $P^O_y+\mathfrak m_OO[x]$. For multidegrees
$b_1,\dots,b_N$ and a fibre component $W$, we write $\towergeoTriv_W(b)$ if there are a field $\Omega'$
over the residue field $k_O$, coordinates $y'$ in $\Omega'$ whose coordinate ideal over $k_O$ is the
reduction of $W$, and a basis of $\Omega'^N$ which is a frame adapted to a full flag of the trivial bundle
of rank $N$ whose graded pieces are the line bundles of multidegrees $b_1,\dots,b_N$.

\begin{definition}[Trivial flag statement]\label{pf:towergeo:def-triv}
The \emph{trivial flag statement} \textup{(TF)} asserts: in the situation above, for every flag
$\mathrm{Fl}$ of $\towergeoHC$ with root lines $(x,e)$ and every $p\in S$ at which $\towergeoHC$ is
Hodge--Tate free, at all places $O$ of $\Omega$ over $p$, for every $\tau\in\Psi_H$ and every fibre
component $W$, $\towergeoTriv_W\bigl((x(\tau,i))_{i<N}\bigr)$ holds.
\formal{GlobalTheory.CompanionTheory.Ar.TrivFlagOnHodge}
\end{definition}

A permutation $\varphi$ of the labels is a \emph{Frobenius at $p$} if for every prime $\mathfrak n$ of
$\OO_{F_c}$ above $p$, every label $s$ and $a\in\OO_E$, the images $u,v$ of $a$ under the embeddings
indexed by $\varphi(s)$ and $s$ satisfy $u\equiv v^p\bmod\mathfrak n$. For $j\le N$ and $m\in\N$ let
$\pi_{j,m}$ be the polynomials in the variables $X_{(\sigma,i)}$, $(\sigma,i)\in\towergeoRI(D)$, given by
\[
\pi_{j,m}(\sigma,1)=\sum_{i<j}X_{(\sigma,i)}^m,\qquad\pi_{j,m}(\sigma,2)=\sum_{i\ge j}(-X_{(\sigma,i)})^m,
\]
and for $\tau\in\Psi_H$, $k\ge1$ put
\begin{equation}\label{pf:towergeo:eq-bracket}
\mathcal B_j(\varphi,p,k,\tau)=k\sum_{s}\towergeoco(\varphi,p,k;(\tau,1),s)\,\pi_{j,k-1}(s),
\end{equation}
the sum over all labels $s$.

\begin{definition}[Chern number comparisons]\label{pf:towergeo:def-fcc}
Let $O$ be a discrete valuation ring, $P\subset O[x]$ and $d,p,j\in\N$.
\begin{enumerate}[label=(\alph*)]
\item For a permutation $\varphi$ of a set of labels $\Psi\times\{1,2\}$ and a local pair $T$ over $O$
  indexed by $\Psi\times\{0,\dots,N-1\}$, the \emph{flag comparison} $\towergeoFC_j(P,d;p,\varphi,T)$ means:
  for all $\tau\in\Psi$, $k\ge1$ and lists $Q$ of multidegrees with $|Q|+k=d+1$,
  \begin{equation}\label{pf:towergeo:eq-fcc}
  \towergeock\bigl(T\circ(\tau,\cdot),P,k,Q\bigr)=\towergeocyc\bigl(\towergeodiv_P(p),T.x,\mathcal B_j(\varphi,p,k,\tau),Q\bigr).
  \end{equation}
\item For a permutation $\varphi$ of a finite set $\Sigma$ and families of local pairs
  $\mathrm{lie}=(\mathrm{lie}_\tau)$, $\mathrm{om}=(\mathrm{om}_\tau)$ indexed by $\tau\in\Sigma$, the
  \emph{label comparison} $\towergeoLC(P,d;p,\varphi,\mathrm{lie},\mathrm{om})$ means: for all $k\ge1$, $Q$
  with $|Q|+k=d+1$ and $\tau\in\Sigma$,
  \begin{equation}\label{pf:towergeo:eq-lcc}
  \towergeock(\mathrm{lie}_\tau,P,k,Q)+\towergeock(\mathrm{om}_\tau,P,k,Q)
  =\sum_{\sigma\in\Sigma}\towergeoco(\varphi,p,k;\tau,\sigma)\,\towergeola_P(p,\mathrm{lie}_\sigma.x,k,Q).
  \end{equation}
\item In the situation of (b), the \emph{relations} $\towergeoRel(P,d;p,\varphi,\mathrm{lie},\mathrm{om})$
  mean: for all $k\ge1$ and $Q$ with $|Q|+k=d+1$ there is $h\colon\Sigma\to\R$ with, for all $\tau$,
  \begin{gather}
  \towergeock(\mathrm{lie}_\tau,P,k,Q)+\towergeock(\mathrm{om}_\tau,P,k,Q)
  =h(\tau)-\frac{\towergeola_P(p,\mathrm{lie}_\tau.x,k,Q)}{p-1},\label{pf:towergeo:eq-rel1}\\
  p^kh(\tau)=h(\varphi\tau)+\towergeola_P(p,\mathrm{lie}_{\varphi\tau}.x,k,Q).\label{pf:towergeo:eq-rel2}
  \end{gather}
\end{enumerate}
\formal{GlobalTheory.CompanionTheory.Ar.Num.FlagChernComparison}
\formal{GlobalTheory.CompanionTheory.Ar.Num.LabelChernComparison}
\formal{GlobalTheory.CompanionTheory.Ar.Num.Relations}
\end{definition}

\begin{definition}[Hodge comparison]\label{pf:towergeo:def-hlfb}
The \emph{Hodge comparison} \textup{(HC)} asserts: in the situation above, for every flag $\mathrm{Fl}$ of
$\towergeoHC$, finite set $S_0$ and local pair $T$ over $\Z$ indexed by $\towergeoRI(D)$ such that
$\mathrm{Fl}$ has two levels $(S_0,T)$, every prime $p\in S\setminus S_0$ at which $\towergeoHC$ is
Hodge--Tate free with $p\ne2$ and $p\nmid d_E$, and every Frobenius $\varphi$ at $p$ on the labels: at all
places $O$ of $\Omega$ over $p$, $\towergeoFC_{p_D}(P^O_y,d;p,\varphi,T\otimes O)$ holds, with $j=p_D$ the first
signature index of $D$.
\formal{GlobalTheory.CompanionTheory.Ar.HodgeLatticeFlagBound}
\end{definition}

\begin{definition}[Positivity of families]\label{pf:towergeo:def-fam}
The \emph{positivity of families} \textup{(F)} asserts: in the situation above with $N>0$, for all bases
$b_s$ of the eigenspaces $V^{\mathrm{gen}}_s$, every $p\in S$ at which $\towergeoHC$ is Hodge--Tate free with
$p\ne2$, $p\nmid d_E$, and every Frobenius $\varphi$ at $p$ on the labels: at all places $O$ of $\Omega$
over $p$ and for every label $t_0$, there are a finite set $\Sigma'$, a map $\epsilon$ from $\Sigma'$ to the
labels with $t_0$ in its image, a permutation $\varphi'$ of $\Sigma'$ with $\epsilon\varphi'=\varphi\epsilon$,
a finite field $F$ of characteristic $p$, an element $w\in O$, embeddings $\sigma_t\colon F\to k_O$, and
\emph{Lie lattice data} $D_f(c)$ in the Segre charts $c$ (groups $M$, $\mathrm{lie}_t$, $W_t$,
$\mathrm{om}_t$, $\mathrm{om}^{\mathrm{HT}}_t$, $\Gamma_t$, $\Delta_t$, $\Theta_t$) forming a
\emph{pre-family} (they are lattice families in the charts, $w^{p-1}$ is a unit times $p$, $p$ is not a
unit in $O$, and structural relations hold in every chart and at all test points), such that in every
chart $\mathrm{lie}_t$ and $W_t$ are the coordinate lattices of $\mathrm{dR}_{\epsilon t}$ and
$\mathrm{HT}_{p,\epsilon t}$ in the basis $b_{\epsilon t}$, and $\mathrm{om}_t$, $\mathrm{om}^{\mathrm{HT}}_t$
are the duals of those of $\mathrm{dR}_{(\epsilon t)^\vee}$ and $\mathrm{HT}_{p,(\epsilon t)^\vee}$.
\formal{GlobalTheory.CompanionTheory.Ar.Num.FamilyExistsPos}
\end{definition}

\subsubsection*{The height constants statement}
The last statement concerns \emph{component data}: coordinates $y$ in a field $\Omega\supseteq\Q$
defining a birational arithmetic variety with smooth generic fibre of dimension $n+1$,
$n=|\Psi|\,j\,(N-j)$, with constants $F'\to\Omega$, a closed torus- and conjugation-invariant boundary
$\mathrm{Bd}$ of the complex points, finitely many normal crossings charts covering the complex points,
rank-$N$ bundles $W_\tau$ ($\tau\in\Psi$) with Hermitian metrics off $\mathrm{Bd}$, a finite set $S$ of
primes with vertical divisors of multidegrees $v_r$ and factors $\phi^{\mathrm{vert}}_r$ ($r\in S$), and a
positive conjugation-invariant Hermitian line $A'$ of multidegree $a'>0$ whose heights on arithmetic
curves are at least $-c_{A'}$ times their degrees; we write $f=|\Psi|$ and $j'=N-j$. \emph{Hodge data} are
subbundles $F_\tau\subseteq W_\tau$ and a section of $\bigotimes_\tau\bigwedge^{j'}(W_\tau/F_\tau)$ of
multidegree $a_\mu$, locally a unit times a wedge of adapted frames; they define the Hermitian line $\mu$ of
multidegree $a_\mu$ and the function $\phi_\mu=\log h_\mu$. A \emph{component origin} realises these
data on a Hodge component $\towergeoHC$ of a model $Y$ with primes $S_0$: $\towergeoHC.X$ covers the
variety with some degree, $W_\tau\cong V^{\mathrm{gen}}_{b(\tau)}\oplus V^{\mathrm{gen}}_{s(\tau)}$ with
$F_\tau=V^{\mathrm{gen}}_{b(\tau)}$ for a choice $\{b(\tau),s(\tau)\}=\{(\tau,1),(\tau,2)\}$ (with $j=p$ or
$j=q$ accordingly), $N=p+q$, the metrics are the Gram matrices of the polarisation on the interior,
$\mathrm{Bd}$ is the complement of the interior, $S\cap S_0=\emptyset$, and conditions on torsion and
Hodge--Tate lattices hold at the primes of $S$. The origin is \emph{admissible} for $B_1\in\N$ if $Y$ is
global, the origin is good, $S_0$ is the set of primes dividing $2d_EM'$, $|S_0|\le4fN^2$,
$4f+2\le B_1$, $S$ is the set of primes $r\le B_1$ with $r\notin S_0$, and $j'\le j$. A function $U$ on the
complex points is the \emph{cusp weight} of the origin if at every interior complex point it equals the
cusp weight $w^{\mathrm{cusp}}$ (\cref{pf:towergeo:sec-nine}) of the Riemann lattice of the fibre.

A \emph{bounded regularisation setting} $S'$ for a weight $w$ consists of a smooth reference line of
multidegree $a_\mu$ with metric $h_{\mathrm{ref}}$, bounded plurisubharmonic data whose singular weight
$\phi_{\mathrm{sing}}$ on the interior is $w-\log h_{\mathrm{ref}}$, a mollifier, cut-off functions and an
error function $e_1$; it is \emph{tied} if $e_1(t)\to0$ as $t\to0^+$ and the regularised weights of
overlapping pieces differ by at most $e_1(t)$ for small $t$. The regularised weight at $t>0$ is
$\towergeoreg_t$. For $B,\alpha,\beta\in\N$ and $t>0$ let $\mathcal M_t=\mathcal M_t(B,\alpha,\beta)$ be the Hermitian line of
multidegree $\beta(a_\mu+\sum_{r\in S,\,r\le B}v_r)+\alpha a'$ with metric
\[
\bigl(h_{\mathrm{ref}}\,e^{\towergeoreg_t}\bigr)^\beta\Bigl(\prod_{r\in S,\,r\le B}\phi^{\mathrm{vert}}_r\Bigr)^\beta
h_{A'}^\alpha,
\]
made smooth and conjugation invariant, and for an arithmetic curve $P$ on the variety put
$h_{\mathcal M_t}(P)=\mathrm{II}(P;\mathcal M_t)/(\beta\deg P)$. For a function $U$ on the complex points let
$\mathrm{mean}(U)=\int U\,d\mu/\deg_\mu$, where the integral is taken with respect to the measure on the
complex points defined by $\phi_\mu$ and $\deg_\mu$ is the geometric degree of $\mu$.

\begin{definition}[Height constants statement]\label{pf:towergeo:def-nine}
For $C_h,C_u\in\R$, the \emph{height constants statement} $\textup{(GH)}_{C_h,C_u}$ asserts: for every
abelian CM field $E$, host $H$, datum $D$, family of moduli spaces $X^o$, level $l$, coefficient field $c$,
model $Y$ of $X^o_l$ over $\OO_{F_c}$, full intersection system $\mathrm{II}$, component data with Hodge
data, component origin with primes $S_0$ admissible for some $B_1$, function $U$ which is the cusp weight
of the origin, complex model and tied bounded regularisation setting for the weight $\phi_\mu+U/2$:
\begin{enumerate}[label=(\alph*)]
\item (positivity) for all $B\ge4f+2$ and $1\le\alpha\le\beta$ there is $t_0>0$ such that for $0<t\le t_0$,
  $\mathcal M_t(B,\alpha,\beta)$ is semipositive, nef on the fibres and ample on the generic fibre;
\item (heights) for all $B\ge4fN^2$ and $1\le\alpha\le\beta$ there is $t_0>0$ such that for $0<t\le t_0$ and every
  arithmetic curve $P$ on the variety, $h_{\mathcal M_t}(P)\ge-C_hfj'\ell_H(B)-\frac\alpha\beta c_{A'}$;
\item (mean) if $1\le j'$ and $3j'+1\le N$, then $\mathrm{mean}(U)\le C_ufj'\ell_H(B)$ for all
  $B\ge4fN^2$.
\end{enumerate}
\formal{GlobalTheory.CompanionTheory.Ar.Cusp.NineGeoU}
\end{definition}

\subsubsection*{The hypotheses}
The \emph{geometric package} $\towergeoG$ is the conjunction of fourteen statements: model clauses (among
them the existence of families of uniformisations of the moduli spaces, and the compatibility of toroidal
models with levels and Hecke operators), the key formula for theta constants, the extension of
homomorphisms of semi-abelian schemes, a numerical Segre statement, resolution of singularities in
projective coordinates, a semi-abelian statement, the Hodge--Tate statement \textup{(HT)} in the form of
Fargues, GAGA for sections, a slope bound, a package of statements on metrics and intersections (the
\emph{metric package}), the Hilbert--Samuel provision, an analytic provision, height bounds for
push-downs to Grassmannians and an antisymmetric perturbation statement. It is proved from the inputs of
\cref{thm:tower} in the proof of \lref{GlobalTheory.T.good_tower_volume_theorem}. The other
hypotheses of \cref{pf:towergeo:main} are: the \emph{integrability of points} (PI), the integrability, over
the base of a good finite projection of an arithmetic variety, of the fibre sums of $\log\|s\|$ against
products of Chern forms of smooth Hermitian lines; the \emph{avoidance statement} (Av), the existence, for
finitely many multiprojective primes, of multihomogeneous elements of every sufficiently large multidegree
avoiding all of them; the \emph{Weil pairing statement} (W), the existence, uniqueness,
bimultiplicativity, alternation and compatibility under $k\mapsto kl$ of the Weil pairings on the
$k$-torsion of polarised abelian schemes over locally Noetherian schemes \cite{Mumford1970}; the
positivity of families (F) of \cref{pf:towergeo:def-fam}; and the \emph{comparison of pairs} (CP): for a
full intersection system, and root models $M,M'$ attached by Hodge ties to Hodge components exchanged by
an automorphism of $X^o$ covering a twist, the normalised arithmetic intersection numbers of the lines
$M.x(t,\cdot)$ and $M'.x(t,\cdot)$, the latter with the CM type permuted, have the same limit as
$t\to0^+$.

Three further statements occur in the proof. The \emph{flat line statements} (FL) assert that on the
geometric components of models there are flat line generators, unique up to nonzero constants, whose sup
norms at conjugate embeddings agree and which are flat generators for the families of uniformisations.
The \emph{flat degree formula} (FD) asserts that the adelic norms $\mathcal N$ of these flat generators
have degree $\frac{p-q}2\Re\log\Omega(E,\chi)+q(2\log\sqrt d-\log\mathrm{rd}_E)$, where
$\log\Omega(E,\chi)=\gamma+\log(2\pi/\mathrm{rd}_E)-\frac{L'}{L}(1,\chi)$ for the Kronecker character
$\chi$ of conductor $d$ (compare \cref{rem:covolume:omega}). The \emph{comparison statement} (CS) asserts
that the difference of the degree-normalised arithmetic intersection numbers of two realisations of an
arithmetic variety through block maps tends to $0$ as $t\to0^+$.

\subsubsection*{Proof of the main theorem}

\begin{proof}[Proof of \cref{pf:towergeo:main}]
Assume the six hypotheses. By \lref{GlobalTheory.T.flat_zero_from_resolution_and_gaga} and
\lref{GlobalTheory.T.flat_first_from_resolution_and_gaga}, resolution in projective coordinates and
GAGA for sections, both part of $\towergeoG$, imply the two flat line statements (FL).

\emph{Existence.} Let $E,H,D,X^o,l,c,Y$ be as in $\sfname{GoodTowerExistsX}$. By
\cref{pf:towergeo:exist}, applied with $\towergeoG$, (W) and (FL), there are a root tower $T$ of type
$(\Psi_H,p+q)$ and a good origin of $T$ on $Y$ satisfying (Ex).

\emph{Volume.} The volume theorem \lref{GlobalTheory.T.volume_theorem_of_inputs} states that
$\towergeoG$, (PI), (W), the existence of $C_h,C_u$ with $\textup{(GH)}_{C_h,C_u}$, (HC), (TF), (FD), (TA),
(CP) and (CS) imply $\sfname{VolumeTheoremX}$. We supply the hypotheses which are not among ours: (FD)
holds by \lref{GlobalTheory.T.FlatDegree_holds}; $\textup{(GH)}_{C_h,C_u}$ for some $C_h,C_u$ follows
from $\towergeoG$, (W) and (FD) by \cref{pf:towergeo:nine}; (HC) follows from (F) by
\cref{pf:towergeo:hodge}; (TF) follows from the Hodge--Tate statement (HT) of $\towergeoG$ by
\cref{pf:towergeo:triv}; (TA) holds by \cref{pf:towergeo:twist}; and (CS) follows from the metric package of
$\towergeoG$, (Av) and resolution in projective coordinates by
\lref{GlobalTheory.T.comparison_of_inputs_avoidance}. Hence $\sfname{VolumeTheoremX}$ holds.
\end{proof}

\begin{remark}\label{pf:towergeo:rem-hyps}
The avoidance statement enters only through (CS), and (PI) and (CP) only through the volume theorem.
In $\textup{(GH)}_{C_h,C_u}$ the constants are chosen before the field $E$ and all other data, so the bounds
are uniform; the same holds for the constant of $\sfname{VolumeTheoremX}$. The existence half uses
neither (F) nor (CP).
\end{remark}

% ---------------------------------------------------------------------------
\subsection{Twisting by automorphisms of the host}\label{pf:towergeo:sec-twist}

We use the notation of \cref{pf:towergeo:sec-statements}. The following lemma is purely algebraic.

\begin{lemma}[Inverse of a pro-isomorphism]\label{pf:towergeo:proiso-inv}
Let $O$ be a commutative ring, $L,L'$ $O$-modules with $\Z$-bilinear forms $B,B'$, and assume that $L'$ is
projective. For every pro-isomorphism $f=(f_n)\colon(L,B)\to(L',B')$ there is a pro-isomorphism
$h=(h_n)\colon(L',B')\to(L,B)$ such that $h_n(f_n(x))\equiv x\bmod nL$ and $f_n(h_n(y))\equiv y\bmod nL'$
for all $n\ge1$, $x\in L$, $y\in L'$.
\formallabel{GlobalTheory.CompanionTheory.Ar.ProIso.exists_isInv}
\end{lemma}

\begin{proof}
Fix $n\ge1$. The $O$-linear map $L\oplus L'\to L'$, $(x,e)\mapsto f_n(x)+ne$, is surjective by the
surjectivity condition on $f_n$; as $L'$ is projective it has an $O$-linear section $\varsigma_n$. Put
$h_n=\mathrm{pr}_1\circ\varsigma_n$; then $f_n(h_n(y))=y-n\,\mathrm{pr}_2(\varsigma_n(y))\equiv y\bmod n$.
If $f_n(u)\equiv f_n(v)\bmod n$, then $f_n(u-v)\in nL'$, hence $u-v\in nL$ by the injectivity condition.
Applying this to $u=h_n(f_n(x))$, $v=x$ gives $h_n(f_n(x))\equiv x\bmod n$. We check that $h$ is a
pro-isomorphism. If $h_n(y)=nz$ then $y\equiv f_n(h_n(y))=nf_n(z)$, so $y\in nL'$; conversely
$h_n(nz)=nh_n(z)$. Every $x$ equals $h_n(f_n(x))-ne$ for some $e$. Writing $f_n(h_n(x))=x+ne_1$ and
$f_n(h_n(y))=y+ne_2$,
\[
B(h_nx,h_ny)-B'(x,y)=\bigl(B'(x+ne_1,y+ne_2)-B'(x,y)\bigr)-\bigl(B'(f_nh_nx,f_nh_ny)-B(h_nx,h_ny)\bigr),
\]
where the first bracket is $n(B'(x,e_2)+B'(e_1,y)+nB'(e_1,e_2))$ and the second is divisible by $n$ since
$f$ preserves the forms modulo $n$. Finally, for $n,m\ge1$, using $f_{nm}\equiv f_n\bmod n$,
$f_n(h_{nm}(y))\equiv f_{nm}(h_{nm}(y))\equiv y\equiv f_n(h_n(y))\bmod n$, hence
$h_{nm}(y)\equiv h_n(y)\bmod n$ by the first observation.
\end{proof}

\begin{lemma}[Twisting back]\label{pf:towergeo:twback}
Let $g$ be a host automorphism and $a\in\towergeoTw(D,g)$. There is $a'\in\towergeoTw(D,g^{-1})$ such that
for every level $l$, base ring with base data, scheme $S$ over it and object $o$ over $S$, whenever the
twists are admissible there are maps of objects $(o^{g,a})^{g^{-1},a'}\to o$ and
$(o^{g^{-1},a'})^{g,a}\to o$ over $\mathrm{id}_S$.
\formallabel{GlobalTheory.T.twist_inverse_iso}
\formallabel{PositivityTheory.UnitarySelfGeneration.tw_back}
\end{lemma}

\begin{proof}
The module $\Lambda^g$ is finitely generated projective of the same rank
(\lref{GlobalTheory.T.twisted_module_finite_projective_rank}: twisting the scalars of a finitely
generated projective module by a ring automorphism preserves these properties and the rank), so
\cref{pf:towergeo:proiso-inv} gives an inverse $h$ of $a$. Since each $h_n$ is linear for the twisted
action, i.e.\ $h_n(\sigma_g(c)x)=ch_n(x)$, the same maps define $a'=(h_n)\in\towergeoTw(D,g^{-1})$.
Consider $(o^{g,a})^{g^{-1},a'}$: its semi-abelian scheme, line bundle and cubical structure are those of
$o$ and its $\OO_E$-action is the action of $o$ composed with $\sigma_g\sigma_g^{-1}=\mathrm{id}$, so the
identity of the group scheme is a map of semi-abelian schemes over $\mathrm{id}_S$ (it is cartesian,
equivariant and compatible with the cubical structure). Its level structures are
$x\mapsto\eta(a_M(h_M(x)))$ and $x\mapsto\eta_A(a_A(h_A(x)))$. Now $a_M(h_M(x))=x+Me$ for some $e$; since
$M\in\mathfrak M\bar{\mathfrak M}\subseteq\mathfrak M$, $Me\in\mathfrak M\Lambda$ is killed by $\eta$, so
$\eta(x+Me)=\eta(x)\eta(Me)=\eta(x)$; likewise $\eta_A(x+Ae)=\eta_A(x)$ because $\eta_A$ kills $A\Lambda$.
Hence the identity is a map of objects. The second assertion is proved in the same way with
$h_n(a_n(x))\equiv x$.
\end{proof}

\begin{theorem}\label{pf:towergeo:twist}
The twisting statement \textup{(TA)} holds.
\formallabel{GlobalTheory.T.TwistAut_holds}
\formallabel{GlobalTheory.CompanionTheory.Ar.TwistAutProof.twistAut}
\end{theorem}

\begin{proof}
Let $E,H,D,l,X^o,g$ be given. The set $\towergeoTw(D,g)$ is nonempty by
\lref{PositivityTheory.UnitarySelfGeneration.tw_iso}, which constructs a pro-isometry
$\Lambda\to\Lambda^g$ from local isometries at all primes. Let $a\in\towergeoTw(D,g)$ and let $a'$ be as in
\cref{pf:towergeo:twback}. By \lref{GlobalTheory.T.twist_ok_all}, the twist $o^{g,a}$ is admissible for
every object $o$, every host automorphism $g$ and every $a\in\towergeoTw(D,g)$; in particular all twists
below are admissible. The scheme $X^o$ is smooth over the field $K_A$, hence reduced.

By the universal property of $X^o$ applied to the reduced scheme $X^o$ and the object
$\mathrm{univ}^{g,a}$, there is $\Phi\colon X^o\to X^o$ over $K_A$ with a map of objects
$m\colon\mathrm{univ}^{g,a}\to\mathrm{univ}$ over $\Phi$; likewise there is $\Psi$ over $K_A$ with
$m'\colon\mathrm{univ}^{g^{-1},a'}\to\mathrm{univ}$ over $\Psi$. Twisting is functorial: a map of objects
$o'\to o$ over $f$ is also a map $o'^{g,a}\to o^{g,a}$ over $f$, since the actions on both sides are
composed with $\sigma_g$ and the level structures on both sides are precomposed with $a_M$, $a_A$. Maps of
objects compose, because maps of semi-abelian schemes compose
(\lref{ReciprocityTheory.ChartRestriction.G.Geo.Tor.sabComp}: the composite of cartesian equivariant
cubical morphisms is again one) and the compatibilities with the level structures are transitive.
Applying the twist by $(g^{-1},a')$ to $m$ and composing with $m'$ gives a map
$(\mathrm{univ}^{g,a})^{g^{-1},a'}\to\mathrm{univ}$ over $\Psi\circ\Phi$. By \cref{pf:towergeo:twback}
there is also a map $(\mathrm{univ}^{g,a})^{g^{-1},a'}\to\mathrm{univ}$ over the identity. Both
$\Psi\circ\Phi$ and $\mathrm{id}$ are morphisms over $K_A$, so the uniqueness in the universal property
gives $\Psi\circ\Phi=\mathrm{id}$. Exchanging the roles of $(g,a)$ and $(g^{-1},a')$ gives
$\Phi\circ\Psi=\mathrm{id}$. Thus $\Phi$ is an automorphism of $X^o$ over $K_A$, and $m$ is the required
map of objects.
\end{proof}

% ---------------------------------------------------------------------------
\subsection{Trivial flags in the special fibres}\label{pf:towergeo:sec-triv}

\begin{definition}[Trivial flags modulo $\varpi$]\label{pf:towergeo:def-trivmod}
Let $o$ be a commutative ring, $\Omega$ an $o$-algebra which is a field, $\varpi\in o$, $y$ coordinates in
$\Omega$ and $b_1,\dots,b_N$ multidegrees. A \emph{trivial flag of type $b$ modulo $\varpi$}, written
$\towergeoTriv_\varpi(o;y,b)$, consists of matrices $g_s\in M_N(\Omega)$ for the Segre indices $s$ with
$y^s\ne0$ such that:
\begin{enumerate}[label=(\roman*)]
\item the entries of $g_s$ lie in $o[y]_s$, and $\det(g_s)\,h=1+\varpi r$ for some $h,r\in o[y]_s$;
\item for all $s,t$ with $y^s,y^t\ne0$ there are $m\in\N$ and matrices $U,Z$ with entries in $o[y]_t$,
  $U$ upper triangular with $U_{ii}=(y^s/y^t)^m\,y_s^{b_i}/y_t^{b_i}$, such that $(y^s/y^t)^mg_s$ has
  entries in $o[y]_t$ and $(y^s/y^t)^mg_s=g_tU+\varpi Z$.
\end{enumerate}
\formal{GlobalTheory.CompanionTheory.Ar.Num.TrivFlagMod}
\end{definition}

Thus, modulo $\varpi$, the frames $g_s$ of the trivial bundle of rank $N$ are invertible on the charts and
change by upper triangular matrices whose diagonal entries are the transition functions of the line
bundles of multidegrees $b_i$, up to a common twist: the trivial bundle carries, modulo $\varpi$, a full
flag with these graded pieces. If $\varpi'\mid\varpi$, say $\varpi=\varpi'c$, a trivial flag modulo
$\varpi$ is one modulo $\varpi'$: replace $r$ by $cr$ and $Z$ by $cZ$. There are two inputs.

The first input is a construction from $p$-adic Hodge theory. Let $A\subseteq K$ be a ring in a field and
$R,N\in\N$. A \emph{Hodge--Tate chart} over $A$ of a $K$-vector space $V$ with an action of $\OO_E$ consists
of an $\OO_E$-stable $A$-lattice $\omega\subseteq V$ which is a projective $A$-module, an alternating
biadditive pairing $e$ on $\OO_E^N$ with values in $\Z/p^R$ which is perfect modulo $p^R$ and satisfies
$e(ax,y)=e(x,\bar ay)$ (a model of the Weil pairing on the $p^R$-torsion), and a \emph{Hodge--Tate map}
$h\colon\OO_E^N\to\omega$ which is $\OO_E$-linear modulo $p^R\omega$. We put $\omega'=A\,h(\OO_E^N)+p^R\omega$, and for
an embedding $\sigma\colon E\to K$ we let $\omega^{\mathrm{HT}}_\sigma$ be the $\sigma$-eigencomponent of $\omega'$
and $q_\sigma\colon A^N\to\omega^{\mathrm{HT}}_\sigma$ the map $x\mapsto\sum_ix_i\pi_\sigma(h(e_i))$, where
$\pi_\sigma$ is the eigenprojection.

\begin{proposition}[Trivial flags from Hodge--Tate charts]\label{pf:towergeo:htcharts}
Assume the Hodge--Tate statement \textup{(HT)}. Let $\towergeoHC$ be a Hodge component of a model $Y$ with
primes $S$ and coordinates $y$ in $\Omega$, $p\in S$, $N>0$, and let $o$ be a ring with a map
$o\to\Omega$ whose image is $\{a/b:a\in\Z,\,b\in\N,\,p\nmid b\}$. If $\towergeoHC$ is Hodge--Tate free at
$p$ and a flag of $\towergeoHC$ has root lines $(x,e)$, then $\towergeoTriv_{p^{R(p)-2}}(o;y,(x(\tau,i))_i)$
holds for every $\tau\in\Psi_H$.
\formallabel{GlobalTheory.CompanionTheory.Ar.HT.trivFlagMod_rootLines_model}
\formallabel{PositivityTheory.HtPt.Ge.triv_flag_mod_of_hodge_tate_free_of_root_lines}
\end{proposition}

\begin{proof}
Write $R=R(p)\ge4$, fix $\tau$, and let $\sigma\colon E\to\Omega$ be $\tau$ followed by the constants of
$\towergeoHC$, so that $\sigma$ restricted to $\OO_E$ is the embedding of the label $(\tau,1)$. The chart rings
$o[y]_s$ are Noetherian, integrally closed, with fraction field $\Omega$, and $\Omega$ has characteristic $0$
(\lref{PositivityTheory.HtPt.Hcptgen.hcpt_generic}: these are general facts on the chart rings of
Hodge components over $o$); in particular they contain $\sigma'(\OO_E)$ for every embedding $\sigma'$.

\emph{Charts and test points.} By \lref{GlobalTheory.T.hodge_tate_free_flag_lattice_lower} there is, for
every Segre index $s$ with $y^s\ne0$, a Hodge--Tate chart $D_s$ over $o[y]_s$ of the dual of the Lie algebra of
$\towergeoHC$, with $N=p_D+q_D$ and $R=R(p)$, such that the charts form a family (the actions and pairings agree,
the lattices agree at the points of overlapping charts, and the Hodge--Tate maps agree there modulo $p^R$), and
every $D_s$ has test points at every $\kappa\in o[y]_s$ which equals $\zeta-1$ for a primitive $p$-th root of
unity $\zeta$ ($\kappa=4$ if $p=2$), provided the Hodge--Tate statement at points holds. The latter follows from
(HT) by \lref{GlobalTheory.T.pointHT_of_fargues_hodge_tate}. Since $\towergeoHC$ is Hodge--Tate free at $p$,
\lref{PositivityTheory.HtPt.Famproj.family_proj} shows that the lattices $\omega_s$ and their eigencomponents
are projective, that the eigencomponents for $\sigma$ and $\sigma\circ c$ ($c$ complex conjugation) admit duality
data, given by an element $d$ generating the different, and that there is such a $\kappa$ in all chart rings.
Since $\kappa$ divides $p^2$ (for $p$ odd, $p=\prod_{i=1}^{p-1}(1-\zeta^i)$ is divisible by $1-\zeta$ in $\Z[\zeta]$, and
$\zeta=\kappa+1$ lies in the chart rings; for $p=2$, $\kappa=4=p^2$), $\kappa p^{R-2}$ divides $p^R$ in every chart ring.

\emph{Exactness modulo $p^{R-2}$.} By \lref{PositivityTheory.HtPt.Ge.chart_sequence}, applied in each chart
with the test points at $\kappa$ and the divisibility $\kappa p^{R-2}\mid p^R$, there is a map $i_s$ which makes
the $\sigma$-eigencomponent of $h_s$ dual, modulo $p^{R-2}$, to the $(\sigma\circ c)$-eigencomponent, such that the
resulting sequence is exact modulo $p^{R-2}$; moreover $\kappa\omega_s\subseteq\omega'_s$, i.e.\ the cokernel of the
Hodge--Tate map is killed by $\kappa$ (this is where (HT) enters, through the test points), hence
$p^2\omega_s\subseteq\omega'_s$; and the maps $q_\sigma$ of the charts are surjective onto
$\omega^{\mathrm{HT}}_{\sigma,s}$ (\lref{PositivityTheory.HtPt.Ga.q_surj}).

\emph{Transitions.} For charts $s,t$, every element of $\omega_s$ lies in $\omega_t$ after multiplication by a
power of $y^s/y^t$ (\lref{GlobalTheory.T.hodge_component_integral_at_prime_in_S}: an element which lies in
the lattices at all points of the chart $s$ in the fibre over $p$ is integral over $o[y]_t$ up to such a power).
Applied to $p^{-R}(h_s(x)-h_t(x))$, which lies in the lattices at all common points by the gluing property, this
gives $n$ with $(y^s/y^t)^n(h_s(x)-h_t(x))\in p^R\omega_t$. Projecting to the $\sigma$-eigencomponent and using
$p^2\omega_t\subseteq\omega'_t$ we obtain $(y^s/y^t)^n\pi_\sigma(h_s(x)-h_t(x))\in p^{R-2}\omega^{\mathrm{HT}}_{\sigma,t}$,
and \lref{GlobalTheory.T.hodge_chart_transition_integral} gives the integrality of the transitions of the
eigencomponents.

\emph{Flag bases and assembly.} By \lref{PositivityTheory.HtPt.Flagb.chart_flag_bases}, applied to the flag
with root lines $(x,e)$, there are generic vectors $w_1,\dots,w_N$ and, in each chart, bases of the dual of
$\omega^{\mathrm{HT}}_{\sigma\circ c,s}$ (of size $p_D$) and of $\omega^{\mathrm{HT}}_{\sigma,s}$ (of size $q_D$)
whose $a$-th element is an upper triangular combination of the $w_{a'}$ with diagonal coefficient the chart
monomial $y_s^{x(\tau,a)}$. Finally \lref{PositivityTheory.HtPt.Ge.trivflag_charts} assembles from the
Hodge--Tate maps, written in these bases, matrices $g_s$ with the properties (i) and (ii) of
\cref{pf:towergeo:def-trivmod} for $\varpi=p^{R-2}$: condition (i) comes from the exactness modulo $p^{R-2}$, and
condition (ii) from the transition property and the triangularity of the flag bases.
\end{proof}

The second input is the reduction of a trivial flag to the special fibre.

\begin{lemma}[{\lref{GlobalTheory.T.trivial_flag_at_prime}}]\label{pf:towergeo:trivat}
Let $O$ be a discrete valuation ring with uniformiser $\varpi$, $\Omega$ an $O$-algebra which is a field,
and assume $\towergeoTriv_\varpi(O;y,b)$. Then $\towergeoTriv_W(b)$ holds for every relevant multihomogeneous
prime $W\subset O[x]$ containing $P^O_y$ and $\varpi$.
\end{lemma}

In its formal proof the reduction of $W$ is the coordinate ideal of a point over an extension of the
residue field, some Segre coordinate of which is nonzero, and the frames $g_s$ reduced modulo $\varpi$ give
the required flag frame.

\begin{theorem}\label{pf:towergeo:triv}
The Hodge--Tate statement \textup{(HT)} implies the trivial flag statement \textup{(TF)}.
\formallabel{GlobalTheory.T.triv_flag_on_hodge_of_fargues}
\formallabel{GlobalTheory.CompanionTheory.Ar.Num.trivFlagOnHodge_of_fargues}
\end{theorem}

\begin{proof}
Let the data of (TF) be given, let $O$ be a place of $\Omega$ over $p$, obtained by localising $A$ at
$\mathfrak m\ni p$, let $\varpi$ be a uniformiser of $O$, $\tau\in\Psi_H$ and $W$ a fibre component. We first
show $\towergeoTriv_\varpi(O;y,x(\tau,\cdot))$. If $N=0$ this holds with empty matrices. Let $N>0$. Since $y$
defines an arithmetic variety, $P_y\cap\Z=0$, so $\Z\to\Omega$ is injective and $\Omega$ has characteristic
$0$. Let $o\subseteq\Omega$ be the subring of fractions $a/b$ with $p\nmid b$. By
\cref{pf:towergeo:htcharts}, $\towergeoTriv_{p^{R(p)-2}}(o;y,x(\tau,\cdot))$ holds. If $p\nmid b$ then $b$
is a unit in $O$: otherwise $b\in\mathfrak m$, and writing $1=ub+vp$ with $u,v\in\Z$ gives
$1\in\mathfrak m$. Hence $o\subseteq O$ inside $\Omega$, every chart ring $o[y]_s$ is contained in $O[y]_s$,
and the matrices witnessing the trivial flag over $o$ witness it over $O$, with the same element
$p^{R(p)-2}$. As $p\in\mathfrak mO=\varpi O$ and $R(p)-2\ge2$, $\varpi$ divides $p^{R(p)-2}$ in $O$, so
$\towergeoTriv_\varpi(O;y,x(\tau,\cdot))$ holds by the remark after \cref{pf:towergeo:def-trivmod}.

The ideal $P^O_y+\mathfrak m_OO[x]$ is multihomogeneous, so its minimal prime $W$ is multihomogeneous
(\lref{GroundTheory.Heights₂.Heights.NumbersDVR.MinimalPrimeMultiHomog.proof}: minimal primes over
multihomogeneous ideals are multihomogeneous); $W$ is relevant and contains $P^O_y$ and
$\varpi\in\mathfrak m_O$. By \cref{pf:towergeo:trivat}, $\towergeoTriv_W(x(\tau,\cdot))$ holds.
\end{proof}

% ---------------------------------------------------------------------------
\subsection{Existence of good towers}\label{pf:towergeo:sec-exist}

Throughout this subsection $E,H,D,X^o,l,c$ are as in $\sfname{GoodTowerExistsX}$, $N=p+q$, and $Y$ is a
model of $X^o_l$ over $\OO_{F_c}$. A \emph{generic point} of $Y.X$ is a point $x$ which is not a proper
specialisation of another point.

\begin{lemma}[Components]\label{pf:towergeo:components}
Assume $\towergeoG$. The set of generic points of $Y.X$ is finite and nonempty.
\formallabel{GlobalTheory.T.generic_points_finite_nonempty}
\formallabel{GlobalTheory.CompanionTheory.Ar.Comps.finite_nonempty_maximal}
\end{lemma}

\begin{proof}
Since $Y.X$ is projective over the Noetherian ring $\OO_{F_c}$, its underlying space is Noetherian
(\lref{GlobalTheory.ModelSum.noetherianSpace_of_isProjOver}: a scheme projective over a Noetherian
ring has a Noetherian underlying space). A Noetherian space has finitely many irreducible components
\cite[I, Prop.~1.5]{Hartshorne1977}, and in a scheme every irreducible closed subset has a unique generic
point \cite[II, Ex.~2.9]{Hartshorne1977}; the generic points of $Y.X$ in our sense are exactly the generic
points of the irreducible components. It remains to show $Y.X\ne\emptyset$, since a nonempty Noetherian
space has an irreducible component. First, $X^o\ne\emptyset$: if $N=0$, the universal property applied to
the zero object over $\Spec K_A$ (\lref{GlobalTheory.Models.RankZero.zeroObj}: for $p+q=0$ the zero
semi-abelian scheme with trivial level structures is an object) gives a morphism $\Spec K_A\to X^o$; if
$N>0$, the function on embeddings $\tau\mapsto p$ for $\tau\in\Psi_H$ and $\tau\mapsto q$ otherwise is the
multiplicity function of a list of embeddings, and the model clauses of $\towergeoG$ provide a family of
uniformisations of $X^o$ of this type, in particular a complex point of $X^o$. Given a point of $X^o$, there
is a point of $X^o\times_{K_A}F_c$ over it, since both factors map to the unique point of $\Spec K_A$, and
the open immersion $X^o_{F_c}\to Y.X\times_{\OO_{F_c}}F_c$ followed by the projection gives a point of
$Y.X$.
\end{proof}

A \emph{reference geometry} $G$ of $Y$ is a Hodge component $\towergeoHC_G$ of $Y$ with empty set of primes
and coordinates $y_G$ in a field $\Omega_G$ of characteristic $0$, with a flag of degree $1$ with root lines
of multidegrees $x_G$, an ideal defining a model with simple normal crossings boundary whose complex points
are exactly the coordinates of the non-interior points, and a block $k_0$ of coordinates which embeds the
generic point. An \emph{analytic realisation} of $G$ consists of singular Hermitian metrics on the root
lines realising the Hodge norms, an analytic system $\mathrm{Sys}$ and a realisation of it by a logarithmic
resolution of $y_G$ with $J=(1)$, of degree $1$, with bijective projection $\pi$, embedded, compatible with
the interior, the fundamental weights and complex conjugation, with reference line of multidegree $e_{k_0}$
and metric $(\sum_i|z_{k_0,i}|^2)^{1/2}$, positive reference potential, and pluriharmonic potentials for the
last roots $(\tau,N-1)$. For $n\colon\towergeoRI(D)\to\Z$, $q'\in\N$ and $s\in\Omega_G$, the triple
$(n,q',s)$ is \emph{effective} for $G$ if at every point of $\towergeoHC_G.X$ outside the fibres over all
primes, in every chart of the block $k_0$ and for every local generator $t$ of the tensor product of the
root lattices with exponents $-n(r)$, the element $s\,y_{k_0,i}^{q'}/t$ is regular, and it is not a unit at
the points outside the abelian locus of the semi-abelian scheme.

We first prove the five constructions used below. The first produces reference geometries; it uses
resolution in projective coordinates, one of the statements of $\towergeoG$.

\begin{proposition}[Reference geometries]\label{pf:towergeo:refgeo}
Assume resolution in projective coordinates. For every point $\xi_0\colon\Spec\Omega_0\to Y.X$ which is the
generic point of a component (in the sense of \cref{pf:towergeo:def-hodgecpt}) there is a reference geometry
$G$ of $Y$ whose point $\xi_0(G)$ has the same generic point as $\xi_0$.
\formallabel{GlobalTheory.T.exists_reference_geometry_same_generic}
\formallabel{GlobalTheory.CompanionTheory.Ar.exists_refGeo}
\end{proposition}

\begin{proof}
The field $\Omega_0$ has characteristic $0$, since the generic points of the components of the flat
$\OO_{F_c}$-scheme $Y.X$ lie over the generic point of $\Spec\OO_{F_c}$. By
\lref{GlobalTheory.T.exists_coords_for_component}, which uses resolution in projective coordinates, there are
coordinates $y_0$ in $\Omega_0$, a scheme $X$ with a point $\xi$ of which every point is a specialisation and whose
local rings are given by the Segre coordinates of $y_0$, a morphism $f\colon X\to Y.X$ with $\xi\circ f=\xi_0$,
constants $F_c\to\Omega_0$ compatible with the structure maps, a semi-abelian scheme $t$ with $\OO_E$-action
over $X$ with a cartesian map to $Y.t$ over $f$, a model of its Lie algebra at $\xi$ with action $\rho$, lattices
$L_s$ for the labels which are vector bundles of the ranks of the labels with generic fibres the $\rho$-eigenspaces
and coincide with the eigenlattices of the Lie algebra outside the fibres over all primes, and an ideal defining a
model of $y_0$ with simple normal crossings boundary whose complex points are exactly the coordinates of the
complex points $z$ of $X$ with $f(z)$ outside the abelian locus. By
\lref{GlobalTheory.T.coords_with_constants_map} these data can be transported to an extension $\Omega$ of
$\Omega_0$ and new coordinates $y$ in $\Omega$, so that in addition there are bases $b_s$ of the eigenspaces
defining full flags of vector bundles of the eigenlattices outside the fibres over all primes (compatible with a
model of the Lie algebra over $\Omega_0$, i.e.\ a generic flag of degree $1$ over $\Omega_0$), and the local rings
outside the fibres over all primes are integrally closed. By
\lref{GlobalTheory.T.generic_coords_descend_subfield} there are coordinates $y'$ in $\Omega$ extending $y$
blockwise, a Hodge component $\towergeoHC$ of $Y$ with empty set of primes and coordinates $y'$, a morphism
$g\colon\towergeoHC.X\to X$ compatible with the points and with $f$, identifying the local rings outside the fibres
over all primes and injective there, and a flag of $\towergeoHC$ of degree $1$ with root lines $(x,e)$; moreover
the point $\xi_0(\towergeoHC)$ has the generic point of $\xi_0$.

It remains to add the block which embeds the generic point. Let $G$ have the coordinates $y'$ together with one
further block $k_0$, indexed by $\prod_kJ'_k$, whose coordinates are the Segre coordinates $y'^s$. The local rings
are still given by the Segre coordinates (the Segre coordinates of the enlarged family are products of those of
$y'$), the root lines keep their generators and have multidegree $0$ in the new block, and the new block embeds
the generic point, because the ratios of Segre coordinates of a coordinate tie in characteristic $0$ generate the
function field (\lref{GlobalTheory.T.embeds_generic_of_char_zero}). The boundary ideal is the image of the
one of $y$ under the renaming of variables; it defines a simple normal crossings model for $y'$
(\lref{GlobalTheory.T.snc_model_reindexed_coords}), and its complex points are the coordinates of the
non-interior points: complex coordinates of $\towergeoHC.X$ restrict to complex coordinates of $X$
(\lref{GlobalTheory.T.complex_coords_reindexed}), and the abelian locus is compatible with the cartesian map
of semi-abelian schemes (\lref{ReciprocityTheory.ChartRestriction.G.Geo.Tor.abLocus_pullback}: a point lies in
the abelian locus of a pull-back if and only if its image lies in the abelian locus).
\end{proof}

\begin{proposition}[Analytic realisations]\label{pf:towergeo:holo}
Assume $\towergeoG$ and \textup{(FL)}, let $Y$ be global with split components, and $p,q\ge1$. Every reference
geometry $G$ of $Y$ has an analytic realisation $\mathcal A$ whose projection $\pi$ has a holomorphic map
$\iota$ with $\iota\circ\pi=\mathrm{id}$.
\formallabel{GlobalTheory.T.exists_holomorphic_section_of_analytic}
\formallabel{AnalyticReference.reference_geometry_holomorphic_inverse}
\end{proposition}

\begin{proof}
By \lref{GlobalTheory.T.hodge_cpt_flag_xdeg}, for every root index $r$ there is a conjugation-invariant
Hermitian line $h_r$ of multidegree $x_G(r)$ which is a \emph{Hodge line} of the flag at $r$: on the interior, its
norm of the generator $e_r$ is the Hodge norm. Let $\chi$ be a simple normal crossings chart of the boundary, with
$a$ boundary coordinates, and $\gamma_\chi$ its parametrisation. By
\lref{AnalyticReference.An.local_bounds}, which uses $\towergeoG$, (FL), globality and split components, for
every $\tau$ and $i$ the function $\phi_{\tau,i}=-\sum_{k\le i}\log h_{(\tau,k)}\circ\gamma_\chi$ has
\emph{log-log growth} at the centre: near $0$ and off the boundary it is smooth, and $|\phi_{\tau,i}|$ and the
norm of its differential are bounded by $C\log\ell(t)$ and $C\ell(t)^Ce^{\ell(t)/2}$, where $\ell(t)$ is the
logarithmic distance to the boundary of the chart; if $a=0$, $\phi_{\tau,i}$ is of class $C^2$ at $0$ with
nonnegative Levi form there, and so is $-\phi_{\tau,N-1}$. By \lref{AnalyticReference.An.of_local}, which uses
the metric package of $\towergeoG$, these local bounds for the lines $h_r$ of a reference geometry (with $p,q\ge1$)
yield an analytic realisation $\mathcal A$ of $G$ together with a holomorphic map $\iota$ with
$\iota\circ\pi=\mathrm{id}$.
\end{proof}

For a weight $\nu\colon\towergeoRI(D)\to\Z$ the \emph{gaps} are $g_\nu(\tau,i)=\nu(\tau,i)-\nu(\tau,i+1)$ for
$i+1<N$, $g_\nu(\tau,N-1)=\nu(\tau,N-1)$, and $g_\nu=0$ at the further index of the analytic system. For
an analytic realisation $\mathcal A$, the \emph{combined line} of exponents $n$ is the tensor product of the
reference lines of $\mathcal A$ with exponents $n$, and $\deg_{\mathcal A}$ is the polynomial in the exponents giving
the degree of the combined line (the top intersection number of the reference lines).

\begin{proposition}[Effective weights]\label{pf:towergeo:exeff}
Assume $\towergeoG$, let $Y$ be global and $p,q\ge1$. Let $G$ be a reference geometry of $Y$ which has an
analytic realisation, and $\nu\colon\towergeoRI(D)\to\Z$ strictly decreasing in $i$ for every $\tau$. Then there
are $k,q'>0$ and $s\ne0$ in $\Omega_G$ such that $(kq'\nu,q',s)$ is effective for $G$.
\formallabel{GlobalTheory.T.exists_effective_weight_section}
\formallabel{GlobalTheory.CompanionTheory.Ar.ExRefEff.exists_eff}
\end{proposition}

\begin{proof}
Let $\mathcal A$ be an analytic realisation. The coordinates $y_G$ are birational
(\lref{GlobalTheory.T.hodge_cpt_birational_arith_var}: the coordinates of a Hodge component with a flag define
a birational arithmetic variety of dimension $\dim\mathrm{Fl}(\Psi_H,N)+1$). By
\lref{GlobalTheory.T.reference_geometry_multihomog_form} there is a multihomogeneous $\beta$ of multidegree
$d_\beta$ with $\beta(y_G)\ne0$, which in every chart is a section of the line bundle of multidegree $d_\beta$ and is
not a unit at the points outside the fibres over all primes and outside the abelian locus.

The gaps $g_\nu$ are $\ge0$ except possibly at the last roots $(\tau,N-1)$: for $i+1<N$, $g_\nu(\tau,i)>0$ by
strict monotonicity. At the last roots both the potentials and their negatives are plurisubharmonic, which is part
of the analytic realisation. Moreover $\deg_{\mathcal A}(g_\nu)>0$ by
\lref{GlobalTheory.T.global_reference_gap_positive}, which uses $\towergeoG$ and globality. Hence the
\emph{nef polarisation} statement of the metric package (applied with the intersection numbers over $\Q$, a field
intersection family by \lref{GroundTheory.Heights₂.Heights.NumbersDVR.IntNumIsFieldIntersection.full}) gives
$k\ge1$ and a nonzero $e\in\Omega_G$ which in every chart is a section of the line bundle of multidegree
$k\,a-\mathbf 1-d_\beta$, where $a$ is the multidegree of the combined line of exponents $g_\nu$ and $\mathbf 1$ the
multidegree with all entries $1$. The multidegree of the combined line is
$a=-\sum_{\tau,i}\nu(\tau,i)\,x_G(\tau,i)$ (\lref{GlobalTheory.CompanionTheory.Ar.ExRefEff.comb_a}: an Abel
summation, the reference lines of the realisation being attached to the fundamental weights).

Put $q'=1$, $n=k\nu$, choose for every block $\kappa$ an index $i_0(\kappa)$ with $y_{\kappa,i_0(\kappa)}\ne0$ (they
exist since $\xi$ lies in some chart), and
\[
s=e\,\beta(y_G)\prod_{\kappa\ne k_0}y_{\kappa,i_0(\kappa)}\prod_re_r^{-n(r)}.
\]
Then $s\ne0$, as all factors are nonzero. The triple $(n,1,s)$ is effective: in a chart of the block $k_0$, a local
generator $t$ of the tensor product of the root lattices with exponents $-n(r)$ is a unit times the chart monomial of
multidegree $-\sum_rn(r)x_G(r)=k\,a$ times $\prod_re_r^{-n(r)}$; the multidegrees of the factors of
$s\,y_{k_0,i}$ add up to $k\,a$, so $s\,y_{k_0,i}/t$ is a unit times a product of elements of the chart ring, hence
regular, and it is not a unit outside the abelian locus because the factor coming from $\beta(y_G)$ is not
(\lref{GlobalTheory.CompanionTheory.Ar.ExRefEff.isEff_of_data} carries out this verification in the charts).
\end{proof}

\begin{proposition}[Root origins from reference geometries]\label{pf:towergeo:refroot}
Let $G$ be a reference geometry of $Y$ with an analytic realisation $\mathcal A$ whose projection $\pi$ has a
holomorphic left inverse $\iota$, assume that $Y$ has split components, and let $k,q'>0$, $s\ne0$ and $\nu$ be such
that $(kq'\nu,q',s)$ is effective for $G$. Then there are fixed data $F_x$ on $Y$ relative to the analytic system
of $\mathcal A$, a root model $R$ and a root origin $o$ of $R$ over $F_x$ with empty set of primes such that: the
generic point of $F_x.\xi_f$ is that of $\xi_0(G)$; $F_x$ has the integers $k_A=k$ and $q'$ and the weights
$kq'\nu$; the flag of $o$ has degree $1$, its realisation has $J=(1)$, is embedded and has bijective projection;
the reference multidegree of $R$ is ample; and condition \textup{(i)} of \textup{(Ex)} holds for $F_x$.
\formallabel{GlobalTheory.T.effective_scaled_weight_of_retraction}
\formallabel{GlobalTheory.CompanionTheory.Ar.RefRoot.refRoot}
\end{proposition}

\begin{proof}
Let $y_f$ be the coordinates of the block $k_0$, numbered $0,\dots,m$. Since the block embeds the generic point,
their ratios generate $\Omega_G$ (\lref{GlobalTheory.CompanionTheory.Ar.RefRoot.EmbedsGeneric.exists_closure}).
For a point $x$ of the analytic system let $\mathrm{pt}(x)\in\C^{m+1}$ be the block $k_0$ of the coordinates of the
lift of $\iota(x)$ in the logarithmic resolution. It is nonzero, since the lift lies in the complex points of the
variety, and injective up to scalars: if $\mathrm{pt}(x)=a\,\mathrm{pt}(x')$, the two lifts lie in the same torus
orbit (\lref{GlobalTheory.T.same_orbit_of_proportional_coords}: complex points whose $k_0$-blocks are
proportional lie in the same orbit, as the block embeds the generic point), hence $\iota(x)=\iota(x')$ because the
realisation is embedded, and $x=\pi\iota(x)=\pi\iota(x')=x'$. The Lie algebra model of $\towergeoHC_G$ is pushed to
a model of the Lie algebra of $Y.t$ at $\xi_0$ (\lref{GlobalTheory.T.lie_model_push_along_map}), $y_G$ is
birational (\lref{GlobalTheory.T.hodge_cpt_birational_arith_var}), and, as $Y$ has split components and the flag
has degree $1$, every element of $\Omega_G$ algebraic over $\Z$ lies in the image of the constants
(\lref{GlobalTheory.T.algebraic_constants_in_range}). These are the data of the fixed data $F_x$ (with
$k_A=k$, $q'$, weights $kq'\nu$ and $s_E=s$), of the root model $R$ (coordinates $y_G$, the lines of $\mathcal A$,
$k_A=k$, $\delta=[F_c:\Q]$) and of the root origin $o$ (the Hodge component, flag, Hodge lines and realisation of
$G$ and $\mathcal A$); the effectivity of $s$ is that of $(kq'\nu,q',s)$, and the generic point of $F_x.\xi_f$ is
that of $\xi_0(G)$ because $\xi_f$ is $\xi_G$ followed by the structure map, which factors through $\xi_0(G)$.
The reference multidegree is the unit vector of the block $k_0$, which is ample because the block embeds the
generic point (\lref{GlobalTheory.T.ample_of_single_coord_block}).

For condition (i) of (Ex), let $x$ be a point of the analytic system, $U=\iota^{-1}(U_{\iota(x)})$ for the
neighbourhood $U_{\iota(x)}$ of the lifts in the resolution, and $\lambda(x')$ the block $k_0$ of the lift at
$\iota(x)$ evaluated at $\iota(x')$. It is holomorphic on $U$ as a composition of holomorphic maps, it is a nonzero
multiple of $\mathrm{pt}(x')$ by the compatibility of the lifts of an embedded realisation, and its affine ratios
$\lambda_i/\lambda_p$ are immersions in the charts: they are immersions on the resolution
(\lref{GlobalTheory.T.holomorphic_atlas_coord_compat}), and composing with the holomorphic inverse $\iota$ of
the bijection $\pi$ preserves this (\lref{GlobalTheory.T.holomorphic_inverse_chart_value}).
\end{proof}

The last construction transports a root origin with empty set of primes to the Hodge components with primes
which are produced at the levels of the tower.

\begin{proposition}[Root origins over covers]\label{pf:towergeo:rootcov}
Assume $\towergeoG$. Let $\mathrm{Sys}$, $F_x$, $R$ and a root origin $o$ of $R$ over $F_x$ with empty set of primes
be given, such that the flag of $o$ has degree $1$, its realisation is embedded and has bijective projection. Let
$S$ be a finite set of primes, $y$ coordinates in a field $\Omega$ of characteristic $0$ which extend $R.y$
blockwise through an embedding $i\colon R.\Omega\to\Omega$ (with an injective map of blocks), $\towergeoHC$ a Hodge
component of $Y$ with primes $S$ and coordinates $y$ with a morphism $h$ to the Hodge component of $o$ which is
compatible with the points through $i$, with the maps to $Y.X$, with the constants and with the models of the Lie
algebras, and $\mathrm{Fl}$ a flag of $\towergeoHC$ with $[\Omega:R.\Omega]=\deg\mathrm{Fl}$ and root lines
$(x,e)$. If $\towergeoHC$ is $\omega$-regular and $y$ has smooth generic fibre, there are a root model $R'$ and a
root origin $o'$ of $R'$ over $F_x$ with primes $S$, whose Hodge component is $\omega$-regular with the same
function $R$, whose realisation has $J=(1)$, and such that $R'.y$ has smooth generic fibre.
\formallabel{GlobalTheory.T.root_origin_coords_over_field}
\formallabel{GlobalTheory.CompanionTheory.Ar.root_model_on_cover}
\end{proposition}

\begin{proof}
The root origin $o'$ has the Hodge component $\towergeoHC$ and the flag $\mathrm{Fl}$; we construct the root model
$R'$ with coordinates $y$ and the realisation. The coordinates $y$ define a birational arithmetic variety of
dimension $\dim\mathrm{Fl}(\Psi_H,N)+1$ (\lref{GlobalTheory.T.hodge_cpt_birational_arith_var}), so, having
smooth generic fibre, they admit an analytic logarithmic resolution $R''$ with $J=(1)$
(\lref{GlobalTheory.T.exists_analytic_log_resolution}). The coordinate ideal of $R.y$ is contained in the image
of that of $y$ under the block map, so by \lref{GlobalTheory.T.log_resolution_functorial_holomorphic} there is
a holomorphic map $h'$ from $R''$ to the resolution of $o$ mapping lifts to lifts up to the torus action. By
\lref{GlobalTheory.T.log_resolution_reindexed} and \lref{GlobalTheory.T.holomorphic_degree_composition}, the
composite $\pi\circ h'$ is holomorphic of degree $\deg\mathrm{Fl}$ (as $\pi$ is bijective and the field degree is
$\deg\mathrm{Fl}$).

\emph{Correction lines.} By \lref{GlobalTheory.T.hodge_comp_descent_first},
\lref{GlobalTheory.T.hodge_comp_descent_second} and \lref{GlobalTheory.T.hodge_comp_descent_third}, the flag
$\mathrm{Fl}$ is the pull-back of the flag of $o$ generically and on the norms (the bases correspond through the
compatibility of the Lie algebra models, and the Hodge norms are pulled back along $h$). Hence by
\lref{GlobalTheory.T.hodge_comp_descent_fourth} and \lref{GlobalTheory.T.degree_reindexed_coords}, for every
root index $r$ the ratio of the generators $e_r$ and $i(e^o_r)$, corrected by the chart monomials, is a positive
function on the complex points which is homogeneous of multidegree $x(r)-x^R(r)$ (with $x^R(r)$ extended by $0$ to
the new blocks), conjugation invariant, and locally the norm of a nonvanishing analytic function. By
\lref{GlobalTheory.T.positive_homogeneous_potential}, it is the metric of a smooth conjugation-invariant
Hermitian line $T_r$ of multidegree $x(r)-x^R(r)$ which is flat on the resolution.

\emph{The root model and the origin.} Let $R'$ have the coordinates $y$, multidegrees $x$, lines
$x'(t,r)=x^R(t,r)\otimes T_r$ (pulled back along the block map), the integer $k_A$ of $R$, the reference line of $R$
pulled back, and $\delta=[F_c:\Q]\deg\mathrm{Fl}>0$. The realisation has the resolution $R''$, $J=(1)$, projection
$\pi\circ h'$ of degree $\deg\mathrm{Fl}$, and reference lines those of $o$ tensored with the corresponding
products of the $T_r$; their potentials are the pull-backs of those of $o$ plus pluriharmonic functions
(\lref{GlobalTheory.T.log_resolution_potential_pullback} and
\lref{GlobalTheory.T.flat_on_log_resolution_corr}), and the weights of the lines differ from those of $o$ by
the same corrections, so the identities of the realisation persist. The singular metrics of $o'$ are those of $o$
pulled back and tensored with $T_r$; they realise the Hodge norms of $\mathrm{Fl}$ by
\lref{GlobalTheory.T.hodge_comp_descent_fifth}, since the interior of $\towergeoHC$ maps to the interior of the
Hodge component of $o$. The embedding of the field of $F_x$ is $i$ composed with that of $o$; its degree is
$1\cdot\deg\mathrm{Fl}$ by multiplicativity of degrees, and the identification of the block $k_0$ and the
effectivity of $s_E$ are transported along $i$. The Hodge component, and hence $\omega$-regularity and the function
$R$, are those of $\towergeoHC$.
\end{proof}

\begin{lemma}[Data on a component]\label{pf:towergeo:comp-data}
Assume $\towergeoG$ and the flat line statements \textup{(FL)}, and let $Y$ be global with split
components, $p,q\ge1$. For every generic point $x$ of $Y.X$ there are an analytic system $\mathrm{Sys}$,
fixed data $F_x$ on $Y$, a root model $R$ and a root origin $o$ of $R$ over $F_x$ with empty set of primes
such that:
\begin{enumerate}[label=(\alph*)]
\item the generic point of $F_x.\xi_f$ is $x$, and $F_x.n(\tau,i)=-k_Aq'\,i$ for all $(\tau,i)$, where
  $k_A$, $q'$ are the integers of $F_x$;
\item the flag of $o$ has degree $1$, its realisation has $J=(1)$, is embedded and has bijective
  projection $\pi$, and the reference multidegree of $R$ is ample;
\item the potentials of $\mathrm{Sys}$ for the last roots $(\tau,N-1)$ are pluriharmonic;
\item every point of $\mathrm{Sys}$ has a neighbourhood with a holomorphic lift of $F_x.\mathrm{pt}$ as in
  condition (i) of (Ex).
\end{enumerate}
\formallabel{GlobalTheory.CompanionTheory.Ar.TopTower.comp_data}
\end{lemma}

\begin{proof}
Since $Y.X$ is normal its local rings are domains; the local ring at the generic point $x$ is a field, equal
to $\kappa(x)$, and the local rings at the points of the component of $x$ inject into it. Hence the
canonical point $\Spec\kappa(x)\to Y.X$ is the generic point of a component in the sense required for
Hodge components. By \cref{pf:towergeo:refgeo} (resolution in projective coordinates is part of $\towergeoG$)
there is a reference geometry $G$ of $Y$ whose point $\xi_0$ has the generic point $x$. By \cref{pf:towergeo:holo},
$G$ has an analytic realisation $\mathcal A$ for which $\pi$ has a holomorphic left inverse. Put $\nu(\tau,i)=-i$;
it is strictly decreasing in $i$. By \cref{pf:towergeo:exeff} there are $k,q'>0$ and $s\ne0$ such that
$(kq'\nu,q',s)$ is effective for $G$. By \cref{pf:towergeo:refroot}, using that $Y$ has split components, we
obtain $F_x$, $R$, $o$ with $F_x.n=kq'\nu$ and properties (a), (b), (d). Property (c) is part of the analytic
realisation $\mathcal A$.
\end{proof}

\begin{lemma}[Data on a level]\label{pf:towergeo:level-data}
Assume $\towergeoG$ and \textup{(W)}. Let $\mathrm{Sys}$, $F_x$, $R$, $o$ be as in
\cref{pf:towergeo:comp-data}(b) (flag of degree $1$, embedded realisation, bijective $\pi$), and let $S$ be a
finite set of primes containing all prime divisors of $d_E$. Then there are a root model $R'$ and a root
origin $o'$ of $R'$ over $F_x$ with primes $S$ such that the Hodge component of $o'$ is $\omega$-regular
with $R(p)=4$ for all $p\in S$, the realisation of $o'$ has $J=(1)$, and $R'.y$ has smooth generic fibre.
\formallabel{GlobalTheory.CompanionTheory.Ar.TopTower.level_data}
\end{lemma}

\begin{proof}
By \lref{GlobalTheory.T.root_origin_good_primes_finite}, which uses $\towergeoG$ and (W), there are
coordinates $y$ in a field $\Omega$ of characteristic $0$ extending those of $R$ blockwise through an
embedding $i\colon R.\Omega\to\Omega$ of degree equal to the degree of a new flag, a Hodge component
$\towergeoHC$ of $Y$ with primes $S$ and coordinates $y$ mapping to the Hodge component of $o$ compatibly
with the points, the constants and the Lie algebra models, and a flag of $\towergeoHC$ with root lines and
compatible bases, such that $\towergeoHC$ is $\omega$-regular, $y$ has smooth generic fibre, and $R(p)=4$
for $p\in S$. By \cref{pf:towergeo:rootcov}, applied to these data, we obtain a root model $R'$ and a root origin
$o'$ over $F_x$ with primes $S$ whose Hodge component is $\omega$-regular and has the same function $R$, with
$J=(1)$ and smooth generic fibre.
\end{proof}

\begin{theorem}\label{pf:towergeo:exist}
Assume $\towergeoG$, \textup{(W)} and \textup{(FL)}. For all $E,H,D,X^o,l,c$ and every global model $Y$
of $X^o_l$ over $\OO_{F_c}$ with split components, with $p,q\ge1$, there are a root tower $T$ of type
$(\Psi_H,p+q)$ and a good origin of $T$ on $Y$ satisfying \textup{(Ex)}.
\formallabel{GlobalTheory.T.exists_good_root_tower}
\formallabel{GlobalTheory.CompanionTheory.Ar.TopTower.exists_good_tower}
\end{theorem}

\begin{proof}
Let $\mathrm{Comp}$ be the set of generic points of $Y.X$; it is finite and nonempty by
\cref{pf:towergeo:components}. For $a\in\mathrm{Comp}$ choose $\mathrm{Sys}(a)$, $F_x(a)$, $R(a)$, $o(a)$ as
in \cref{pf:towergeo:comp-data}. Since $M\ge3$ and $A\ge1$, $M'=\lcm(M,A)\ne0$, and $d_E\ne0$, so
$n_0=2|d_E|M'\ne0$. Let $\mathrm{bad}$ be the set of prime factors of $n_0$; thus $r\in\mathrm{bad}$ if and
only if $r$ is a prime dividing $2d_EM'$, and $2\in\mathrm{bad}$. For $B\in\N$ let
$S_B=\{r\le B\text{ prime}\}\cup\mathrm{bad}$, so that $r\in S_B$ if and only if $r$ is prime and $r\le B$
or $r\in\mathrm{bad}$; every prime divisor of $d_E$ lies in $\mathrm{bad}\subseteq S_B$. For each $B$ and $a$
apply \cref{pf:towergeo:level-data} to $o(a)$ and $S_B$, obtaining $R'(B,a)$ and $o'(B,a)$.

Define the root tower $T$ by $T.j=p$ (so $1\le j<N$ as $p,q\ge1$), $T.\mathrm{Comp}=\mathrm{Comp}$,
$T.\mathrm{bad}=\mathrm{bad}$, $T.\nu(\tau,i)=-i$ (a rational weight) and $T.M(B,a)=R'(B,a)$. The data
$j=p$, $\mathrm{bad}$, $S_B$, $\mathrm{Sys}(a)$, $F_x(a)$ and $\mathrm{Or}(B,a)=o'(B,a)$ form a tower
origin: the characterisations of $\mathrm{bad}$ and $S_B$ were shown above; the points $F_x(a).\xi_f$ have
generic point $a$ by \cref{pf:towergeo:comp-data}(a), so they are pairwise distinct and exhaust the generic
points of $Y.X$; and $F_x(a).n(\tau,i)=k_Aq'\cdot(-i)=k_Aq'\,\nu(\tau,i)$. It is a good origin with reference
root models $R(a)$ and root origins $o(a)$: $\omega$-regularity, $J=(1)$ and smoothness of the generic
fibres hold by \cref{pf:towergeo:level-data}, the pluriharmonicity of the last-root potentials by
\cref{pf:towergeo:comp-data}(c), and the conditions on $R(a)$, $o(a)$ by \cref{pf:towergeo:comp-data}(b).
Finally condition (i) of (Ex) is \cref{pf:towergeo:comp-data}(d), and condition (ii) holds since $R(p)=4$
for all $p\in S_B$ by \cref{pf:towergeo:level-data}.
\end{proof}

\begin{remark}\label{pf:towergeo:rem-tower}
In the tower constructed above the weight is $\nu(\tau,i)=-i$ for every component, the levels $B$ enter
only through the sets of primes $S_B$, and the reference root models $R(a)$ do not depend on $B$. The
integers $k_A$, $q'$ may depend on the component.
\end{remark}

% ---------------------------------------------------------------------------
\subsection{Chern numbers of Hodge flags}\label{pf:towergeo:sec-hodge}

In this subsection we deduce the Hodge comparison (HC) from the positivity of families (F). The heart of
the argument is the relation theorem for lattice families (\cref{pf:towergeo:relations}); from it, the
Chern number identity \eqref{pf:towergeo:eq-fcc} is obtained by summing along Frobenius orbits and by a
descent from a ramified extension. We begin with the self-contained steps.

\subsubsection*{Summation along Frobenius orbits}

\begin{lemma}[The coefficients $\towergeoco$]\label{pf:towergeo:co-series}
Let $\varphi$ be a permutation of a finite set $\Sigma$, $p\ge2$, $k\ge1$, $L\colon\Sigma\to\R$ and
$\tau\in\Sigma$. Then the series defining $\towergeoco(\varphi,p,k;\tau,\sigma)$ converge absolutely and
\[
\sum_{\sigma\in\Sigma}\towergeoco(\varphi,p,k;\tau,\sigma)L(\sigma)=\sum_{n\ge0}p^{-k(n+1)}L(\varphi^{n+1}\tau)-\frac{L(\tau)}{p-1}.
\]
\formallabel{GlobalTheory.CompanionTheory.Ar.Num.co_series_three_step}
\end{lemma}

\begin{proof}
Put $c=p^{-k}\in(0,1)$. The $n$-th term of the series for $\towergeoco(\varphi,p,k;\tau,\sigma)$ is $c^{n+1}$
or $0$, so the series is dominated by the convergent geometric series $\sum_nc^{n+1}$. Hence the finite
sum over $\sigma$ may be interchanged with the summation over $n$, and for fixed $n$,
$\sum_\sigma[\varphi^{n+1}\tau=\sigma]c^{n+1}L(\sigma)=c^{n+1}L(\varphi^{n+1}\tau)$. The term
$-[\sigma=\tau]/(p-1)$ contributes $-L(\tau)/(p-1)$.
\end{proof}

\begin{lemma}[Local comparison along orbits]\label{pf:towergeo:orbits}
Let $O$ be a discrete valuation ring, $P\subset O[x]$, $d\in\N$, $p\ge2$, $\varphi$ a permutation of a finite
set $\Sigma$, and $\mathrm{lie}$, $\mathrm{om}$ families of local pairs indexed by $\Sigma$. Assume that for
every $\tau\in\Sigma$ there are a finite set $\Sigma'$, a map $e\colon\Sigma'\to\Sigma$ with
$\tau\in e(\Sigma')$ and a permutation $\varphi'$ of $\Sigma'$ with $e\varphi'=\varphi e$ such that
$\towergeoRel(P,d;p,\varphi',\mathrm{lie}\circ e,\mathrm{om}\circ e)$ holds. Then
$\towergeoLC(P,d;p,\varphi,\mathrm{lie},\mathrm{om})$ holds.
\formallabel{GlobalTheory.T.local_comparison_of_orbits_holds}
\formallabel{GlobalTheory.CompanionTheory.Ar.Num.local_comparison_of_orbits}
\end{lemma}

\begin{proof}
Fix $k\ge1$, $Q$ with $|Q|+k=d+1$ and $\tau\in\Sigma$; choose $\Sigma',e,\varphi'$ as in the hypothesis and
$t_0\in\Sigma'$ with $e(t_0)=\tau$, and let $h\colon\Sigma'\to\R$ satisfy \eqref{pf:towergeo:eq-rel1} and
\eqref{pf:towergeo:eq-rel2} for $\mathrm{lie}\circ e$, $\mathrm{om}\circ e$, $\varphi'$. Put
$L(\sigma)=\towergeola_P(p,\mathrm{lie}_\sigma.x,k,Q)$ and $c=p^{-k}\in(0,1)$. By
\eqref{pf:towergeo:eq-rel2}, $h(t)=c\,\bigl(h(\varphi't)+L(e(\varphi't))\bigr)$ for all $t$, and by induction
on $m$,
\[
h(t)=\sum_{n<m}c^{n+1}L\bigl(e(\varphi'^{n+1}t)\bigr)+c^mh(\varphi'^mt).
\]
Since $\Sigma'$ is finite, $|h|$ and $|L\circ e|$ are bounded, so the remainder tends to $0$ and the series
converges absolutely; as $e\varphi'^n=\varphi^ne$, $h(t_0)=\sum_{n\ge0}p^{-k(n+1)}L(\varphi^{n+1}\tau)$. By
\eqref{pf:towergeo:eq-rel1} at $t_0$ and \cref{pf:towergeo:co-series},
\[
\towergeock(\mathrm{lie}_\tau,P,k,Q)+\towergeock(\mathrm{om}_\tau,P,k,Q)=h(t_0)-\frac{L(\tau)}{p-1}
=\sum_\sigma\towergeoco(\varphi,p,k;\tau,\sigma)L(\sigma).\qedhere
\]
\end{proof}

\begin{lemma}[From labels to flags]\label{pf:towergeo:flag-from-labels}
Let $O$ be a discrete valuation ring, $P$ a model over $O$ of relative dimension $d$, $p\in\N$ with
$p\ne0$ in $O$, $\Psi$ finite, $j,j'\in\N$, $\varphi$ a permutation of $\Psi\times\{1,2\}$, $T$ a local pair
over $O$ indexed by $\Psi\times\{0,\dots,j+j'-1\}$, and $\mathrm{lie}_s,\mathrm{om}_s$ local pairs indexed by
the labels $s$, of sizes $j,j'$ for $s=(\tau,1)$ and $j',j$ for $s=(\tau,2)$. Assume
\begin{enumerate}[label=(\roman*)]
\item $\mathrm{lie}_{(\tau,1)}=T\circ(\tau,i)_{i<j}$ and $\mathrm{om}_{(\tau,1)}=T\circ(\tau,j+i)_{i<j'}$;
\item $\towergeola_P(p,\mathrm{lie}_{(\tau,2)}.x,k,Q)=\towergeola_P\bigl(p,(-T.x(\tau,j+i))_{i<j'},k,Q\bigr)$
  whenever $|Q|+k=d+1$;
\item $\towergeoLC(P,d;p,\varphi,\mathrm{lie},\mathrm{om})$.
\end{enumerate}
Then $\towergeoFC_j(P,d;p,\varphi,T)$ holds.
\formallabel{GlobalTheory.T.flag_m3_from_local}
\formallabel{GlobalTheory.CompanionTheory.Ar.Num.flag_m3_of_local_m3}
\end{lemma}

\begin{proof}
Fix $\tau$, $k\ge1$, $Q$ with $|Q|+k=d+1$, and put $Z=\towergeodiv_P(p)$. We claim that for every label
$s$,
\begin{equation}\label{pf:towergeo:eq-la}
\towergeola_P(p,\mathrm{lie}_s.x,k,Q)=k\,\towergeocyc\bigl(Z,T.x,\pi_{j,k-1}(s),Q\bigr).
\end{equation}
For $s=(\sigma,1)$, $\mathrm{lie}_s.x=(T.x(\sigma,i))_{i<j}$ by (i), and
\[
\towergeola_P(p,\mathrm{lie}_s.x,k,Q)=k\,\towergeocyc\Bigl(Z,\mathrm{lie}_s.x,\sum_{i<j}X_i^{k-1},Q\Bigr),
\]
which is \eqref{pf:towergeo:eq-la} after renaming the variables. For $s=(\sigma,2)$, by (ii) the left side
equals $k\,\towergeocyc(Z,(-T.x(\sigma,j+i))_i,\sum_iX_i^{k-1},Q)$. The polynomial $\sum_iX_i^{k-1}$ is
homogeneous of degree $k-1$ with $(k-1)+|Q|=d$, and $p\ne0$ in $O$, so by
\lref{GlobalTheory.T.divisor_cycle_substitution} (cycle numbers of the divisor of a nonzero
multihomogeneous element on a model of relative dimension $d$ transform under an integral linear
substitution of the classes by the corresponding substitution of the variables in a homogeneous polynomial
of degree $d-|Q|$) this equals $k\,\towergeocyc(Z,(T.x(\sigma,j+i))_i,\sum_i(-X_i)^{k-1},Q)$, which is
\eqref{pf:towergeo:eq-la} after renaming. Now the Chern number of $T\circ(\tau,\cdot)$ is the sum of those
of $T\circ(\tau,i)_{i<j}$ and $T\circ(\tau,j+i)_{i<j'}$, so by (i), (iii), \eqref{pf:towergeo:eq-la} and the
linearity of cycle numbers,
\begin{align*}
\towergeock\bigl(T\circ(\tau,\cdot),P,k,Q\bigr)
&=\sum_s\towergeoco(\varphi,p,k;(\tau,1),s)\,k\,\towergeocyc\bigl(Z,T.x,\pi_{j,k-1}(s),Q\bigr)\\
&=\towergeocyc\bigl(Z,T.x,\mathcal B_j(\varphi,p,k,\tau),Q\bigr).\qedhere
\end{align*}
\end{proof}

\begin{lemma}[Descent through a ramified extension]\label{pf:towergeo:descent}
Let $O$ be a discrete valuation ring with uniformiser $\varpi$, $\Omega$ an $O$-algebra which is a field of
characteristic $0$, $y$ coordinates in $\Omega$ such that $P=P^O_y$ is a model over $O$ of relative
dimension $d$ and $O$ is integrally closed in $\Omega$. Let $O'$ be a discrete valuation ring which is a
local algebraic extension of $O$, and $\Omega'\supseteq\Omega$ a field containing $O'$ compatibly. Let
$p\ne0$ be a natural number and $T$ a local pair over $O$ which is vertical. If
$\towergeoFC_j(P^{O'}_y,d;p,\varphi,T\otimes O')$ holds, then $\towergeoFC_j(P,d;p,\varphi,T)$ holds.
\formallabel{GlobalTheory.T.model_descent_integral_closure}
\formallabel{GlobalTheory.CompanionTheory.Ar.Num.model_descends_from_dvr_ext}
\end{lemma}

\begin{proof}
By \lref{GlobalTheory.T.model_over_dvr_extension} (for such an extension, $P^{O'}_y$ is the extension
of $P$ and is a model over $O'$ of relative dimension $d$), $P'=P^{O'}_y$ is a model over $O'$. Since
$P\cap O=0$, the maps $O\to\Omega$ and $O\to O'$ are injective, so $\varpi\ne0$ in $O'$; write
$\varpi=u\varpi'^e$ with $u\in O'^\times$, $\varpi'$ a uniformiser of $O'$ and $e\in\N$. Here $e\ge1$:
otherwise $\varpi$ would be a unit of $O'$, hence of $O$ because $O\to O'$ is local. By
\lref{GlobalTheory.T.model_over_ramified_extension}, for every multihomogeneous $f\notin P$ and every
list $l$ of $d$ multidegrees,
\begin{equation}\label{pf:towergeo:eq-scale}
\sum_{W'}\towergeodiv_{P'}(f)(W')\,\langle l\rangle_{W'}=e\sum_W\towergeodiv_P(f)(W)\,\langle l\rangle_W .
\end{equation}
Both sides of \eqref{pf:towergeo:eq-fcc} are linear combinations of such numbers. On the left, the Chern
number of $T\circ(\tau,\cdot)$ involves the divisors of the $f_i$ and $g_i$, which are multihomogeneous and
not in $P$ because $T$ is vertical, against lists of length $(k-1)+|Q|=d$. On the right the divisor of $p$
occurs, and $p\notin P$ because $P\cap O=0$ and $p\ne0$ in $\Omega$; the polynomial
$\mathcal B_j(\varphi,p,k,\tau)$ is homogeneous of degree $k-1$, so again the lists have length $d$. By
\eqref{pf:towergeo:eq-scale} both sides of \eqref{pf:towergeo:eq-fcc} for $(P',T\otimes O')$ are $e$ times
the corresponding sides for $(P,T)$; since $e>0$, the identity for $(P',T\otimes O')$ implies that for
$(P,T)$.
\end{proof}

\subsubsection*{Lattice families}
The main input is a relation between the Chern numbers of the Lie and $\omega$ lattices of a family of
finite flat $F$-vector space schemes over a discrete valuation ring of mixed characteristic, in coordinates
in the spirit of Raynaud's description of such group schemes \cite{Raynaud1974}. A \emph{lattice family}
over $O$ in the charts of coordinates $x$ consists of Lie lattice data $D(c)$ in the Segre charts $c$ which
are families of lattices in the chart rings, an element $w\in O$ with $w^{p-1}$ a unit times $p$, $p$ a
nonunit, characters $\chi_\tau\colon F\to O$ with $\chi_\tau(0)=0$, $\chi_\tau(1)=1$, and structural
relations in every chart. \emph{Chart bundles} $\mathrm{Lie}_\tau,W_\tau$ (of rank $r_\tau$) and
$\mathrm{Om}_\tau,\mathrm{Om}^{\mathrm{HT}}_\tau$ (of rank $s_\tau$) \emph{realise} the lattice family if their
lattices in the charts are those of $D(c)$. A pair $(B,B')$ of chart bundles is a \emph{flag pair} for a
local pair $(x,x',f,g)$ and a uniformiser $\varpi$ if both carry full flags whose graded pieces are the line
bundles of multidegrees $x_i$, resp.\ $x'_i$, related generically by $f_i/g_i$ with vertical divisor.

For chart bundles $B,B'$ of rank $r$ over $O$ with the same generic fibre, the Chern number
$\towergeock(B,B',k,L,Q)$ and the root number $\towergeola(B,c,G,L,Q)$ (for $c\in O$ and $G\in\Z[X_1,\dots,X_r]$)
are computed on a \emph{flag model} of $(B,B')$, i.e.\ a refinement of the coordinates over which both bundles
acquire full flags with line bundle graded pieces, by the formulas of \cref{pf:towergeo:def-numbers} with the
multidegrees of the graded pieces, $Q(L)$ for the remaining classes and the power sum $p_k=\sum_iX_i^k$ (they
are $0$ if there is no flag model). Coordinates $x'$ \emph{refine} $x$ if they are obtained by adjoining further
blocks of coordinates in $\Omega$ and $P^O_{x'}$ is again a model over $O$ of relative dimension $d$; a lattice
family over $x$ \emph{pulls back} to the family over $x'$ whose lattices in a chart of $x'$ are generated over
its chart ring by those of the restricted chart of $x$.

\begin{theorem}[Bundle relations for lattice families]\label{pf:towergeo:ckrel}
Let $O$ be a discrete valuation ring of characteristic $0$, and $x$ coordinates in an $O$-algebra $\Omega$ which
is a field, such that $P^O_x$ is a model over $O$ of relative dimension $d$ and the ratios of $x$ generate $\Omega$
over $O$. Let $p$ be a prime and $(D,w,\chi,\varphi)$ a lattice family over $O$ in the charts of $x$, realised by
chart bundles $\mathrm{Lie},W,\mathrm{Om},\mathrm{Om}^{\mathrm{HT}}$. Then the \emph{bundle relations} hold: for
all $1\le k\le d+1$, multidegrees $L_1,\dots,L_q$ and homogeneous $Q\in\Z[X_1,\dots,X_q]$ of degree $d+1-k$
there is $h\colon\Sigma\to\R$ such that for all $\tau$
\begin{gather*}
\towergeock(\mathrm{Lie}_\tau,W_\tau,k,L,Q)+\towergeock(\mathrm{Om}_\tau,\mathrm{Om}^{\mathrm{HT}}_\tau,k,L,Q)
=h(\tau)-\frac{k\,\towergeola(\mathrm{Lie}_\tau,p,p_{k-1},L,Q)}{p-1},\\
p^kh(\tau)=h(\varphi\tau)+k\,\towergeola(\mathrm{Lie}_{\varphi\tau},p,p_{k-1},L,Q).
\end{gather*}
\formallabel{GlobalTheory.T.lattice_family_ckRelations}
\formallabel{GlobalTheory.CompanionTheory.Ar.Num.CkTop.ckRelations_of_model}
\end{theorem}

\begin{proof}
\emph{Refinements.} Let $x'$ refine $x$. The pulled-back family is a lattice family over $x'$, and the goodness
properties used below pass from $x$ to $x'$ (\lref{GlobalTheory.T.lattice_family_injective_model}); the chart
rings of $x'$ are Noetherian, contain those of the restricted charts of $x$, and the ratios of $x'$ still generate
$\Omega$ (\lref{GlobalTheory.T.chart_rings_noetherian_fraction}).

\emph{Step 1 (descent).} If the bundle relations hold for all chart bundles over a refinement $x'$ realising the
pulled-back family, they hold for all chart bundles over $x$ realising $D$. Indeed, let
$\mathrm{Lie},W,\mathrm{Om},\mathrm{Om}^{\mathrm{HT}}$ realise $D$. By
\lref{GlobalTheory.T.lattice_family_injective_model} their pull-backs to $x'$ realise the pulled-back family.
The pairs $(\mathrm{Lie}_\tau,W_\tau)$ and $(\mathrm{Om}_\tau,\mathrm{Om}^{\mathrm{HT}}_\tau)$ have the same generic
fibres, by the inclusions between their lattices with factors $1$ and $w\ne0$ which are part of the structural
relations of the family. For such pairs, \lref{GlobalTheory.T.model_over_prime_perm_module_lattice} shows that
the bundle relations for the pull-backs to a refinement with the same function field imply the bundle relations
over $x$ (a flag model of the pull-backs is a flag model of the original pairs, so the Chern numbers and root
numbers agree).

\emph{Step 2 (successive refinements).} We refine five times, keeping the previous properties by the remark on
refinements. (a) By \lref{GlobalTheory.T.lattice_family_model_facts} there is a refinement over which the
family is \emph{good} in every chart: the auxiliary lattice of functions on $F^n$ attached to the data is locally
free of rank $|F|^n$, and the images of the character maps and the vanishing and character submodules are direct
summands. (b) By \lref{GlobalTheory.T.model_refinement_exists} there is a refinement whose chart rings are
integrally closed in $\Omega$; over it, $w$ times the standard lattice is contained in the annihilator lattice of
every character in every chart (\lref{GlobalTheory.T.hodge_annihilator_lattice_holds}). (c) By
\lref{GlobalTheory.T.lattice_family_model_comparison} there is a refinement over which the family satisfies
the second goodness conditions: the annihilator submodules are direct summands of the annihilator lattices, and the
kernel, step and image lattices attached to the characters are locally free of the expected ranks. (d) By
\lref{GlobalTheory.T.lattice_family_model_chern} there is a refinement over which the character ideals of the
family are line bundles, i.e.\ generated in every chart by a fixed element times a chart monomial. (e) By
\lref{GlobalTheory.T.model_refinement_exists} once more, a final refinement has integrally closed chart rings.

\emph{Step 3 (the computation).} Over the final refinement the chart rings are Noetherian, integrally closed, with
fraction field $\Omega$, the structural facts on the lattice data hold in every chart
(\lref{GlobalTheory.T.lattice_facts_of_normal_holds}), and
\lref{GlobalTheory.T.model_over_mod_p_lattice_eigen_split} gives the bundle relations: it computes the Chern
numbers of the four bundles on the special fibre through the splitting into eigenspaces of the characters of the
lattices modulo $p$. Applying Step 1 to the five refinements of Step 2, in reverse order, gives the bundle relations
over $x$.
\end{proof}

\begin{proposition}[Relations for lattice families]\label{pf:towergeo:relations}
Let $O$ be a discrete valuation ring of characteristic $0$ with uniformiser $\varpi$, and $x$ coordinates in
an $O$-algebra $\Omega$ which is a field, such that $P^O_x$ is a model over $O$ of relative dimension $d$
and the ratios of $x$ generate $\Omega$ over $O$. Let $p$ be a prime, $F$ a finite set with elements $0,1$,
and $(D,w,\chi,\varphi)$ a lattice family over $O$ realised by chart bundles
$\mathrm{Lie},W,\mathrm{Om},\mathrm{Om}^{\mathrm{HT}}$. If $(\mathrm{Lie}_\tau,W_\tau)$ and
$(\mathrm{Om}_\tau,\mathrm{Om}^{\mathrm{HT}}_\tau)$ are flag pairs for local pairs $\mathrm{lie}_\tau$ and
$\mathrm{om}_\tau$, then $\towergeoRel(P^O_x,d;p,\varphi,\mathrm{lie},\mathrm{om})$ holds.
\formallabel{GlobalTheory.T.arithmetic_relations_statement_holds}
\formallabel{GlobalTheory.CompanionTheory.Ar.Num.TopRel.relations}
\end{proposition}

\begin{proof}
By \cref{pf:towergeo:ckrel}, the chart bundles satisfy the bundle relations. Next,
$w\ne0$: otherwise $0=w^{p-1}=up$ with $u$ a unit, contradicting $p\ne0$ in $O$. In every chart, the
lattices of $\mathrm{Lie}_\tau$ and $W_\tau$ (resp.\ of $\mathrm{Om}_\tau$ and $\mathrm{Om}^{\mathrm{HT}}_\tau$)
have the same generic span: by the inclusions which are part of the structural relations of the lattice
family, the product of every element of one of them with $1$ or with $w$ lies in the other. Under this
condition and $p\ne0$ in $O$, \lref{GlobalTheory.T.uniformizer_model_perm_module_lattice} (the bundle
relations for chart bundles which are flag pairs for local pairs $\mathrm{lie}_\tau$, $\mathrm{om}_\tau$
imply the relations for these local pairs, because the Chern numbers and root numbers of flag pairs of
chart bundles equal those of the local pairs) gives $\towergeoRel(P^O_x,d;p,\varphi,\mathrm{lie},\mathrm{om})$.
\end{proof}

\begin{remark}\label{pf:towergeo:rem-relations}
\cref{pf:towergeo:ckrel} carries most of the weight of (HC). Its final step,
\lref{GlobalTheory.T.model_over_mod_p_lattice_eigen_split} (about $4{,}800$ lines in \cite{LeanSource2026}),
computes the Chern numbers of the four bundles on the special fibre from the structure of the $F$-vector space
schemes. The relations
\eqref{pf:towergeo:eq-rel1}--\eqref{pf:towergeo:eq-rel2} express that the defect $h$ is transported by
Frobenius with the factor $p^k$, which is how the coefficients $\towergeoco$ arise in
\cref{pf:towergeo:orbits}.
\end{remark}

The chart bundles to which \cref{pf:towergeo:relations} is applied are read off from lattices on Hodge
components. Let $\xi\colon\Spec\Omega\to X$ be a point of a scheme $X$ of which every point is a specialisation and
whose local rings are given by the Segre coordinates of $y$, $V$ an $\Omega$-vector space, $b_1,\dots,b_r\in V$
linearly independent, and $Q=(x,x',f,g)$ a local pair over $\Z$ indexed by $\{1,\dots,r\}$. For a set $T$ of points,
lattices $L,L'$ in $V$ (submodules $L_q$ over the local rings $\mathcal O_q$, $q\in X$) form a \emph{lattice flag
pair} on $T$ with $b$ and $Q$ if on $T$ both are flagged by $b$ (the partial spans of $b$ give full flags of
subbundles), their graded pieces are the line bundles of multidegrees $x_i$, resp.\ $x'_i$, with generic generators
$e_i,e'_i$, and $f_i(y)=(e'_i/e_i)\,g_i(y)$ with $g_i(y)\ne0$. The \emph{coordinate lattice} of $L$ over $O$ in the
chart $c$ is the set of $z\in\Omega^r$ with $\sum_iz_ib_i\in L_q$ for all points $q$ of the chart $c$ whose local
ring contains $O$.

\begin{proposition}[Chart bundles of lattice flag pairs]\label{pf:towergeo:dedekind}
Let $y$ define an arithmetic variety of dimension $d+1$ and $\xi$, $X$ be as above, $A\subseteq\Omega$ a Dedekind
domain finite over $\Z$ containing the constants of $y$, $p$ a prime, $\mathfrak m\ni p$ a prime of $A$,
$O=A_{\mathfrak m}$ with uniformiser $\varpi$. Let $(L,L')$ be a lattice flag pair on the points over $p$ with $b$
and $Q$, such that $L_q=L'_q$ at the points $q$ over $p$ outside the fibre over $p$. Then there are chart bundles
$B,B'$ over $O$ whose lattices in every chart are the coordinate lattices of $L$, resp.\ $L'$, and which form a
flag pair for $Q\otimes O$ and $\varpi$.
\formallabel{GlobalTheory.T.arith_variety_dedekind_constants}
\formallabel{GlobalTheory.CompanionTheory.Ar.CoordRead.flagPair_coord}
\end{proposition}

\begin{proof}
\emph{Reduction to the standard basis.} The map $c_b\colon\Omega^r\to V$, $z\mapsto\sum_iz_ib_i$, is injective.
The preimages $c_b^{-1}(L)$, $c_b^{-1}(L')$ are lattices in $\Omega^r$ which form a lattice flag pair with the
standard basis and the same $Q$: $c_b$ maps the partial spans of the standard basis onto those of $b$, so the flags
and their graded pieces correspond. They agree where $L$ and $L'$ agree, and their coordinate lattices are those of
$L$ and $L'$. Hence we may assume that $V=\Omega^r$ and $b$ is the standard basis.

\emph{The bundles.} By \lref{GlobalTheory.T.dedekind_constants_degree_one}, applied to $L$ and to $L'$, there
are chart bundles $B,B'$ over $O$ whose lattices in every chart are the coordinate lattices of $L$, resp.\ $L'$,
and which carry the full flags given by the bases $v=(e_ib_i)_i$, resp.\ $v'=(e'_ib_i)_i$, with graded pieces the line
bundles of multidegrees $x_i$, resp.\ $x'_i$. By \lref{GlobalTheory.T.localization_model_over_place}, $P^O_y$ is
a model over $O$ of relative dimension $d$.

\emph{Same generic fibre.} A point $q$ whose local ring contains $O$ lies over $p$, since every prime $p'\ne p$ is a
unit in $O$; if moreover $\varpi$ is a unit at $q$, then $q$ is not in the fibre over $p$, because $p$ is a unit times
a power of $\varpi$ in $O$. At such points $L_q=L'_q$. Now let $z$ lie in the lattice of $B$ in the chart $c$. Then
$z\in L'_q$ at all points $q$ of the chart over $O$ at which $\varpi$ is a unit, and a regular function at these
points becomes regular in the chart after multiplication by a power of $\varpi$; by the flag structure of $B'$
(\lref{GlobalTheory.CompanionTheory.Ar.Num.LatFlag.flag_sections}), $\varpi^nz$ lies in the lattice of $B'$ for
some $n$. By symmetry $B$ and $B'$ have the same generic fibre.

\emph{Flag pair.} The change of basis from $v$ to $v'$ is diagonal with entries $e'_i/e_i$, and
$f_i(y)=(e'_i/e_i)g_i(y)$ with $f_i,g_i$ multihomogeneous of degrees differing by $x_i-x'_i$. At a point of a chart
$c$ over $O$ at which $\varpi$ is a unit, the graded pieces of $L$ and $L'$ agree, so the ratio
$y_c^{x'_i}e'_i/(y_c^{x_i}e_i)$ of their local generators is a unit there; hence $f_i/g_i$, written in the chart, is
a unit at every relevant prime of the chart ring not containing $\varpi$, i.e.\ its divisor is vertical. Hence
$(B,B')$ is a flag pair for $Q\otimes O$ and $\varpi$.
\end{proof}

\subsubsection*{The comparison}

\begin{theorem}\label{pf:towergeo:hodge}
The positivity of families \textup{(F)} implies the Hodge comparison \textup{(HC)}.
\formallabel{GlobalTheory.T.model_on_hodge_from_family}
\formallabel{GlobalTheory.CompanionTheory.Ar.Num.HodgeModel.model_on_hodge_of_family_pos}
\end{theorem}

\begin{proof}
Let the data of (HC) be given: $\towergeoHC$, $d$, a flag $\mathrm{Fl}$ with two levels $(S_0,T)$,
$p\in S\setminus S_0$ Hodge--Tate free with $p\ne2$, $p\nmid d_E$, a Frobenius $\varphi$ at $p$, and a place
$O$ of $\Omega$ over $p$, the localisation of $A$ at $\mathfrak m$, with uniformiser $\varpi$. Write
$P=P^O_y$ and $T_O=T\otimes O$.

\emph{Step 0.} If $N=0$, both sides of \eqref{pf:towergeo:eq-fcc} vanish: the local pair
$T\circ(\tau,\cdot)$ has no indices, and $\pi_{j,k-1}(s)=0$ for all labels, so $\mathcal B_j=0$. Assume
$N>0$. Then $p$ is prime (as $p\in S$) and $\Omega$ has characteristic $0$.

\emph{Step 1 (the model over $O$).} By \lref{GlobalTheory.T.arith_variety_coord_ideal}, $A$ is a
Dedekind domain, finite over $\Z$, contained in $\Omega$ and containing the constants of $y$, and for every
place $O$, $O$ is integrally closed in $\Omega$ and the ratios of $y$ generate $\Omega$ over $O$. By
\lref{GlobalTheory.T.localization_model_over_place}, $P$ is the extension of $P^A_y$, it is a model
over $O$ of relative dimension $d$, and every local pair over $\Z$ whose divisors are supported over a finite
set of primes, with $f_i/g_i$ a unit elsewhere, becomes vertical over $O$. Since $\mathrm{Fl}$ has two levels
$(S_0,T)$, this applies to $T$ with the primes of $S\setminus S_0$; hence $T_O$ is vertical.

\emph{Step 2 (families).} The bases $b_s$ of $\mathrm{Fl}$ are linearly independent and span
$V^{\mathrm{gen}}_s$: they are linearly independent by the definition of a flag, they lie in
$V^{\mathrm{gen}}_s$, and their number is the rank of the vector bundle $V_s$, whose generic fibre is
$V^{\mathrm{gen}}_s$. Apply (F) with these bases at the place $O$: for every label $t$ we obtain $\Sigma'_t$,
$\epsilon_t$, $\varphi'_t$, a finite field $F_t$ of characteristic $p$, $w_t\in O$, $\sigma_t$ and Lie lattice
data $D_{f,t}(c)$ forming a pre-family related to $\mathrm{dR}$ and $\mathrm{HT}_p$ as in
\cref{pf:towergeo:def-fam}.

\emph{Step 3 (roots of unity).} Let $n=\prod_t(|F_t|-1)>0$. By
\lref{GlobalTheory.T.exists_dvr_extension_field} there are a discrete valuation ring $O'$, a local
algebraic extension of $O$, a field $\Omega'\supseteq\Omega$ containing $O'$, and a primitive $n$-th root of
unity $\zeta\in O'$ with $\Omega'=\Omega(\zeta)$. Then $\zeta^{n/(|F_t|-1)}$ is a primitive
$(|F_t|-1)$-th root of unity. Let $\varpi'$ be a uniformiser of $O'$, $P'=P^{O'}_y$ and $T_{O'}=T\otimes O'$.
By \lref{GlobalTheory.T.model_over_dvr_extension}, $P'$ is the extension of $P$ and a model over $O'$ of
relative dimension $d$. The ratios of $y$ generate $\Omega'$ over $O'$: the subfield of $\Omega'$ generated by
$O'$ and the ratios contains the image of $O$ and the ratios, hence $\Omega$ (which they generate), and it
contains $\zeta\in O'$, hence $\Omega(\zeta)=\Omega'$.

\emph{Step 4 (lattice flag pairs).} By \lref{GlobalTheory.T.hodge_flag_outside_prime_fibres},
$\mathrm{HT}_{p,s}=\mathrm{dR}_s$ at the points outside the fibres over all primes; in particular the two
lattices agree at the points over $p$ outside the fibre over $p$. For a local pair $T'=(x,x',f,g)$ indexed
by $\{0,\dots,r-1\}$ put $T'^\dagger=(-x\circ\rho,-x'\circ\rho,g\circ\rho,f\circ\rho)$, $\rho$ the reversal of
the indices; $\dagger$ is an involution. By \lref{GlobalTheory.T.hodge_flag_cm_type_degree}, for every
label $s$ the pair $(\mathrm{dR}_s,\mathrm{HT}_{p,s})$ is a \emph{lattice flag pair} over $p$ with the bases
$b_s$ and the local pair $\lambda_s$ defined by
\[
\lambda_{(\tau,1)}=T\circ(\tau,i)_{i<p_D},\qquad\lambda_{(\tau,2)}=\bigl(T\circ(\tau,p_D+i)_{i<q_D}\bigr)^\dagger:
\]
both lattices are flagged by $b_s$, their graded pieces are line bundles of multidegrees $\lambda_s.x$ and
$\lambda_s.x'$, and $\lambda_s.f(y)=(e'/e)\lambda_s.g(y)$ for the generic generators. Put
$\omega_s=\lambda_{s^\vee}^\dagger$.

\emph{Step 5 (chart bundles).} By \cref{pf:towergeo:dedekind}, for every $s$
there are chart bundles $B_s,B'_s$ over $O$, whose lattices in each chart are the coordinates of
$\mathrm{dR}_s$ and $\mathrm{HT}_{p,s}$ in the basis $b_s$, forming a flag pair for $\lambda_s\otimes O$. By
\lref{GlobalTheory.T.flag_pair_chart_bundle_compat} (the duals of a flag pair form a flag pair for the
pair transformed by $\dagger$), the dual chart bundles $C_s,C'_s$ of $B_{s^\vee},B'_{s^\vee}$ form a flag pair
for $\omega_s\otimes O$. By \lref{GlobalTheory.T.model_descent_valuation_fibre} (the spans over the
chart rings of $O'$ of the base change of a flag pair form a flag pair for $\varpi'$), the $O'$-spans
$LB_s,LB'_s,LC_s,LC'_s$ of the base changes are flag pairs for $\lambda_s\otimes O'$ and $\omega_s\otimes O'$.

\emph{Step 6 (relations on orbits).} Fix a label $t$. By
\lref{GlobalTheory.T.model_descent_generic_fibre}, applied to the pre-family of Step 2 and the root of
unity of Step 3, there are characters $\chi$ and a lattice family $D'_{f,t}$ over $O'$ whose lattices are the
$O'$-spans of the base changes of those of $D_{f,t}$. By the relation of $D_{f,t}$ with $\mathrm{dR}$ and
$\mathrm{HT}_p$ and Step 5, the chart bundles $LB_{\epsilon_t\tau}$, $LB'_{\epsilon_t\tau}$,
$LC_{\epsilon_t\tau}$, $LC'_{\epsilon_t\tau}$ ($\tau\in\Sigma'_t$) realise $D'_{f,t}$, and
\cref{pf:towergeo:relations} gives
$\towergeoRel(P',d;p,\varphi'_t,(\lambda_{\epsilon_t\tau}\otimes O')_\tau,(\omega_{\epsilon_t\tau}\otimes O')_\tau)$.

\emph{Step 7 (labels).} Since $\epsilon_t\varphi'_t=\varphi\epsilon_t$, $t\in\epsilon_t(\Sigma'_t)$ and
$p\ge2$, \cref{pf:towergeo:orbits} gives $\towergeoLC(P',d;p,\varphi,\lambda\otimes O',\omega\otimes O')$.

\emph{Step 8 (flags).} Apply \cref{pf:towergeo:flag-from-labels} over $O'$ with $j=p_D$, $j'=q_D$,
$\mathrm{lie}=\lambda\otimes O'$ and $\mathrm{om}=\omega\otimes O'$: condition (i) holds by the definition
of $\lambda$ and because $\dagger$ is an involution, so that
$\omega_{(\tau,1)}=\lambda_{(\tau,2)}^\dagger=T\circ(\tau,p_D+i)_{i<q_D}$; condition (ii) holds because
$\lambda_{(\tau,2)}.x$ is the reversal of $(-T.x(\tau,p_D+i))_i$ and root numbers are invariant under
permutations of the indices; condition (iii) is Step 7. As $p\ne0$ in $O'$,
$\towergeoFC_{p_D}(P',d;p,\varphi,T_{O'})$ holds.

\emph{Step 9 (descent).} By \cref{pf:towergeo:descent}, applied with Steps 1 and 3, Step 8 gives
$\towergeoFC_{p_D}(P,d;p,\varphi,T_O)$.
\end{proof}

% ---------------------------------------------------------------------------
\subsection{Nefness of root lines on the special fibres}\label{pf:towergeo:sec-nef}

This subsection proves two nefness statements for the root lines of root models which come from Hodge
components: at the primes of $S$ at which the Hodge component is Hodge--Tate free, from the trivial flags of
\cref{pf:towergeo:sec-triv}, and at the primes outside $S$, from a statement on semi-abelian schemes over curves
in characteristic $p>2$. They are used in the proof of the volume theorem and in \cref{pf:towergeo:sec-nine}.

Let $M$ be a root model of type $(\Psi,N)$. A \emph{weight} is a function $\lambda\colon\Psi\times\{0,\dots,N-1\}\to\R$;
it is \emph{dominant} if $\lambda(\tau,b)\le\lambda(\tau,a)$ for $a\le b$; $\|\lambda\|$ is its supremum norm, and
$\mathrm{gap}_j(\lambda,\tau)=\lambda(\tau,j-1)-\lambda(\tau,j)$ for $1\le j<N$. For $b\in\R$, $\mathcal N_j(b,\lambda)$
means $4|\Psi|\,\|\lambda\|\le(b-2)\,\mathrm{gap}_j(\lambda,\tau)$ for all $\tau$. The lines of $M$ are
$\ell_t(r)=x(t,r)$ for the root indices $r$ and $\ell_t(\infty)=A$, with coefficients
$\mathrm{coef}(\lambda,s)(r)=-\lambda(r)$ and $\mathrm{coef}(\lambda,s)(\infty)=s/k_A$. A \emph{vertical curve of
characteristic $p$} is a relevant multihomogeneous prime $W\supseteq P_M$ with $W\cap\Z=p\Z$ and
$\dim\Z[x]/W=1+|K|$; $\deg(a\cdot W)$ denotes the degree on $W$ of the line bundle of multidegree $a$, which is
additive in $a$. As in the proof of the volume theorem, $\mathrm{Nef}_p(\lambda,s)$ means
$\sum_r\mathrm{coef}(\lambda,s)(r)\deg(\ell_1(r)\cdot W)\ge0$ for every vertical curve $W$ of characteristic $p$.
A multidegree $a$ is \emph{nef on the fibres} of $P_M$ if $\deg(a\cdot W)\ge0$ for all vertical curves $W$ of all
characteristics.

\begin{lemma}[Abel summation]\label{pf:towergeo:abel}
Let $\lambda_0\ge\lambda_1\ge\dots\ge\lambda_{N-1}$ and $x_0,\dots,x_{N-1}$ be real numbers with
$X_i=\sum_{l<i}x_l\ge0$ for $0\le i\le N$ and $X_N=0$. Then $\sum_l\lambda_lx_l\ge0$.
\formallabel{PositivityTheory.HtPt.Gf.abel_sum}
\end{lemma}

\begin{proof}
Since $x_l=X_{l+1}-X_l$ and $X_0=X_N=0$,
$\sum_{l=0}^{N-1}\lambda_lx_l=\sum_{i=1}^{N-1}(\lambda_{i-1}-\lambda_i)X_i\ge0$.
\end{proof}

\begin{lemma}[Trivial flags are nef]\label{pf:towergeo:trivnef}
Let $M$ be a root model of type $(\Psi,N)$ with $M.\Omega$ of characteristic $0$, $p$ a prime, and $o\subseteq M.\Omega$
the ring of fractions $a/b$ with $a\in\Z$, $b\in\N$, $p\nmid b$. If $\towergeoTriv_p(o;M.y,(x(\tau,i))_i)$ holds for
every $\tau$, then $M$ satisfies $\mathrm{Nef}_p(\mu,0)$ for every dominant $\mu$.
\formallabel{PositivityTheory.HtPt.Gf.trivflag_nef}
\end{lemma}

\begin{proof}
Let $W$ be a vertical curve of characteristic $p$ and $\tau\in\Psi$. By
\lref{PositivityTheory.HtPt.Gf.trivflag_deg}, the trivial flag modulo $p$ restricts to a flag of the trivial
bundle of rank $N$ on $W$ whose graded pieces are the line bundles of multidegrees $x(\tau,i)$, so that
$\sum_{l<i}\deg(x(\tau,l)\cdot W)\le0$ for $0\le i\le N$ (a subbundle of a trivial bundle on a curve has
nonpositive degree) and $\sum_{l<N}\deg(x(\tau,l)\cdot W)=0$. \Cref{pf:towergeo:abel}, applied to
$\lambda_l=\mu(\tau,l)$ and $x_l=-\deg(x(\tau,l)\cdot W)$, gives $\sum_l-\mu(\tau,l)\deg(x(\tau,l)\cdot W)\ge0$.
Summing over $\tau$, and noting that the reference line has coefficient $0/k_A=0$, gives $\mathrm{Nef}_p(\mu,0)$
at $W$.
\end{proof}

\begin{proposition}[Nefness at Hodge--Tate free primes]\label{pf:towergeo:rootw}
Assume the Hodge--Tate statement \textup{(HT)}. Let $Y$ be a model of $X^o$ over a ring $O$, $S$ a finite set of
primes, $M$ a root model of type $(\Psi_H,N)$, $\towergeoHC$ a Hodge component of $Y$ with primes $S$ and
coordinates $M.y$ and $\mathrm{Fl}$ a flag of $\towergeoHC$ with root lines $(M.x,e)$. If $p\in S$ and $\towergeoHC$
is Hodge--Tate free at $p$, then $M$ satisfies $\mathrm{Nef}_p(\mu,0)$ for every dominant $\mu$.
\formallabel{GlobalTheory.T.root_lines_weight_bound_of_free}
\formallabel{GlobalTheory.CompanionTheory.Ar.HT.rootModel_nefAt_of_hodgeCpt}
\end{proposition}

\begin{proof}
If $N=0$, the only line is the reference line, with coefficient $0/k_A=0$, and the assertion is trivial. Let
$N>0$. The prime $p$ is prime as $p\in S$, and $M.\Omega$ has characteristic $0$ since it contains the image of the
constants $F_c$. Let $o\subseteq M.\Omega$ be the ring of fractions $a/b$ with $p\nmid b$. By
\cref{pf:towergeo:htcharts}, $\towergeoTriv_{p^{R(p)-2}}(o;M.y,(x(\tau,i))_i)$ holds for every $\tau$. Since
$R(p)\ge4$, $p$ divides $p^{R(p)-2}$, so $\towergeoTriv_p(o;M.y,(x(\tau,i))_i)$ holds by the remark after
\cref{pf:towergeo:def-trivmod}, and \cref{pf:towergeo:trivnef} gives the assertion.
\end{proof}

At the primes outside $S$ the input is a statement on curves. A geometric fibre of a semi-abelian scheme is
\emph{power-free} if it is an extension of an abelian variety $B$ by a torus such that the pull-back of the line
bundle of the semi-abelian scheme is isomorphic to the pull-back of a line bundle on $B$ with trivial group
$K(\cdot)$. For a field $K$ of characteristic $p$, a \emph{curve} over $K$ is a complete regular curve given by its
places, and for a model $V$ over $K$ of the Lie algebra of a semi-abelian scheme $t$ along a curve $C$ and a subspace
$U\subseteq V$, $d(U)$ denotes the degree on $C$ of the saturation of $U$ in the Lie algebra lattice of $t$ along
$C$. A family $\chi$ of ring homomorphisms $\OO_E\to k$, indexed by $\Psi\sqcup\Psi$, is \emph{split} if its members
are pairwise distinct.

\begin{definition}[Curve statement]\label{pf:towergeo:def-curve}
For a prime $p$, integers $j,q\ge1$ and $\mu\colon\Psi\times\N\to\R$, the \emph{gap condition} $\mathrm{GC}_p(j,q,\mu)$
means: $\mu(\tau,i+1)\le\mu(\tau,i)$ for $1\le i<j+q$, and there are $C,g\in\R$ with $|\mu(\tau,i)|\le C$ for
$1\le i\le j+q$, $g\le\mu(\tau,j)-\mu(\tau,j+1)$ for all $\tau$, and $4|\Psi|\,C\le(p-2)g$. The \emph{curve
statement} \textup{(CC)} asserts: for every host $H$ on a CM field $E$, semi-abelian scheme $t$ with
$\OO_E$-action over a scheme $X_0$ whose power maps are bijective on global functions, fields $k\subseteq K$ with
$K$ of characteristic $p>2$, curve $C$ over $K$ with a morphism to $X_0$ at whose places the fibres of $t$ are
power-free and on which the elements of $k$ have valuation $\le1$, model $V$ of the Lie algebra of $t$ at the
generic point of $C$ with the action of $\OO_E$, split family $\chi$ with $\chi_{(\tau,2)}=\chi_{(\tau,1)}\circ c$,
$j,q\ge1$, chains of subspaces $0=W_\tau(0)\subseteq\dots\subseteq W_\tau(j)=V_{\chi(\tau,1)}$ and
$0=L_\tau(0)\subseteq\dots\subseteq L_\tau(q)=V_{\chi(\tau,2)}$ of the eigenspaces, and $\mu$ with
$\mathrm{GC}_p(j,q,\mu)$:
\begin{multline}\label{pf:towergeo:eq-curve}
\sum_\tau\Bigl(\sum_{i<j}\mu(\tau,i+1)\bigl(d(W_\tau(i+1))-d(W_\tau(i))\bigr)\\
+\sum_{m<q}\mu(\tau,j+m+1)\bigl(d(L_\tau(q-m-1))-d(L_\tau(q-m))\bigr)\Bigr)\le0.
\end{multline}
\formal{GlobalTheory.CompanionTheory.Ar.Conj.CurveSemiAbelianConjecture}
\end{definition}

The curve statement holds by \lref{GlobalTheory.T.curveSemiAbelianConjecture}. It is applied through de Rham
ties. Let $M$ be a root model of type $(\Psi,N)$, $1\le j\le N$ and $p$ a prime. A \emph{de Rham tie} of $M$ at
$(j,p)$ consists of: a scheme $X$ with a point $\xi\colon\Spec M.\Omega\to X$ of which every point is a
specialisation and whose local rings are given by the Segre coordinates of $M.y$; a semi-abelian scheme $t$ with
$\OO_E$-action over $X$ whose power maps are bijective on global functions and whose fibres at the points in the
fibre over $p$ are power-free; a model $V$ of its Lie algebra at $\xi$ with the action $r$ of $\OO_E$; pairwise
distinct embeddings $\sigma_s\colon\OO_E\to M.\Omega$ ($s\in\Psi\sqcup\Psi$) with $2|\Psi|=[E:\Q]$,
$\sigma_{(\tau,2)}=\sigma_{(\tau,1)}\circ c$ and values regular at the points over $p$; the condition $p\nmid d_E$; and
bases $b_1(\tau)$ of size $j$ and $b_2(\tau)$ of size $N-j$ of the $\sigma_{(\tau,1)}$- and
$\sigma_{(\tau,2)}$-eigenspaces of $V$ which define full flags of the eigenlattices of the Lie algebra at the points
in the fibre over $p$, with graded pieces the line bundles of multidegree $x(\tau,i)$ at position $i<j$, resp.\
$-x(\tau,i)$ at position $N-1-i$ for $i\ge j$.

\begin{lemma}[Reduction to curves]\label{pf:towergeo:m2curve}
Assume \textup{(CC)}. Let $M$ be a root model of type $(\Psi,N)$, $1\le j<N$, $p>2$ a prime, and assume that $M$
has a de Rham tie at $(j,p)$. Then $M$ satisfies $\mathrm{Nef}_p(\mu,0)$ for every dominant $\mu$ with
$\mathcal N_j(p,\mu)$.
\formallabel{GlobalTheory.T.rootLineDeRhamModel_of_curveSemiAbelian}
\formallabel{GlobalTheory.CompanionTheory.Ar.Conj.m2Model_of_m2Curve}
\end{lemma}

\begin{proof}
Let $W$ be a vertical curve of characteristic $p$; $M.\Omega$ has characteristic $0$ as $M.y$ defines an arithmetic
variety.

\emph{The curve.} By \lref{GlobalTheory.T.function_field_model_exists} there are a function field $K$ of
transcendence degree one over $\F_p$ and coordinates $w$ in $K$, generating $K$ and with a nonzero entry in every
block, whose coordinate ideal is $W$. By \lref{GlobalTheory.T.function_field_curve_exists}, $K$ is the function
field of a curve $C$ (with at least one place) on which the elements algebraic over $\F_p$ have valuation $\le1$.
Since $P_M\subseteq W$, \lref{GlobalTheory.T.curve_through_coordinates} gives a $K$-point $z$ of $X$ with
coordinates $w$ and a curve in $X$ along $C$ through $z$, all of whose points lie in the fibre over $p$ (as $p=0$
in $K$).

\emph{Characters and flags.} The values of the $\sigma_s$ at $z$ are regular; reducing them gives, by
\lref{GlobalTheory.T.split_characters_unramified} (as $p\nmid d_E$ and there are $[E:\Q]$ pairwise distinct
$\sigma_s$), a split family $\chi_s\colon\OO_E\to\bar\F_p\cap K$ with $\chi_{(\tau,2)}=\chi_{(\tau,1)}\circ c$; the
elements of $\bar\F_p\cap K$ have valuation $\le1$ on $C$. Let $V_z$ be the Lie algebra of $t$ at $z$, a
finite-dimensional $K$-vector space with an action of $\OO_E$. The flags of the tie at the points of the curve give
chains $W_\tau(i)$ ($i\le j$) and $L_\tau(i)$ ($i\le N-j$) of subspaces of the $\chi$-eigenspaces of $V_z$, from
$0$ to the eigenspace, such that the saturations have the degrees of the graded pieces: by
\lref{GlobalTheory.T.curve_lattice_degree_formula} the degree along $C$ of a line bundle of multidegree $a$ on
the fibre over $p$ is the coordinate degree of $w$ in $a$, degrees along $C$ are unique
(\lref{GlobalTheory.T.curve_degree_unique}), and the coordinate degree equals $\deg(a\cdot W)$
(\lref{GlobalTheory.T.function_field_degree_eq}). Writing $x_r=\deg(x(r)\cdot W)$ we get
\[
x_{(\tau,i)}=d(W_\tau(i+1))-d(W_\tau(i))\quad(i<j),\qquad
x_{(\tau,i)}=d(L_\tau(N-1-i))-d(L_\tau(N-i))\quad(i\ge j).
\]

\emph{The weights.} Put $\mu'(\tau,i)=\mu(\tau,i-1)$ for $1\le i\le N$ (and $\mu'(\tau,i)=\mu(\tau,N-1)$ for larger
$i$). Then $\mathrm{GC}_p(j,N-j,\mu')$ holds: the monotonicity is the dominance of $\mu$, $|\mu'|\le C=\|\mu\|$, and
with $g=\min_\tau\mathrm{gap}_j(\mu,\tau)=\min_\tau(\mu'(\tau,j)-\mu'(\tau,j+1))$ the inequality $4|\Psi|C\le(p-2)g$
is $\mathcal N_j(p,\mu)$. Substituting the degrees, the left side of \eqref{pf:towergeo:eq-curve} with $q=N-j$ is
\[
\sum_\tau\Bigl(\sum_{i<j}\mu(\tau,i)x_{(\tau,i)}+\sum_{m<N-j}\mu(\tau,j+m)x_{(\tau,j+m)}\Bigr)=\sum_r\mu(r)x_r
\]
By (CC) it is $\le0$, i.e.\
$\sum_r-\mu(r)\deg(x(r)\cdot W)\ge0$; as the reference line has coefficient $0$, this is $\mathrm{Nef}_p(\mu,0)$ at
$W$.
\end{proof}

\begin{proposition}[Nefness at good primes]\label{pf:towergeo:nefgood}
Let $X^o$ be a family of moduli spaces, $l$ a level which is good at every prime, and $Y$ a model of $X^o_l$ over
$\OO_{F_c}$ which dominates compatible systems of local models over the valuation rings of $\C$ above all primes
whose fibres are split and satisfy the further condition required in the definition of global models (for
instance, $Y$ global). Let $p_D,q_D\ge1$, $S$ a finite set of primes, $M$ a root model of type $(\Psi_H,N)$,
$\towergeoHC$ a Hodge component of $Y$ with primes $S$ and coordinates $M.y$, and $\mathrm{Fl}$ a flag with root lines
$(M.x,e)$. Let $p>2$ be a prime with $p\notin S$ and $p\nmid d_E$. Then $M$ satisfies $\mathrm{Nef}_p(\mu,0)$ for every
dominant $\mu$ with $\mathcal N_{p_D}(p,\mu)$.
\formallabel{GlobalTheory.T.hodge_cpt_global_nef_bound}
\formallabel{GlobalTheory.CompanionTheory.Ar.Conj.hodge_cpt_global_nefness}
\end{proposition}

\begin{proof}
By \lref{GlobalTheory.T.global_model_power_free}, the domination of the local models implies that the power
maps of $Y.t$ are bijective on global functions and that its geometric fibres at all points of characteristic $p$
are power-free. Hence, as $p\notin S$ and $p\nmid d_E$, \lref{GlobalTheory.T.root_lines_unramified_primes}
provides a de Rham tie of $M$ at $(p_D,p)$, whose data are taken from $\towergeoHC$ and the flag (outside $S$ the
patched lattices of $\towergeoHC$ are its de Rham lattices). As $1\le p_D<N$ and $p>2$,
\cref{pf:towergeo:m2curve} and the curve statement \textup{(CC)} give the assertion.
\end{proof}

\begin{lemma}[Nefness of the block sums]\label{pf:towergeo:blocknef}
Assume $\towergeoG$, let $Y$ be global, $p_D,q_D\ge1$, $S$ a finite set of primes, $M$ a root model of type
$(\Psi_H,N)$, $\towergeoHC$ a Hodge component of $Y$ with primes $S$ and coordinates $M.y$ which is Hodge--Tate free
at all $p\in S$, and $\mathrm{Fl}$ a flag with root lines $(M.x,e)$. Assume that every prime $r\notin S$ satisfies
$r\nmid d_E$, $r\ge4f+2$ and $r>2$. Then the multidegrees $\sum_\tau\sum_{i\ge p_D}x(\tau,i)$ and
$-\sum_\tau\sum_{i<p_D}x(\tau,i)$ are nef on the fibres of $P_M$.
\formallabel{GlobalTheory.T.nefOnFibres_block_xdeg_sums_of_rootLines}
\formallabel{GlobalTheory.CompanionTheory.Relay.nefOnFibres_block_xdeg_sums_of_rootLines}
\end{lemma}

\begin{proof}
Let $\mu$ be dominant with $|\mu(\tau,i)|\le1$ and $\mathrm{gap}_{p_D}(\mu,\tau)=1$ for all $\tau$, and $r$ a prime.
Then $M$ satisfies $\mathrm{Nef}_r(\mu,0)$: if $r\in S$ by \cref{pf:towergeo:rootw} (the statement (HT) is part of
$\towergeoG$); if $r\notin S$, then $\mathcal N_{p_D}(r,\mu)$ holds because
$4f\|\mu\|\le4f\le r-2=(r-2)\,\mathrm{gap}_{p_D}(\mu,\tau)$, and \cref{pf:towergeo:nefgood} applies. Take
$\mu=-\mathbf 1_{i\ge p_D}$: it is dominant with gap $0-(-1)=1$, and $\mathrm{Nef}_r(\mu,0)$ says
$\sum_\tau\sum_{i\ge p_D}\deg(x(\tau,i)\cdot W)\ge0$ for all vertical curves $W$ of characteristic $r$. Take
$\mu=\mathbf 1_{i<p_D}$: it is dominant with gap $1-0=1$, and $\mathrm{Nef}_r(\mu,0)$ says
$-\sum_\tau\sum_{i<p_D}\deg(x(\tau,i)\cdot W)\ge0$. By additivity of $\deg(\cdot\cdot W)$ in the multidegree, these are
the assertions.
\end{proof}

% ---------------------------------------------------------------------------
\subsection{Geometric height constants}\label{pf:towergeo:sec-nine}

We use the notation of \cref{pf:towergeo:def-nine}. Throughout, $f=|\Psi_H|$, $N=p+q$ and $j'=N-j$.

\begin{lemma}[Numbers of a host]\label{pf:towergeo:host-numbers}
For a host $H$ on $E$: $\mathrm{rd}_E\ge2$, $\Psi_H\ne\emptyset$ and $2|\Psi_H|=[E:\Q]$.
\formallabel{GlobalTheory.CompanionTheory.Ar.Top.HostNumbers.host_numbers}
\end{lemma}

\begin{proof}
For an embedding $\tau\colon E\to\C$, $\tau(\sqrt{-d})^2=-d$, so $\tau(\sqrt{-d})=\pm\rho_0$, and $\rho_0\ne0$
is purely imaginary, so $\bar\rho_0=-\rho_0\ne\rho_0$. Hence $\tau\mapsto\bar\tau$ maps $\Psi_H$ bijectively
onto its complement, and $2|\Psi_H|=[E:\Q]>0$. Let $F$ be the maximal real subfield of $E$ and $m=[F:\Q]$.
As $E/F$ is unramified (part of the definition of a host), the relative different is trivial and
$|d_E|=|d_F|^{[E:F]}=|d_F|^2$ \cite[Ch.~III, \S2]{Neukirch1999}. If $m=1$ this gives $|d_E|=1$, which is
impossible for a quadratic field. If $m\ge2$, Minkowski's bound for the totally real field $F$ gives
$|d_F|\ge(m^m/m!)^2$ \cite[Ch.~III, \S2]{Neukirch1999}, and $(m^m/m!)^2\ge2^m$: by induction
$2^mm!\le2m^m$ for $m\ge1$ (the step uses $(1+1/m)^m\ge2$), hence $4^m(m!)^2\le4m^{2m}$, and $4\le2^m$
gives $2^m(m!)^2\le m^{2m}$. Therefore $|d_E|\ge4^m=2^{[E:\Q]}$, i.e.\ $\mathrm{rd}_E\ge2$.
\end{proof}

\subsubsection*{The cusp weight on fundamental sets}
Let $E$ be a CM field and $g\in\N$. A \emph{Riemann lattice} $R$ of genus $g$ on an $\OO_E$-module $\Lambda$
which is free of rank $2g$ over $\Z$ consists of a unimodular alternating form $\langle\,,\rangle$ on
$\Lambda$ with $\langle ax,y\rangle=\langle x,\bar ay\rangle$; an additive map $P\colon\Lambda\to\C^g$; a
positive definite Hermitian matrix $H_R$ with $\Im H_R(Px,Py)=\langle x,y\rangle$; and an action of $\OO_E$
on $\C^g$ compatible with $P$. Its \emph{first
minimum} $\lambda_1(R)$ is the infimum of the $H_R$-norms of the vectors $P(v)$, $v\ne0$. A \emph{cusp} of $R$ is
a nonzero saturated isotropic $\OO_E$-submodule $I\subseteq\Lambda$; for a family $\sigma$ of embeddings it
has a normalised covolume $\mathrm{nz}_\sigma(I)>0$, and the \emph{cusp weight} is
$w^{\mathrm{cusp}}_\sigma(R)=\max\bigl(0,\sup_I(-\log\mathrm{nz}_\sigma(I))\bigr)$, the supremum over the
cusps. A \emph{universal family} $F$ of type $(\Psi,j,j')$ consists of a CM family
$\sigma\colon\Psi\to\Hom(\OO_E,\C)$ (so that $2|\Psi|=[E:\Q]$), a fixed form on $\Lambda$, and a holomorphic
 family
$t\mapsto R(t)$ of Riemann lattices with this form over the \emph{ball domain} $\mathcal D\subseteq\C^n$,
$n=|\Psi|jj'$, of signature $(j,j')$ at $\sigma$, with eigenframes whose log-determinant is the ball
potential, such that every Riemann lattice with this form and signature is equivalent to $R(t)$ for exactly
one $t\in\mathcal D$. The \emph{ball density} is $\delta(t)=\exp((j+j')\,\phi_{\mathrm{ball}}(t))$. An
\emph{arithmetic group} $\Gamma$ is a subgroup of finite index of the automorphism group of the form, and a
\emph{fundamental set} is a measurable $\mathcal F\subseteq\mathcal D$ containing exactly one point of every
$\Gamma$-equivalence class.

A \emph{group dictionary} for $F$ consists of a Hermitian form $h$ on $E^{j+j'}$ and an $\OO_E$-linear embedding of
$\Lambda$ whose image is a lattice integral for $h$ (up to a fixed denominator), orthogonality for $h$ corresponding to
the vanishing of the form of $F$ on $\OO_E$-multiples; a positive definite majorant of $h$, with length function
$|\cdot|_M$; the archimedean unitary group $G$ of $h$ and the discrete subgroup $G_\Lambda$ of elements stabilising the
lattice; a continuous map $\pi$ from $G$ onto $\mathcal D$ whose fibres are the cosets of the stabiliser of the
majorant; and $\kappa>0$ such that the squared norm of a lattice vector $x$ in $R(\pi(a))$ is $\kappa|a^{-1}x|_M^2$.
For $C\in\N$, the \emph{short set} $V_C$ is the set of nonzero lattice vectors $x$ with $h(x,x)=0$ whose index in the
saturation of the line $\OO_Ex$ is at most $C$. The following lemma is a general fact on measures.

\begin{lemma}[Mean of the logarithm of a first minimum]\label{pf:towergeo:meanlog}
Let $\mu$ be a measure on a measurable space $X$ and $S\subseteq X$ measurable with $0<\mu(S)<\infty$. Let
$\lambda\colon X\to(0,\infty)$ and $c(t,\cdot)\colon X\to[0,\infty]$ $(t>0)$ be measurable, $a\ge1$ an integer, $k\ge3a$
and $r>0$ real numbers, such that for all $t>0$: $\int_Sc(t,\cdot)\,d\mu<\infty$; $c(t,z)=0$ if $t\le\lambda(z)$;
$1\le c(t,z)\le(3t/\lambda(z))^a$ if $\lambda(z)<t\le r$; and $\int_Sc(\theta t,\cdot)\,d\mu=\theta^k\int_Sc(t,\cdot)\,d\mu$
for $\theta>0$. Then $\int_S\log^+(r/\lambda)\,d\mu\le3\mu(S)$.
\formallabel{GlobalTheory.CompanionTheory.Ar.Cusp.lintegral_logPlus_le_of_mass}
\formallabel{GroundTheory.Heights₂.Heights.LowerBound.MeanLogPlus.full}
\formallabel{GroundTheory.Heights₂.Heights.LowerBound.TailBound.full}
\end{lemma}

\begin{proof}
Let $\nu=\mu(S)^{-1}\mu|_S$, a probability measure, and $m(t)=\int c(t,\cdot)\,d\nu<\infty$; then $m(\theta t)=\theta^km(t)$.

\emph{Tail bound: $\nu(\lambda<e^{-x}r)\le\min(1,2\cdot36^ae^{-kx})$ for $x\ge0$.} For $0<t\le r$ we have
$c(t,\cdot)\ge1$ on $\{\lambda<t\}$, so $\nu(\lambda<t)\le m(t)=(t/r)^km(r)$; it suffices to show $M:=m(r)\le2\cdot36^a$.
Put $G(u)=\nu(\lambda<r/u)$ for $u\ge1$, so $G(u)\le\min(1,u^{-k}M)$. As $c(r,\cdot)\le(3r/\lambda)^a$ on $\{\lambda<r\}$
and vanishes elsewhere, the layer-cake formula gives
\[
M\le3^a\int_{\{\lambda<r\}}(r/\lambda)^a\,d\nu=3^a\Bigl(G(1)+\int_1^\infty au^{a-1}G(u)\,du\Bigr).
\]
If $M\le1$ we are done. Otherwise let $u_0=M^{1/k}>1$; splitting the integral at $u_0$,
\[
\int_1^\infty au^{a-1}G(u)\,du\le(u_0^a-1)+Ma\int_{u_0}^\infty u^{a-1-k}\,du=u_0^a-1+\frac a{k-a}u_0^a,
\]
using $Mu_0^{-k}=1$. Hence $M\le3^au_0^a\frac k{k-a}\le\frac32\,3^aM^{a/k}$, since $k\ge3a$. As $1-a/k\ge2/3$ and $M>1$,
$M^{2/3}\le M^{1-a/k}\le\frac32\,3^a$, so $M\le(3/2)^{3/2}3^{3a/2}\le2\cdot36^a$.

\emph{The mean.} Let $F=\log^+(r/\lambda)$, $C_a=2\cdot36^a$ and $x_0=\log(C_a)/k\ge0$. Then $\{F>x\}=\{\lambda<e^{-x}r\}$ for
$x>0$, and
\begin{multline*}
\int F\,d\nu=\int_0^\infty\nu(F>x)\,dx\le x_0+\int_{x_0}^\infty C_ae^{-kx}\,dx\\
=\frac{\log2+a\log36+1}{k}\le\frac{\log2+1}3+\frac{\log36}3<3,
\end{multline*}
using $k\ge3a\ge3$. Multiplying by $\mu(S)$ gives the assertion.
\end{proof}

\begin{proposition}[First minima on a fundamental domain]\label{pf:towergeo:fmin}
Let $F$ be a universal family of type $(\Psi,j,j')$ with a group dictionary and a Haar measure $\mu$ on $G$, and let
$v>0$. Assume $1\le j'$, $3j'+1\le j+j'$, and that there are $a_1,\dots,a_{[E:\Q]}\in\OO_E$, linearly independent over
$\Q$, with $|\sigma(a_i)|\le v$ for all embeddings $\sigma$. Then there is a fundamental domain $D_\infty$ of
$G_\Lambda$ in $G$ with $0<\mu(D_\infty)<\infty$ and
\[
\int_{D_\infty}\log^+\bigl(v^{-1/2}/\lambda_1(R(\pi(a)))\bigr)\,d\mu(a)\le3\mu(D_\infty).
\]
\formallabel{GlobalTheory.T.first_minimum_integral_bound}
\formallabel{GlobalTheory.CompanionTheory.Ar.Cusp.group_mean_logPlus_le}
\end{proposition}

\begin{proof}
Let $f=|\Psi|$, so $[E:\Q]=2f$, $r=v^{-1/2}$, $\lambda(a)=\lambda_1(R(\pi(a)))$ and
$C=\lfloor(v\sqrt{2f})^{2f}\rfloor$. By \lref{GlobalTheory.T.first_minimum_attained_with_bounds}, which uses
the small elements $a_i$, $1\le j'$ and the signature of $R(\pi(a))$: $\lambda(a)>0$ is attained by a lattice vector
$x_a$; lattice vectors of norm $<r$ are pairwise orthogonal for the form; and for $\lambda(a)<t\le r$ there are at
most $(3t/\lambda(a))^{2fj'}$ lattice vectors of norm $<t$. By \lref{GroundTheory.Heights₂.Heights.FiniteCovolume.ScalingLaw.full}
there is a fundamental domain $D_\infty$ of $G_\Lambda$ such that the counting functions
\[
c(t,a)=\#\{x\in V_C:\sqrt\kappa\,|a^{-1}x|_M<t\},
\]
the numbers of elements of the short set of norm $<t$ in $R(\pi(a))$, have finite integrals over $D_\infty$, homogeneous
of degree $k=2f(j+j'-1)$ in $t$. By \lref{GroundTheory.Heights₂.Heights.FiniteCovolume.ArithmeticCovolumeFinite.full},
$\mu(D_\infty)<\infty$, and $\mu(D_\infty)>0$, since otherwise the countably many translates of $D_\infty$ by
$G_\Lambda$, which cover $G$, would be null, whereas nonempty open sets have positive Haar measure.

We apply \cref{pf:towergeo:meanlog} with $S=D_\infty$, $a=2fj'\ge1$ and $k=2f(j+j'-1)\ge3a$ (as $3j'\le j+j'-1$).
The function $\lambda$ is continuous, hence measurable
(\lref{GlobalTheory.T.first_minimum_continuous_cusp_weight_measurable}). If $t\le\lambda(a)$, no element of the
short set has norm $<t$, so $c(t,a)=0$. If $\lambda(a)<t\le r$, then $c(t,a)\le(3t/\lambda(a))^{2fj'}$ since the short
set consists of lattice vectors, and $c(t,a)\ge1$: the minimal vector $x_a$ has norm $\lambda(a)<r$, so it is orthogonal
to itself, i.e.\ isotropic, and its index is at most $C$ by
\lref{GroundTheory.Heights₂.Heights.FiniteCovolume.ShortestVectorIndex.full} (a shortest nonzero vector of the
lattice at a point of $G$ has index at most $(v\sqrt{2f})^{2f}$ in its saturated line, given the small elements $a_i$), so
$x_a\in V_C$ is counted. The lemma gives the assertion.
\end{proof}

\begin{proposition}[The cusp weight on fundamental sets]\label{pf:towergeo:cusp-fund}
Let $F$ be a universal family of type $(\Psi,j,j')$, $\Gamma$ an arithmetic group and $\mathcal F$ a
fundamental set, and assume $1\le j'$ and $3j'+1\le j+j'$. Then $\delta$ is integrable on $\mathcal F$ with
$\int_{\mathcal F}\delta>0$, the function $w^{\mathrm{cusp}}_\sigma(R(t))\,\delta(t)$ is integrable on
$\mathcal F$, and
\[
\int_{\mathcal F}w^{\mathrm{cusp}}_\sigma(R(t))\,\delta(t)\,dt\le
fj'\bigl(6+\log(2f\,\mathrm{rd}_E)+\log(2fj'\,\mathrm{rd}_E)\bigr)\int_{\mathcal F}\delta(t)\,dt,\qquad f=|\Psi|.
\]
\formallabel{GlobalTheory.T.cusp_weight_integrable_on_fundamental_set}
\formallabel{GlobalTheory.CompanionTheory.Ar.Cusp.fund_set_mean_cuspW}
\end{proposition}

\begin{proof}
Put $v=2f\,\mathrm{rd}_E$ and $r=v^{-1/2}$. Since $|d_E|\ge1$ we have $\mathrm{rd}_E\ge1$, so $v\ge1$ and
$r\le1$. We use the following results.
\begin{enumerate}[label=(\arabic*)]
\item By \lref{GlobalTheory.T.exists_group_dictionary}, $F$ has a \emph{group dictionary}: the
  archimedean unitary group $G$ of the form with a Haar measure $\mu$, a discrete subgroup $G_\Lambda$ of $G$
  defined by the lattice, and a map $\pi$ from $G$ onto $\mathcal D$.
\item By \lref{GlobalTheory.T.cm_field_small_integral_basis}, there are $a_1,\dots,a_{[E:\Q]}\in\OO_E$,
  linearly independent over $\Q$, with $|\sigma(a_i)|\le[E:\Q]\,\mathrm{rd}_E=v$ for all embeddings
  $\sigma$.
\item By \cref{pf:towergeo:fmin}, given (2), $1\le j'$ and $3j'+1\le j+j'$,
  there is a fundamental domain $D_\infty$ of $G_\Lambda$ in $G$ with $0<\mu(D_\infty)<\infty$ and
  $\int_{D_\infty}\log^+\bigl(r/\lambda_1(R(\pi(a)))\bigr)\,d\mu(a)\le3\mu(D_\infty)$.
\item By \lref{GlobalTheory.T.fundamental_set_to_fundamental_domain}, there is $C\in(0,\infty)$ with
  $\int_{\mathcal F}\Phi\,\delta=C\int_{D_\infty}\Phi\circ\pi\,d\mu$ for every measurable
  $\Phi\colon\mathcal D\to[0,\infty]$ which is constant on the equivalence classes for the automorphism group
  of the form.
\item By \lref{GlobalTheory.T.first_minimum_continuous_cusp_weight_measurable}, $\mathcal D$ is open,
  $t\mapsto\lambda_1(R(t))$ is continuous on $\mathcal D$, $t\mapsto w^{\mathrm{cusp}}_\sigma(R(t))$ (extended
  by $0$) is measurable, and both are constant on equivalence classes.
\item By \lref{GlobalTheory.T.first_minimum_and_cusp_weight_bounds}, for $t\in\mathcal D$,
  $\lambda_1(R(t))>0$ and
  $0\le w^{\mathrm{cusp}}_\sigma(R(t))\le fj'\bigl(2\max(0,\log(1/\lambda_1(R(t))))+\log(2fj'\,\mathrm{rd}_E)\bigr)$.
\end{enumerate}
Write $\lambda(t)=\lambda_1(R(t))$ and $w(t)=w^{\mathrm{cusp}}_\sigma(R(t))$. Taking $\Phi=1$ in (4),
$V:=\int_{\mathcal F}\delta=C\mu(D_\infty)\in(0,\infty)$; as $\delta>0$, $\delta$ is integrable on $\mathcal F$
with positive integral. Taking $\Phi=\log^+(r/\lambda)$ on $\mathcal D$, which is measurable and constant on
equivalence classes by (5), (4) and (3) give
\[
\int_{\mathcal F}\log^+(r/\lambda)\,\delta=C\int_{D_\infty}\log^+\bigl(r/\lambda(\pi(a))\bigr)\,d\mu(a)\le3C\mu(D_\infty)=3V.
\]
For $\lambda>0$ we have $\max(0,\log(1/\lambda))\le\log^+(r/\lambda)+\frac12\log v$: if $\lambda\ge1$ the left
side vanishes, and otherwise $\log(1/\lambda)=\log(r/\lambda)+\frac12\log v$. By (6),
\[
0\le w\le fj'\bigl(2\log^+(r/\lambda)+\log v+\log(2fj'\,\mathrm{rd}_E)\bigr)\quad\text{on }\mathcal F.
\]
The right side times $\delta$ is integrable on $\mathcal F$, so $w\delta$ is integrable, and
\[
\int_{\mathcal F}w\,\delta\le fj'\bigl(6+\log v+\log(2fj'\,\mathrm{rd}_E)\bigr)V,
\]
which is the assertion since $\log v=\log(2f\,\mathrm{rd}_E)$.
\end{proof}

\subsubsection*{The three bounds}

\begin{proposition}[Mean of the cusp weight]\label{pf:towergeo:meanU}
Assume $\towergeoG$. Then part (c) of $\textup{(GH)}_{C_h,C_u}$ holds with $C_u=4$, for all data as in
\cref{pf:towergeo:def-nine}.
\formallabel{GlobalTheory.T.cusp_volume_sum}
\formallabel{GlobalTheory.CompanionTheory.Ar.Cusp.top_mean_cusp_weight}
\end{proposition}

\begin{proof}
Let the data be given, and assume $1\le j'$, $3j'+1\le N$ and $B\ge4fN^2$. Note that $j+j'=N$.

\emph{Decomposition into fundamental sets.} By \lref{GlobalTheory.T.cusp_volume_left}, which uses
$\towergeoG$, the origin admits a \emph{decomposition into pieces}: finitely many universal families
$F_k$ with arithmetic groups $\Gamma_k$, fundamental sets $\mathcal F_k$ and ball densities $\delta_k$,
multiplicities $d_k\in\N$ not all zero, and maps $\gamma_m$ from parameter pieces $V_m$ to the complex
points of the variety. By \lref{GlobalTheory.T.ptsIntegral_eq_sum_ballDens_integrals} there is
$\kappa>0$ such that for every function $G$ on the complex points which is invariant under the torus action
and is given on the pieces by functions $G_k$ (i.e.\ $G(\gamma_m(t))=G_k(t)$ for $t\in V_m$, $k$ the family
of the piece $m$), with $G_k\delta_k$ integrable on $\mathcal F_k$,
\begin{equation}\label{pf:towergeo:eq-pieces}
\int G\,d\mu=\kappa\sum_kd_k\int_{\mathcal F_k}G_k(t)\,\delta_k(t)\,dt,
\end{equation}
where $\mu$ is the measure on the complex points defined by $\phi_\mu$. By
\lref{GlobalTheory.T.cusp_volume_middle}, the cusp weight $U$ is torus invariant and is given on the
pieces by the cusp weights $w^{\mathrm{cusp}}(R_k(t))$ of the lattices of the families $F_k$.

\emph{The mean.} By \cref{pf:towergeo:cusp-fund}, applied to each $F_k$ (with $1\le j'$ and
$3j'+1\le j+j'$), $\delta_k$ and $w^{\mathrm{cusp}}(R_k)\delta_k$ are integrable on $\mathcal F_k$,
$\int_{\mathcal F_k}\delta_k>0$, and
$\int_{\mathcal F_k}w^{\mathrm{cusp}}(R_k)\delta_k\le C\int_{\mathcal F_k}\delta_k$ with
$C=fj'(6+\log(2f\,\mathrm{rd}_E)+\log(2fj'\,\mathrm{rd}_E))$. The degree $\deg_\mu$ is by definition the
integral of the constant function $1$, so \eqref{pf:towergeo:eq-pieces} with $G=1$ gives
$\deg_\mu=\kappa\sum_kd_k\int_{\mathcal F_k}\delta_k>0$, because $\kappa>0$, all terms are nonnegative and one
$d_k$ is positive. With $G=U$,
\[
\int U\,d\mu=\kappa\sum_kd_k\int_{\mathcal F_k}w^{\mathrm{cusp}}(R_k)\delta_k\le
C\kappa\sum_kd_k\int_{\mathcal F_k}\delta_k=C\deg_\mu,
\]
so $\mathrm{mean}(U)\le C$.

\emph{Numerics.} By \cref{pf:towergeo:host-numbers}, $\mathrm{rd}_E\ge2$ and $f\ge1$. From $j'\ge1$ and
$3j'+1\le N$ we get $N\ge4$, hence $B\ge4fN^2\ge64$, and $2fj'\le4fN^2\le B$ since $j'\le N^2$. Then
$\log(2f\,\mathrm{rd}_E)\le\log(B\,\mathrm{rd}_E)$ and $\log(2fj'\,\mathrm{rd}_E)\le\log(B\,\mathrm{rd}_E)$;
moreover $B\,\mathrm{rd}_E\ge128$, so $\log(B\,\mathrm{rd}_E)\ge7\log2>3$. Hence
$6+\log(2f\,\mathrm{rd}_E)+\log(2fj'\,\mathrm{rd}_E)\le4\log(B\,\mathrm{rd}_E)\le4\ell_H(B)$, and
$\mathrm{mean}(U)\le C\le4fj'\ell_H(B)$.
\end{proof}

For an admissible origin and $B\in\N$ let $S_B=\{r\in S:r\le B\}$. The next two results give part (a).

\begin{proposition}[Nefness of the twisted Hodge degree]\label{pf:towergeo:idu}
Assume $\towergeoG$ and \textup{(W)}. For a component origin which is admissible for $B_1$ and every $B\ge4f+2$, the
multidegree $a_\mu+\sum_{r\in S_B}v_r$ is nef on the fibres of the variety.
\formallabel{GlobalTheory.T.intersection_degree_upper}
\formallabel{GlobalTheory.CompanionTheory.Relay.intersection_degree_upper}
\end{proposition}

\begin{proof}
By \lref{GlobalTheory.T.compact_datum_intersection}, the Hodge component $\towergeoHC$ of the origin, with the
primes $S_0\cup S_B$ and the lattices patched at these primes, is again a Hodge component (the elements of
$S_0\cup S_B$ are primes $p$ with $R(p)\ge4$ and full $p^{R(p)}$-torsion, and the patched lattices are vector bundles
of the ranks of the labels). At the generic point choose a chart; by the local description of the Hodge data there
are frames $b_\tau$ of $W_\tau$ and $\lambda\ne0$ such that the Hodge section is $\lambda$ times the tensor product
of the wedges of the last $j'$ vectors of $b_\tau$ modulo $F_\tau$. Let $s_\tau(k)\in V^{\mathrm{gen}}_{s(\tau)}$ be
the component of the $k$-th of these vectors in the decomposition
$W_\tau\cong V^{\mathrm{gen}}_{b(\tau)}\oplus V^{\mathrm{gen}}_{s(\tau)}$. By
\lref{GlobalTheory.T.intersection_datum_degree_sum}, the $s_\tau$ are linearly independent and the tensor
product of the determinant lattices of the patched lattices $V_{s(\tau)}$ in the frames $s_\tau$ has multidegree
$a_\mu+\sum_{r\in S_B}v_r$; by counting dimensions the $s_\tau$ span $V^{\mathrm{gen}}_{s(\tau)}$. The choice
$s(\tau)$ is the same label side for all $\tau$, and $p_D,q_D\ge1$ by the side condition of the origin.

By \lref{GlobalTheory.T.hodge_comp_cover_cm_signs}, which uses $\towergeoG$ and (W), for a Hodge component of a
global model which covers the variety and is Hodge--Tate free at all its primes, such that every other prime $r$
satisfies $r\nmid d_E$, $r\ge4f+2$ and $r>2$, and for spanning families of the eigenspaces of one side, the multidegree
of the tensor product of the determinant lattices is nef on the fibres (its proof passes to a root model covering the
component and applies \cref{pf:towergeo:blocknef}). The hypotheses hold: $Y$ is global since the origin is
admissible, and $\towergeoHC.X$ covers the variety (part of the component origin); $\towergeoHC$ is Hodge--Tate free
at the primes of $S_0\cup S_B$ by goodness; and if
$r\notin S_0\cup S_B$ is a prime, then $r\nmid2d_EM'$ (otherwise $r\in S_0$), so $r\nmid d_E$ and $r\ne2$, and
$r\ge4f+2$, since otherwise $r<4f+2\le\min(B,B_1)$ would give $r\in S$ and $r\in S_B$.
\end{proof}

\begin{proposition}[Positivity from nefness]\label{pf:towergeo:posu}
Assume the metric package of $\towergeoG$. Let $\mathrm{II}$ be a full intersection system, component data with
Hodge data, a complex model and a tied bounded regularisation setting for a weight $w$. If $a_\mu+\sum_{r\in S_B}v_r$
is nef on the fibres for every $B\ge4f+2$, then for all $B\ge4f+2$ and $1\le\alpha\le\beta$ there is $t_0>0$ such that
for $0<t\le t_0$, $\mathcal M_t(B,\alpha,\beta)$ is semipositive, nef on the fibres and ample on the generic fibre.
\formallabel{GlobalTheory.T.posU_of_full_intersection_system}
\formallabel{GlobalTheory.CompanionTheory.Ar.Cusp.posU_of_nefOnFibres}
\end{proposition}

\begin{proof}
For a function $\phi$ on a chart let $\mathcal L\phi(t;v)$ be its Levi form at $t$ in the direction $v$. By
\lref{GlobalTheory.T.uniform_levi_positivity_on_charts} there is $\delta_0>0$ with
$\mathcal L(\log h_{A'}\circ\gamma)(t;v)\ge\delta_0|v|^2$ for all charts $\gamma$ and $|t|\le\frac12$. Put
$\delta=\alpha\delta_0/(2\beta)>0$. Since the setting is tied, for all sufficiently small $t>0$ the metric of
$\mathcal M_t$ is given by the formula of \cref{pf:towergeo:def-nine}
(\lref{GlobalTheory.T.regularized_weight_dominates_eventually}) and
$\mathcal L((\log h_{\mathrm{ref}}+\towergeoreg_t)\circ\gamma)(t';v)\ge-\delta|v|^2$ on the charts
(\lref{GlobalTheory.T.regularized_levi_lower_bound_eventually}); choose $t_0>0$ such that both hold for
$0<t\le t_0$. Fix such $t$, a chart and $|t'|\le\frac12$. By the formula for the metric,
\[
\mathcal L(\log h_{\mathcal M_t}\circ\gamma)=\beta\,\mathcal L((\log h_{\mathrm{ref}}+\towergeoreg_t)\circ\gamma)
+\alpha\,\mathcal L(\log h_{A'}\circ\gamma)
\]
(\lref{GlobalTheory.CompanionTheory.Ar.Cusp.chart_levi_split}). With $V=|v|^2$, $L_F\ge-\delta V$ and
$L_A\ge\delta_0V$ for the two Levi forms on the right,
\[
\beta L_F+\alpha L_A\ge-\tfrac12\alpha\delta_0V+\alpha\delta_0V\ge0,\qquad
2(\beta L_F+\alpha L_A)-\alpha L_A\ge-\alpha\delta_0V+\alpha\delta_0V=0.
\]
The first inequality shows that $\mathcal M_t$ is semipositive, and the second that $\mathcal M_t^{\otimes2}\otimes
A'^{\otimes-\alpha}$ is semipositive (\lref{GlobalTheory.T.semipositive_levi_of_log_metric}: a smooth Hermitian
line whose metric has nonnegative Levi form on all charts is semipositive). The multidegree of $\mathcal M_t$ is
$\beta(a_\mu+\sum_{r\in S_B}v_r)+\alpha a'$, nef on the fibres as a sum of a nef multidegree and a nonnegative one
(degrees on curves are additive in the multidegree, and nonnegative multidegrees have nonnegative degrees). Finally $\mathcal M_t^{\otimes2}=(\mathcal M_t^{\otimes2}\otimes
A'^{\otimes-\alpha})\otimes A'^{\otimes\alpha}$ with the first factor semipositive and $A'$ positive of multidegree
$a'>0$, so $\mathcal M_t$ is ample on the generic fibre
(\lref{GlobalTheory.T.ample_on_generic_fibre_of_semipositive}, which uses the full intersection system and the
metric package).
\end{proof}

\begin{proposition}[Positivity]\label{pf:towergeo:posU}
Assume $\towergeoG$ and \textup{(W)}. Then part (a) of $\textup{(GH)}_{C_h,C_u}$ holds for all data as in
\cref{pf:towergeo:def-nine}.
\formallabel{GlobalTheory.T.intersection_degree_lower}
\formallabel{GlobalTheory.CompanionTheory.Relay.intersection_degree_lower}
\end{proposition}

\begin{proof}
By \cref{pf:towergeo:idu}, for an admissible origin and every $B\ge4f+2$ the multidegree
$a_\mu+\sum_{r\in S_B}v_r$ is nef on the fibres of the variety. By \cref{pf:towergeo:posu}, for a full intersection
system and a tied setting this nefness implies part (a).
\end{proof}

\subsubsection*{Heights of integral points}
Let $L$ be a number field and $z$ an $\OO_L$-point of the Hodge component $\towergeoHC$ of a component origin.
\emph{Integral linear point data} at $z$ consist of a model $V$ over $L$ of the Lie algebra of the abelian part of the
semi-abelian scheme at $z$, an action $\rho$ of $\OO_E$ on $V$, a bijection $s\mapsto\sigma_s$ from the labels to the
embeddings $\OO_E\to L$ with eigenspaces $V_s$, a finitely generated $\rho$-stable $\OO_L$-lattice $\Lambda$ spanning
$V$, a $\rho$-stable subgroup $I\subseteq\Lambda$ spanned by an $L$-linearly independent family, and a Hermitian
$\OO_L$-bundle $B_q$ with a map $\Lambda\to B_q$ with kernel $\OO_LI$, surjective up to scalars, under which distinct
eigenspaces become orthogonal; they are compatible with $z$ and the Hodge component. An \emph{$\OO$-family} is a
family $v_1,\dots,v_r\in I$, linearly independent over $\OO_E$, whose $\OO_E$-span has finite index $[I:\OO_Ev]$ in
$I$. A family $q_1,\dots,q_k\in\Lambda\cap V_s$ is \emph{adapted} to $v$ at $s$ if the $\sigma_s$-eigencomponents of
the $v_a$ followed by the $q_c$ form a basis $\mathrm{fam}_s(v,q)$ of $V_s$. For $\OO_L$-lattices
$F_\tau\subseteq V_{s(\tau)}$ and adapted families $q_\tau$ at $s(\tau)$, the \emph{$\phi$-degree} is
\begin{multline*}
\deg_\phi=\sum_\tau\frac1{[L:\Q]}\Bigl(\log[F_\tau:\mathrm{fam}_{s(\tau)}(v,q_\tau)]
-\frac12\sum_{\iota\colon L\to\C}\log\mathrm{Gram}_\iota(q_\tau)\Bigr)\\
-\frac12\log[I:\OO_Ev]+\frac{fr}2\log(4\pi),
\end{multline*}
where $[F:b]$ is the generalised index of the $\OO_L$-span of $b$ in $F$ and $\mathrm{Gram}_\iota(q)$ is the Gram
determinant of the images of the $q_c$ in $B_q$ for the metric at $\iota$. Below $F_\tau$ is the fibre at $z$ of the
patched lattice $V_{s(\tau)}$ of $\towergeoHC$ with the primes $S_0\cup S_B$, transported to $V$.

\begin{proposition}[Height comparison]\label{pf:towergeo:phideg}
Assume $\towergeoG$ and \textup{(FD)}. There is $C\ge0$ such that for all data as in \cref{pf:towergeo:def-nine} with
a component origin admissible for $B_1$, every $B\ge4fN^2$, every integral linear point at an $\OO_L$-point $z$ whose
Hermitian bundle has rank at most $fN$, every $\OO$-family $v$ of length $r$ and adapted families $q_\tau$ of length
$k$ at the labels $s(\tau)$ with $r+k=j'$, one has $\deg_\phi\le C\,fj'\,\ell_H(B)$.
\formallabel{GlobalTheory.T.phiDeg_fibLat_le_of_isOFamily_of_isAdapted}
\formallabel{LevelTowerAlt.height_comparison_global_model}
\end{proposition}

\begin{proof}
\emph{Constants.} By \lref{GlobalTheory.CompanionTheory.Ar.Cusp.U4.flat_saturation_of_degree}, $\towergeoG$ and
(FD) give the \emph{flat saturation bound}: there is $C_f\ge0$ such that for every principal semi-abelian scheme with
$\OO_E$-action of signature type $(p',q')$ over $\OO_L$, the Hermitian bundle $B$ of its Lie algebra with the Hodge
metric, and conjugate embeddings $s_1,s_2=\bar s_1$, the eigen-subbundles satisfy
$|\deg B_{s_1}-\deg B_{s_2}|\le C_f(\mathrm{rk}B_{s_1}+\mathrm{rk}B_{s_2})(1+\log\mathrm{rd}_E+|L_H|)$, degrees being
normalised by $[L:\Q]$. By \lref{GlobalTheory.T.sum_log_prime_div_pred_le_log} there is $C_1\ge0$ with
$\sum_{p\in S'}\frac{\log p}{p-1}\le C_1\log B$ for $B\ge2$ and every set $S'$ of primes each of which is $\le B$ or in a
set $S''$ with $|S''|\le B$. Let $C_s$ be the constant of the slope bound of $\towergeoG$ (the maximal slope of the
Hermitian bundle of such a semi-abelian scheme is at most $C_s\log(2+\mathrm{rk}B+\max(-\deg B,0))$), and
$C_0=C_s+C_f(2+C_1)$. By \lref{LevelTowerAlt.abel_abstract} there is $C_a\ge0$ such that for every Hermitian
$\OO_L$-bundle $B$, sub-bundles $F_\tau,F'_\tau$ ($\tau\in\Psi$), $m,N\in\N$ and $\ell\in\R$ with $fm\ge1$, $\ell\ge1$,
$fN\le e^\ell$, $\mathrm{rk}B\le fN$, $\mathrm{rk}F_\tau\le m$, maximal slope of $B$ at most
$C_0\log(2+\mathrm{rk}B+\max(-\deg B,0))$, $\sum_\tau\deg F_\tau+\sum_\tau\deg F'_\tau\le\deg B$ and
$|\deg F_\tau-\deg F'_\tau|\le C_0N\ell$, one has $\sum_\tau\deg F_\tau\le C_afm\ell$. Put $C=2C_1+3+C_a$.

\emph{Numerics.} Let $\ell=\ell_H(B)=\log B+\log\mathrm{rd}_E+|L_H|$. By \cref{pf:towergeo:host-numbers},
$\mathrm{rd}_E\ge2$ and $f\ge1$; $N\ge2$ and $j'\ge1$. Hence $B\ge4fN^2\ge16$, $fN\le B$, $|S_0|\le4fN^2\le B$,
$1\le\log B\le\ell$, $\log(4\pi)\le\ell$ and $fN\le B\le e^\ell$.

\emph{The Hermitian bundle.} Let $B=B_q$ and, for a label $s$, $B_s$ the saturation of the image of $V_s$. The ranks
of the $B_s$ are those of the images $\mathrm{img}(s)$ of $\Lambda\cap V_s$, $\deg\mathrm{img}(s)\le\deg B_s$, and
$\sum_s\deg B_s\le\deg B$ (\lref{LevelTowerAlt.eig_facts}); $\mathrm{rk}B_{s(\tau)}=\dim V_{s(\tau)}=r+k=j'$, and
$\mathrm{rk}B_s$ is at most the rank of the label $s$ (\lref{GlobalTheory.T.adapted_family_exists_on_eigenspace}).
The labels at $z$ correspond to the embeddings of $\OO_E$, conjugate labels to conjugate embeddings
(\lref{GlobalTheory.T.cusp_correspondence_bound}), and the abelian part has the signature of $D$
(\lref{GlobalTheory.T.abelian_part_lie_signature}). By the slope bound, the maximal slope of $B$ is at most
$C_0\log(2+\mathrm{rk}B+\max(-\deg B,0))$, and by the flat saturation bound
\[
|\deg B_{(\tau,1)}-\deg B_{(\tau,2)}|\le C_fN(1+\log\mathrm{rd}_E+|L_H|)\le C_fN(2+C_1)\ell\le C_0N\ell,
\]
since $1\le\ell$ and $\log\mathrm{rd}_E+|L_H|\le\ell$. The choice of $C_a$ with $m=j'$, $F_\tau=B_{s(\tau)}$ and
$F'_\tau=B_{b(\tau)}$ gives $T:=\sum_\tau\deg\mathrm{img}(s(\tau))\le\sum_\tau\deg B_{s(\tau)}\le C_afj'\ell$.

\emph{Fibre lattices.} The patched Hodge component with primes $S_0\cup S_B$ is Hodge--Tate free at these primes,
which include all prime divisors of $d_E$. By \lref{GlobalTheory.T.abelian_part_integral_point}, which uses
$\towergeoG$, for every $\tau$ there is a nonzero ideal $\mathfrak a$ of $\OO_L$ with
$\mathfrak aF_\tau\subseteq\Lambda\cap V_{s(\tau)}$, $\mathfrak a(\Lambda\cap V_{s(\tau)})\subseteq F_\tau$ and
$\log\mathrm N\mathfrak a/[L:\Q]\le X:=2\sum_{p\in S_0\cup S_B}\frac{\log p}{p-1}+2\log\mathrm{rd}_E$. Hence
(\lref{LevelTowerAlt.phiDeg_le_parts}, which compares the indices and uses the orthogonality of the eigenspaces)
\[
\deg_\phi\le f(r+k)X-fr\log\mathrm{rd}_E+T+\frac{fr}2\log(4\pi).
\]
By the choice of $C_1$, $X\le2C_1\log B+2\log\mathrm{rd}_E\le(2C_1+2)\ell$. With $r+k=j'$, $r\le j'$,
$\log\mathrm{rd}_E\ge0$ and $\log(4\pi)\le\ell$ we obtain $\deg_\phi\le(2C_1+2+C_a+1)fj'\ell=Cfj'\ell$.
\end{proof}

A number $x$ is a \emph{weight of the origin} at a complex point $z'$ with coordinates $w'$ if there are a chart $c$
and a unit $u$ in which the Hodge section is $u$ times the standard section near $z'$, frames of the patched lattices
$V_{s(\tau)}$ at $z'$ (the small part), and the corresponding volume $\nu>0$ of the small part for the polarisation,
such that $x=\log\bigl(|u(w')|\prod_k|w'_{k,c_k}|^{a_\mu(k)}\bigr)+\frac12\log\nu$. For $w$ with coordinates in $L$
let
\[
h^{\mathrm{pt}}(w)=\frac1{[L:\Q]}\Bigl(\sum_{\mathfrak p}\sum_km_k\log\max_i|w_{k,i}|_{\mathfrak p}
+\sum_{\sigma\colon L\to\C}\sum_{p\in S_B}\log\phi^{\mathrm{vert}}_p(\sigma w)\Bigr),\qquad m=a_\mu+\sum_{p\in S_B}v_p,
\]
the sum over the finite places $\mathfrak p$ of $L$ with their normalised absolute values.

\begin{proposition}[Heights of integral linear points]\label{pf:towergeo:cih}
Assume $\towergeoG$. Let a good component origin with primes $S_0$ and $B\in\N$ be given, $z$ an $\OO_L$-point of its
Hodge component whose generic point has coordinates $w$ in a chart, integral linear point data at $z$, an
$\OO$-family $v$ of length $r$ and adapted families $q_\tau$ of length $k$ at the labels $s(\tau)$. Then there is
$x\colon\Hom(L,\C)\to\R$ such that $x(e)$ is a weight of the origin at $e(z)$ for every $e$, and
\[
h^{\mathrm{pt}}(w)+\frac1{[L:\Q]}\sum_ex(e)=-\deg_\phi .
\]
\formallabel{GlobalTheory.T.cusp_intersection_height}
\formallabel{GlobalTheory.CompanionTheory.Ar.Cusp.PtHt.g_height}
\end{proposition}

\begin{proof}
\emph{Local data.} By \lref{GlobalTheory.CompanionTheory.Ar.Cusp.PtHt.datum} there are a chart $c$ with
$w_{k,c_k}\ne0$, a unit $u$ in which the Hodge section is $u$ times the standard section near $z$, with value
$u_L\ne0$ at $z$, and frames $s_\tau$ of the patched lattices $V_{s(\tau)}$ at the generic point of $z$ which form a
small part there. No prime of $S_0\cup S_B$ vanishes in $L$, so at the generic point of $z$ the patched lattices are
the de Rham lattices. By \lref{GlobalTheory.T.abelian_part_integral_point_alt} the frames $s_\tau$ give bases
$s^L_\tau$ of $V_{s(\tau)}$; comparing dimensions, $r+k=N-j=j'$. Let $C_\tau$ be the matrix expressing
$\mathrm{fam}_{s(\tau)}(v,q_\tau)$ in the basis $s^L_\tau$; it is invertible since both are bases.

\emph{Finite places.} By \lref{GlobalTheory.T.abelian_part_lattice_eigen_compat}, for every finite place
$\mathfrak p$ there is a matrix $T_\tau(\mathfrak p)$ with $\det T_\tau(\mathfrak p)\ne0$, integral with unit
determinant for almost all $\mathfrak p$, such that $F_\tau$ consists of the elements of $V_{s(\tau)}$ which lie,
locally at every $\mathfrak p$, in the span of $s^L_\tau T_\tau(\mathfrak p)$; $T_\tau(\mathfrak p)$ is the value at $z$
of the matrix expressing a frame of $V_{s(\tau)}$ at the point of $z$ over $\mathfrak p$ in the frame $s_\tau$. By
\lref{GlobalTheory.T.glued_lattice_from_local_matrices},
\begin{equation}\label{pf:towergeo:eq-idx}
\log[F_\tau:\mathrm{fam}_{s(\tau)}(v,q_\tau)]=\sum_{\mathfrak p}\log|\det T_\tau(\mathfrak p)|_{\mathfrak p}
+\sum_\sigma\log|\sigma(\det C_\tau)|.
\end{equation}
For $p\in S_B$, \lref{GlobalTheory.T.exists_val_at_vert_factor_eq_norm_mul_prod} gives $t_p\in L^\times$, the value
at $z$ of the vertical function of $p$ in the chart $c$, with
$\phi^{\mathrm{vert}}_p(\sigma w)=|\sigma t_p|\prod_k|\sigma w_{k,c_k}|^{v_p(k)}$ for all $\sigma$. Put
$g=u_L\prod_{p\in S_B}t_p\ne0$. By \lref{GlobalTheory.T.intersection_datum_degree_sum}, the inverse of the tensor
product of the determinant lattices of the $V_{s(\tau)}$ in the frames $s_\tau$ is the line bundle of multidegree $m$,
with a generator whose value at $z$ is $g$ up to the chart monomial; hence by
\lref{GlobalTheory.T.tensor_lattice_coord_degree}, at every finite place,
$\prod_\tau|\det T_\tau(\mathfrak p)|_{\mathfrak p}\cdot\rho_{\mathfrak p}=|g|_{\mathfrak p}$, where
$\rho_{\mathfrak p}=\prod_k\bigl(\max_i|w_{k,i}|_{\mathfrak p}/|w_{k,c_k}|_{\mathfrak p}\bigr)^{m_k}$. Summing over
$\mathfrak p$ and using the product formula (\lref{GlobalTheory.CompanionTheory.Ar.RootHt.RatioSum.ratio_sum}),
\begin{equation}\label{pf:towergeo:eq-fin}
\sum_\tau\sum_{\mathfrak p}\log|\det T_\tau(\mathfrak p)|_{\mathfrak p}
=-\sum_{\mathfrak p}\sum_km_k\log\max_i|w_{k,i}|_{\mathfrak p}-\sum_\sigma\Bigl(\log|\sigma g|+\sum_km_k\log|\sigma w_{k,c_k}|\Bigr).
\end{equation}

\emph{Archimedean places.} By \lref{GlobalTheory.T.abelian_part_integral_lattice}, for every $e\colon L\to\C$ we
have $\mathrm{Gram}_e(q_\tau)>0$, and
$\nu(e)=[I:\OO_Ev]\prod_\tau(4\pi)^{-r}\mathrm{Gram}_e(q_\tau)\,|e(\det C_\tau)|^{-2}$ is the volume of the small part at
$e(z)$ in the frames $s_\tau$. Hence $x(e)=\log\bigl(|e(u_L)|\prod_k|e(w_{k,c_k})|^{a_\mu(k)}\bigr)+\frac12\log\nu(e)$ is a
weight of the origin at $e(z)$.

\emph{The identity.} Since $m=a_\mu+\sum_{p\in S_B}v_p$ and $\log|eg|=\log|e(u_L)|+\sum_p\log|e(t_p)|$, the formula for
$\phi^{\mathrm{vert}}_p$ gives
\begin{multline*}
[L:\Q]\,h^{\mathrm{pt}}(w)+\sum_ex(e)=\sum_{\mathfrak p}\sum_km_k\log\max_i|w_{k,i}|_{\mathfrak p}\\
+\sum_e\Bigl(\log|eg|+\sum_km_k\log|e(w_{k,c_k})|\Bigr)+\frac12\sum_e\log\nu(e).
\end{multline*}
By \eqref{pf:towergeo:eq-fin} the first two terms equal
$-\sum_\tau\sum_{\mathfrak p}\log|\det T_\tau(\mathfrak p)|_{\mathfrak p}$, and
\begin{multline*}
\frac12\sum_e\log\nu(e)=[L:\Q]\Bigl(\frac12\log[I:\OO_Ev]-\frac{fr}2\log(4\pi)\Bigr)\\
+\frac12\sum_\tau\sum_e\log\mathrm{Gram}_e(q_\tau)-\sum_\tau\sum_e\log|e(\det C_\tau)|.
\end{multline*}
By \eqref{pf:towergeo:eq-idx} the total is
\begin{multline*}
-\sum_\tau\Bigl(\log[F_\tau:\mathrm{fam}_{s(\tau)}(v,q_\tau)]-\frac12\sum_e\log\mathrm{Gram}_e(q_\tau)\Bigr)\\
+[L:\Q]\Bigl(\frac12\log[I:\OO_Ev]-\frac{fr}2\log(4\pi)\Bigr)=-[L:\Q]\deg_\phi.\qedhere
\end{multline*}
\end{proof}

\begin{proposition}[Heights of points]\label{pf:towergeo:htU}
Assume $\towergeoG$ and \textup{(FD)}. There is $C_h\ge0$ such that part (b) of $\textup{(GH)}_{C_h,C_u}$
holds with this $C_h$ for all data as in \cref{pf:towergeo:def-nine}.
\formallabel{GlobalTheory.T.cusp_height_upper}
\formallabel{GlobalTheory.CompanionTheory.Ar.Cusp.TopHt.htU_of_given}
\end{proposition}

\begin{proof}
Let $C\ge0$ be the constant of \cref{pf:towergeo:phideg}, which uses $\towergeoG$ and (FD): for admissible origins
and $B\ge4fN^2$, the $\phi$-degree of the fibre lattices at an integral linear point, with respect to an
$\OO$-family and adapted families of the eigenspaces, is at most $Cfj'\ell_H(B)$. Put $C_h=C$. Let the data be
given, $B\ge4fN^2$ and $1\le\alpha\le\beta$.

\emph{Choice of $t_0$.} By \lref{GlobalTheory.T.regularized_weight_dominates_eventually}, since the
setting is tied, there is $u>0$ such that for $0<t<u$: $\phi_{\mathrm{sing}}\le\towergeoreg_t$, and
$\mathcal M_t=\mathcal M_t(B,\alpha,\beta)$ is smooth, of multidegree $\beta(a_\mu+\sum_rv_r)+\alpha a'$, with
$\log h(z)=\beta\bigl(\log h_{\mathrm{ref}}(z)+\towergeoreg_t(z)+\sum_r\log\phi^{\mathrm{vert}}_r(z)\bigr)+\alpha\log h_{A'}(z)$
at all complex points $z$ of the variety. Put $t_0=u/2$ and let $0<t\le t_0$.

\emph{Reduction to a point.} Let $P$ be an arithmetic curve on the variety. By
\lref{GlobalTheory.T.height_of_zero_dim_variety}, $P$ is the coordinate ideal of a point $w$ with
coordinates in a number field $L_0$, and $\mathrm{II}(P;\mathcal L)/\deg P=h_{\mathcal L}(w)$ for every smooth
Hermitian line $\mathcal L$, where $h_{\mathcal L}$ is the height of points. Hence
$h_{\mathcal M_t}(P)=h_{\mathcal M_t}(w)/\beta$, and $h_{A'}(w)\ge-c_{A'}$ by the height condition in the component
data. By the definition of heights as sums of local contributions
(\lref{GlobalTheory.CompanionTheory.Ar.Cusp.TopHt.lineHt_decomp}: the height for a metric of the form
above splits accordingly), for any number field $L$ over which $w$ is defined,
\begin{equation}\label{pf:towergeo:eq-htdec}
h_{\mathcal M_t}(w)=\beta\Bigl(h^{\mathrm{pt}}(w)+\frac1{[L:\Q]}\sum_{\sigma\colon L\to\C}
\bigl(\log h_{\mathrm{ref}}+\towergeoreg_t\bigr)(\sigma w)\Bigr)+\alpha\,h_{A'}(w),
\end{equation}
where $h^{\mathrm{pt}}$ is the point height defined before \cref{pf:towergeo:cih} (with $\sum_r$ over $r\in S_B$);
heights do not change under extension of the number field
(\lref{GlobalTheory.T.line_height_invariant_scaling}).

\emph{Integral linear points.} By \lref{GlobalTheory.T.small_points_over_number_field}, there are a
finite extension $L$ of $L_0$, an $\OO_L$-point $z$ of the Hodge component of the origin with coordinates
$w$, the abelian part of the semi-abelian scheme at $z$, a model of its Lie algebra, and integral linear point
data whose eigenspaces have the ranks of the labels and whose lattice has rank at most $fN$. By
\lref{GlobalTheory.T.cusp_correspondence_bound}, the labels at $z$ correspond bijectively to the
embeddings of $\OO_E$, and $2|\Psi_H|=[E:\Q]$, so $\Psi_H\ne\emptyset$. By
\lref{GlobalTheory.T.adapted_family_exists_on_eigenspace}, there are an $\OO$-family $v$ of some length
$r_0$ and, for every $\tau$, an adapted family $\theta_\tau$ of length $j'-r_0$ in the eigenspace of the label
$s(\tau)$, whose dimension is $j'$.

\emph{The bound.} By \cref{pf:towergeo:cih}, there is
$x\colon\Hom(L,\C)\to\R$ such that $x(\sigma)$ is a weight of the origin at the complex point $\sigma(z)$ for
every $\sigma$, and $h^{\mathrm{pt}}(w)+[L:\Q]^{-1}\sum_\sigma x(\sigma)=-\deg_\phi$, the $\phi$-degree of the fibre
lattices with respect to $v$ and $(\theta_\tau)$. By \cref{pf:towergeo:phideg}, $\deg_\phi\le Cfj'\ell_H(B)$. By
\lref{GlobalTheory.T.cusp_intersection_finite}, since $U$ is the cusp weight,
\[
x(\sigma)\le\phi_{\mathrm{sing}}(\sigma w)+\log h_{\mathrm{ref}}(\sigma w)
\le\log h_{\mathrm{ref}}(\sigma w)+\towergeoreg_t(\sigma w).
\]
Inserting these bounds in \eqref{pf:towergeo:eq-htdec},
\begin{align*}
h_{\mathcal M_t}(P)&=h^{\mathrm{pt}}(w)+\frac1{[L:\Q]}\sum_\sigma(\log h_{\mathrm{ref}}+\towergeoreg_t)(\sigma w)
+\frac\alpha\beta\,h_{A'}(w)\\
&\ge-\deg_\phi-\frac\alpha\beta c_{A'}\ge-Cfj'\ell_H(B)-\frac\alpha\beta c_{A'}.\qedhere
\end{align*}
\end{proof}

\begin{theorem}\label{pf:towergeo:nine}
$\towergeoG$, \textup{(W)} and \textup{(FD)} imply that $\textup{(GH)}_{C_h,C_u}$ holds for some
$C_h,C_u\in\R$.
\formallabel{GlobalTheory.T.exists_geometric_height_constants}
\formallabel{GlobalTheory.CompanionTheory.Ar.Cusp.TopNine.nineGeoU_of_given}
\end{theorem}

\begin{proof}
Let $C_h$ be the constant of \cref{pf:towergeo:htU} and $C_u=4$ (\cref{pf:towergeo:meanU}); both are chosen
before the data. For data as in \cref{pf:towergeo:def-nine}, part (a) holds by \cref{pf:towergeo:posU},
part (b) with $C_h$ by \cref{pf:towergeo:htU}, and part (c) with $C_u$ by \cref{pf:towergeo:meanU}.
\end{proof}

\begin{remark}\label{pf:towergeo:rem-nine}
The tied regularisation setting is not used for part (c); (W) is used only for part (a), and (FD) only for
part (b). The value $C_u=4$ is explicit, while $C_h$ comes from the bound for $\deg_\phi$.
\end{remark}

% ---------------------------------------------------------------------------
\subsection{Two auxiliary statements}\label{pf:towergeo:sec-aux}

The following two results of the formal development are used in the proof of the volume theorem.
For a field $F$, an \emph{intersection family} over $F$ assigns to every multihomogeneous prime
$V\subset F[x]$ defining a projective variety of dimension $n$ and multidegrees $a_1,\dots,a_n$ a number
$I(V;a_1,\dots,a_n)$; it is a \emph{field intersection family} if it satisfies the axioms of intersection
numbers of line bundles on multiprojective varieties of the formal development (normalisation and the
inductive formula through divisors of sections). For $F=\Q$ we write $\mathrm{int}_\Q$ for the intersection
numbers of the generic fibres of arithmetic varieties.

\begin{lemma}[Intersection numbers over a field]\label{pf:towergeo:fieldnumbers}
Let $F$ be a field. (a) There is a field intersection family over $F$, and any two agree on all projective
varieties. (b) Every field intersection family is symmetric in $a_1,\dots,a_n$. (c) Every field intersection
family is additive in each $a_i$.
\formallabel{GlobalTheory.CompanionTheory.Ar.PB.Known.fieldNumbers}
\end{lemma}

\begin{proof}
(a) is \lref{GroundTheory.Heights₂.Heights.Targets.ExistsUnique.full}, which constructs the family from
the intersection numbers of the varieties over $F$ and proves uniqueness. (b) is
\lref{ReciprocityTheory.ChartRestriction.G.AI.I1.intersection_symmetric} (field intersection families
are invariant under permutations of the multidegrees), and (c) is
\lref{ReciprocityTheory.ChartRestriction.G.AI.I1.multilin_thm} (they are additive in each multidegree).
\end{proof}

\begin{lemma}[Nonnegativity of degrees]\label{pf:towergeo:ptsint}
Assume the Kleiman criterion over fields: for every field $F$, field intersection family $I$ over $F$,
projective variety $V$ of dimension $n$ over $F$ and multidegrees $a_1,\dots,a_n$ each of which has
nonnegative intersection number with all curves on $V$, one has $I(V;a_1,\dots,a_n)\ge0$. Let $\mathrm{II}$ be
a full intersection system, $W\subset\Z[x]$ an arithmetic variety of dimension $e+1$ and $L_1,\dots,L_e$
smooth Hermitian lines such that every $L_j$ has nonnegative geometric degree on every arithmetic surface
$W'\supseteq W$. Then the integral over the complex points of $W$ of the product of the first Chern forms
of $L_1,\dots,L_e$ is nonnegative.
\formallabel{GlobalTheory.T.ptsIntegral_log_nonneg}
\end{lemma}

\begin{proof}
For a full intersection system, the integral over the complex points of an arithmetic variety of the
product of the first Chern forms of smooth Hermitian lines equals the intersection number over $\Q$ of
its generic fibre with their multidegrees (\lref{GlobalTheory.CompanionTheory.Ar.Giv.degree_of_system}
and \lref{GlobalTheory.CompanionTheory.Ar.Giv.degreeGeneric_of_degree}: the integral is the sum of the
geometric intersection numbers of the components of the fibre over a prime, which is the generic
intersection number). The generic fibre $W_\Q$ is a projective variety of dimension $e$
(\lref{GroundTheory.Heights₂.Heights.HilbertSamuel.isProjVarF_map_of_isArithVar}), and
$\mathrm{int}_\Q$ is a field intersection family
(\lref{GroundTheory.Heights₂.Heights.NumbersDVR.IntNumIsFieldIntersection.full}). Every curve
$V\subseteq W_\Q$ is the generic fibre of an arithmetic surface $W'\supseteq W$
(\lref{GroundTheory.Heights₂.Heights.HilbertSamuel.NefAmpleGeneric.isArithVar_comap_of_isProjVarF}),
and by the same degree identity in dimension one, $\mathrm{int}_\Q(V;a_j)$ is the geometric degree of $L_j$
on $W'$, which is nonnegative. By the Kleiman criterion,
$\mathrm{int}_\Q(W_\Q;a_1,\dots,a_e)\ge0$, which is the assertion.
\end{proof}

\begin{remark}\label{pf:towergeo:rem-aux}
The formal statement of \cref{pf:towergeo:ptsint} has two further hypotheses, the existence and uniqueness
of intersection numbers over $\Q$ and a statement on the support of divisors of sections; they are not
used in its proof.
\end{remark}

